% =====================================================================
% DOCUMENT 2 — RESEARCH NOTES VERSION (Master Version 1.2)
% Complete mathematical notebook. NOT for publication.
% Step 1 (kinematics) + Step 2 (constitutive) + Step 3 (variational).
% FINAL FROZEN mathematical foundation (Sections 1-3).
% All corrections from the mathematical verification audits integrated.
% Compile with pdflatex + bibtex (or Overleaf).
% =====================================================================
\documentclass[11pt]{article}

\usepackage[margin=1in]{geometry}
\usepackage{amsmath,amssymb,bm}
\usepackage{booktabs}
\usepackage{longtable}
\usepackage{cite}
\usepackage{array}

\DeclareMathOperator{\tr}{tr}
\DeclareMathOperator{\sym}{sym}
\newcommand{\gradz}{\nabla_{0}}
\newcommand{\bX}{\mathbf{X}}
\newcommand{\bx}{\mathbf{x}}
\newcommand{\bp}{\boldsymbol{\varphi}}
\newcommand{\bu}{\mathbf{u}}
\newcommand{\bF}{\mathbf{F}}
\newcommand{\bC}{\mathbf{C}}
\newcommand{\bb}{\mathbf{b}}
\newcommand{\bE}{\mathbf{E}}
\newcommand{\bep}{\boldsymbol{\varepsilon}}
\newcommand{\bH}{\mathbf{H}}
\newcommand{\bG}{\bm{\mathcal{G}}}
\newcommand{\bR}{\mathbf{R}}
\newcommand{\bS}{\mathbf{S}}
\newcommand{\bSig}{\boldsymbol{\Sigma}}
\newcommand{\bD}{\mathbb{D}}

\newcommand{\gradz}{\nabla_{0}}
\newcommand{\bX}{\mathbf{X}}
\newcommand{\bx}{\mathbf{x}}
\newcommand{\bp}{\boldsymbol{\varphi}}
\newcommand{\bphi}{\boldsymbol{\varphi}}
\newcommand{\bu}{\mathbf{u}}
\newcommand{\bF}{\mathbf{F}}
\newcommand{\bC}{\mathbf{C}}
\newcommand{\bb}{\mathbf{b}}
\newcommand{\bE}{\mathbf{E}}
\newcommand{\bep}{\boldsymbol{\varepsilon}}
\newcommand{\bH}{\mathbf{H}}
\newcommand{\bG}{\bm{\mathcal{G}}}
\newcommand{\bR}{\mathbf{R}}
\newcommand{\bU}{\mathbf{U}}
\newcommand{\bV}{\mathbf{V}}
\newcommand{\bth}{\boldsymbol{\theta}}
\newcommand{\bK}{\mathbf{K}}
\newcommand{\bS}{\mathbf{S}}
\newcommand{\bSig}{\boldsymbol{\Sigma}}
\newcommand{\bD}{\mathbb{D}}

\title{Research Notes: Finite-Strain Kinematic Framework for the
Geometrically Nonlinear $C^{1}$ BFS Mindlin Form-II Strain-Gradient
Extension \\
\large Master Version 1.2 (Final Frozen Mathematical Foundation, Steps 1--3, with Audit Derivations)}

\author{Internal working document -- not for publication}
\date{}

\begin{document}

\maketitle

\tableofcontents

% =====================================================================
\section{Document history and revision note}
% =====================================================================

This Master Version 1.2 consolidates the finite-strain kinematic
framework (Step~1), the constitutive framework (Step~2), and the
variational (Hamilton-principle) formulation (Step~3) for a
geometrically nonlinear, displacement-based, $C^{1}$ Bogner--Fox--Schmit
(BFS) formulation of Mindlin Form-II strain-gradient elasticity. All
approved corrections from the technical audit and the subsequent
revision logs have been integrated. Version~1.2 additionally corrects
four issues identified by the independent mathematical verification
audit (two critical sign errors in the strong form and its small-strain
recovery, one major error in boundary-term identifications, and one
major tensor-type error in the higher-order boundary work); the
complete proofs of these corrections are recorded in
Appendix~\ref{app:proofs}. A final self-audit confirmed exact
compatibility with the validated 2022 benchmark at every level
(kinematic, constitutive, variational, and strong-form). The
mathematical foundation (Sections~\ref{sec:kinematics}--\ref{sec:variational}) is now FROZEN: henceforth only editorial
shortening for the journal manuscript may alter the text.

% =====================================================================
\section{Purpose, scope, and the mandatory recovery rule}
% =====================================================================

\textbf{Purpose.} Establish the complete finite-strain kinematic
framework, the constitutive framework, and the variational principle
that support a geometrically nonlinear, displacement-based, $C^{1}$ BFS
finite-element formulation of Mindlin Form-II strain-gradient elasticity.

\textbf{Scope (strict).} The following are derived and discussed: motion,
reference and current configurations, material and spatial coordinates,
deformation map, displacement field, deformation gradient, right and
left Cauchy--Green tensors, Green--Lagrange strain, its small-strain
limit, its material gradient, the kinematic assumptions of Mindlin
Form-II finite-strain theory, the Helmholtz free-energy density, and the
Hamilton variational principle (action, kinetic energy, internal virtual
work, external and inertia virtual work, and the derived
strong form with its boundary conditions).

\textbf{Explicitly excluded} (later steps): the weak form, finite-element
matrices, numerical implementation, and the explicit closed forms of the
stress tensors beyond their definitions as energy derivatives.

\textbf{Mandatory recovery rule.} Every formulation must reduce exactly
to the validated 2022 benchmark under the appropriate limiting
assumption. This is verified at the kinematic level (Section~\ref{sec:recovery})
and at the constitutive and variational levels (Sections~\ref{sec:constitutive} and~\ref{sec:variational}): the
linear limit of the present framework furnishes exactly the kinematic
variables ($\bep$, $\gradz\bep$) and the energy density of the 2022
small-strain Mindlin Form-II theory. Full recovery also requires
consistent choices in the constitutive and weak-form steps; the
reduction established here is a necessary and here-satisfied
prerequisite.

% =====================================================================
\section{Reference shelf and cross-checked notation}
% =====================================================================

Every definition below was verified against the standard
finite-strain continuum-mechanics literature
\cite{BonetWood2008,Belytschko2000,Holzapfel2000,SimoHughes1998,Bathe2006}.
The five principal references and the convention each adopts are:

\begin{itemize}
\item \textbf{Bonet \& Wood} \cite{BonetWood2008}: motion $\bp$, material
$\bX$, spatial $\bx$, $\bF=\partial\bp/\partial\bX$,
$\bC=\bF^{T}\bF$, $\bb=\bF\bF^{T}$, $\bE=\tfrac12(\bC-\bI)$,
infinitesimal strain $\bep$.

\item \textbf{Belytschko, Liu \& Moran} \cite{Belytschko2000}: motion
$\bp$ (also $\bphi$), $\bF=\partial\bx/\partial\bX$, $\bC=\bF^{T}\bF$,
Green strain $\bE$; they introduce the displacement gradient
$\bH=\partial\bu/\partial\bX$.

\item \textbf{Holzapfel} \cite{Holzapfel2000}: motion $\bp$,
$\bF=\gradz\bp$, $J=\det\bF$, $\bC=\bF^{T}\bF$, $\bb=\bF\bF^{T}$,
$\bE=\tfrac12(\bC-\bI)$; strict upper-case material / lower-case
spatial indices.

\item \textbf{Simo \& Hughes} \cite{SimoHughes1998}: rigorous TL
framework, $\bF$, $\bC$, $\bb$, Green strain $\bE$.

\item \textbf{Bathe} \cite{Bathe2006}: distinctive notation
${}^{t}_{0}\mathbf{X}$ for the deformation gradient and
${}^{t}_{0}\boldsymbol{\epsilon}$ (lower-case) for the Green strain,
with material coordinates ${}^{0}X_{i}$ and spatial ${}^{t}x_{i}$.
\end{itemize}

\textbf{Selected convention and justification.} We adopt the majority
convention ($\bp,\bF,\bC,\bb,\bE,\bep$, upper-case material /
lower-case spatial indices). \emph{We deliberately reject Bathe's
notation}: his ${}^{t}_{0}\boldsymbol{\epsilon}$ for the Green strain
collides with the near-universal $\boldsymbol{\varepsilon}$ for the
infinitesimal strain, and his ${}^{t}_{0}\mathbf{X}$ for $\bF$ is
redundant with the standard $\bF$.

\begin{table}[h]
\centering
\caption{Master notation cross-check.}\label{tab:master}
\begin{tabular}{llllll}
\toprule
Concept & Bonet \& Wood & Belytschko & Holzapfel & Simo \& Hughes & Bathe \\
\midrule
Motion & $\bp$ & $\bphi/\bp$ & $\bp$ & $\bp$ & $\bx(\bX,t)$ \\
Material coord & $\bX$ & $\bX$ & $\bX$ & $\bX$ & ${}^{0}X_{i}$ \\
Spatial coord & $\bx$ & $\bx$ & $\bx$ & $\bx$ & ${}^{t}x_{i}$ \\
Def.\ gradient & $\bF$ & $\bF$ & $\bF$ & $\bF$ & ${}^{t}_{0}\mathbf{X}$ \\
Right CG & $\bC$ & $\bC$ & $\bC$ & $\bC$ & $\bC$ \\
Left CG & $\bb$ & $\bb$ & $\bb$ & $\bb$ & $\bb$/$\bc$ \\
Green strain & $\bE$ & $\bE$ & $\bE$ & $\bE$ & ${}^{t}_{0}\boldsymbol{\epsilon}$ \\
Inf.\ strain & $\bep$ & $\bep$ & $\bep$ & $\bep$ & ${}^{t}_{0}\mathbf{e}$ \\
Disp.\ grad. & $\gradz\bu$ & $\bH$ & $\gradz\bu$ & $\gradz\bu$ & -- \\
Strain grad. & $\gradz\bE$ & -- & $\gradz\bE$ & $\gradz\bE$ & -- \\
\bottomrule
\end{tabular}
\end{table}

\textbf{Index convention (fixed).} Material indices
$I,J,K,A,B,C\in\{1,2,3\}$ (upper-case); spatial indices
$i,j,k\in\{1,2,3\}$ (lower-case). Repeated indices are summed. The
Kronecker delta is $\delta_{IJ}$ (or $\delta_{ij}$); $\bI$ is the
$3\times3$ identity. The displacement components are $u_{i}$ (the
index $i$ labels the spatial component), and its material gradient is
$u_{i,J}=\partial u_{i}/\partial X_{J}$ (component index $i$,
derivative index $J$). The strain-gradient tensor is
$\bm{\mathcal{G}}=\gradz\bE$ with components
$\mathcal{G}_{IJK}=\partial E_{IJ}/\partial X_{K}=\mathcal{G}_{JIK}$.

\textbf{Gradient-operator convention.} We write $\gradz(\cdot)=
\partial(\cdot)/\partial\bX$ for the \emph{material} gradient and
$\nabla(\cdot)=\partial(\cdot)/\partial\bx$ for the \emph{spatial}
gradient. Several authors use a capital ``$\mathrm{Grad}$'' to emphasise
the material gradient \cite{BonetWood2008}; we follow that convention
in prose and write $\mathrm{Grad}(\mathbf{E})\equiv\gradz\mathbf{E}$.

% =====================================================================
\section{Tensor-identity ledger (material / spatial / two-point)}
\label{sec:ledger}
% =====================================================================

To guarantee notational consistency, every tensor is tagged by its frame
and transformation class:

\begin{itemize}
\item \textbf{Material (reference) tensors} -- transform under a change of
material basis only; invariant under a spatial observer change
(objective). Examples: $\bC$, $\bE$, $\bU$,
$\bm{\mathcal{G}}=\gradz\bE$, and the material gradient of the rotation
$\gradz\bR$.

\item \textbf{Spatial (current) tensors} -- defined in the current
configuration; transform under a superposed rigid-body motion as proper
spatial tensors. Examples: $\bb$, $\bV$, and (later) the Cauchy stress.

\item \textbf{Two-point tensors} -- carry one material and one spatial
index; map between reference and current frames. Examples:
$\bF$ ($F_{iJ}$) and the polar rotation $\bR$ ($R_{iJ}$). They are
\emph{not} objective as bare two-point objects, but their combinations
$\bC=\bF^{T}\bF$ and $\bE$ are material and objective.
\end{itemize}

\noindent\textbf{Reference identities used throughout.}
\begin{align}
\bC &= \bF^{T}\bF = \bU^{2}, \qquad \bb = \bF\bF^{T} = \bV^{2},
\label{eq:id_Cb}\\
\bb &= \bF\bC\bF^{-1}, \qquad \bC = \bF^{-1}\bb\bF,
\label{eq:id_pushpull}\\
\bF^{-1} &= \frac{\partial\bX}{\partial\bx}
\quad\text{(spatial deformation gradient / inverse)},
\label{eq:id_invF}\\
\gradz f &= \bF^{T}\bigl(\nabla f\bigr)
\quad\text{(exact material/spatial gradient chain rule for any scalar $f$)},
\label{eq:id_gradpull}
\end{align}

\noindent (For a pulled-forward tensor field the spatial gradient involves
the additional term $\partial\bF/\partial\bX$; the exact scalar
chain-rule \eqref{eq:id_gradpull} suffices to show that a spatial
gradient is recovered from the material one only through
$\bF^{-T}$.)

% =====================================================================
\section{Detailed derivations (items 1--14)}
\label{sec:kinematics}
% =====================================================================

\subsection{Configurations and coordinates (items 1--5)}
The body in its undeformed state occupies the \emph{reference
configuration} $\mathcal{B}_{0}\subset\mathbb{R}^{3}$. Each material
point is labelled by its \emph{material (reference) coordinates}
$\bX\in\mathcal{B}_{0}$, fixed for all time. At time $t$ the body
occupies the \emph{current (spatial) configuration}
$\mathcal{B}_{t}\subset\mathbb{R}^{3}$, whose points are described by
\emph{spatial (current) coordinates} $\bx\in\mathcal{B}_{t}$
\cite{Holzapfel2000,BonetWood2008}.

\begin{align}
\mathcal{B}_{0} &\subset \mathbb{R}^{3},\qquad \bX\in\mathcal{B}_{0},
\label{eq:n_ref}\\
\mathcal{B}_{t} &\subset \mathbb{R}^{3},\qquad \bx\in\mathcal{B}_{t}.
\label{eq:n_cur}
\end{align}

\subsection{Motion and deformation map (item 6)}
The \emph{motion} (deformation map) is
\begin{align}
\bx &= \bp(\bX,t),
\label{eq:n_motion}
\end{align}
with inverse (diffeomorphism, assumption A1)
\begin{align}
\bX &= \bp^{-1}(\bx,t).
\label{eq:n_inv}
\end{align}

We use $\bp$ \cite{BonetWood2008,Holzapfel2000,SimoHughes1998} rather
than Belytschko's $\bphi$ \cite{Belytschko2000} or Bathe's
$\bx(\bX,t)$ \cite{Bathe2006}; all denote the same map.

\subsection{Displacement field (item 7)}
\begin{align}
\bu(\bX,t) &= \bp(\bX,t) - \bX.
\label{eq:n_disp}
\end{align}

Expressed as a function of $\bX$ (Lagrangian), because the formulation
is Total-Lagrangian (A4).

\subsection{Deformation gradient and Jacobian (item 8)}
\begin{align}
\bF(\bX,t) &= \gradz\bp = \frac{\partial\bx}{\partial\bX},
\qquad F_{iJ}=\frac{\partial x_{i}}{\partial X_{J}},
\label{eq:n_F}\\
J &= \det\bF > 0.
\label{eq:n_J}
\end{align}

Using $\bx=\bX+\bu$,
\begin{align}
F_{iJ} &= \frac{\partial}{\partial X_{J}}(X_{i}+u_{i})
        = \delta_{iJ} + u_{i,J},
\label{eq:n_Fdec}\\
\bF &= \bI + \gradz\bu.
\label{eq:n_Fdec_t}
\end{align}

$\bF$ is a two-point tensor \cite{Holzapfel2000}. $J=\det\bF$ gives the
local volume ratio $dv=J\,dV$; the requirement $J>0$ (A1) means the
deformation is orientation-preserving and invertible, so $\bF^{-1}$ and
all pull-back/push-forward operations exist.

\subsection{Right Cauchy--Green deformation tensor (item 9)}
\begin{align}
\bC &= \bF^{T}\bF,\qquad C_{IJ}=F_{kI}F_{kJ}.
\label{eq:n_C}
\end{align}

$\bC$ measures the squared length of a material line element:
$|d\bx|^{2}=d\bX^{T}\bC d\bX$, hence it is a material tensor
\cite{Holzapfel2000}. Symmetric, and for any $\mathbf{v}\neq\mathbf{0}$,
$\mathbf{v}\cdot\bC\mathbf{v}=\|\bF\mathbf{v}\|^{2}>0$ (invertible
$\bF$), so $\bC$ is symmetric positive definite (SPD) \checkmark.

\subsection{Left Cauchy--Green deformation tensor (item 10)}
\begin{align}
\bb &= \bF\bF^{T},\qquad b_{ij}=F_{iK}F_{jK}.
\label{eq:n_b}
\end{align}

$\bb=\bF\bC\bF^{-1}$ (push-forward of $\bC$), a spatial tensor
\cite{Holzapfel2000}. Symmetric SPD \checkmark. \emph{Not} the Finger
tensor ($\bb^{-1}$); Ogden's $\mathbf{B}$ is avoided because
$\mathbf{B}$ is the FE strain-displacement matrix.

\subsection{Green--Lagrange strain tensor (item 11)}
\begin{align}
\bE &= \tfrac12(\bC-\bI)=\tfrac12(\bF^{T}\bF-\bI),
\label{eq:n_E}\\
E_{IJ} &= \tfrac12\bigl(F_{kI}F_{kJ}-\delta_{IJ}\bigr).
\label{eq:n_Ecomp}
\end{align}

Displacement form (full expansion):
$F_{kI}=\delta_{kI}+u_{k,I}$, $F_{kJ}=\delta_{kJ}+u_{k,J}$, so
$F_{kI}F_{kJ}=\delta_{IJ}+u_{I,J}+u_{J,I}+u_{k,I}u_{k,J}$, giving
\begin{align}
E_{IJ} &= \tfrac12\bigl(u_{I,J}+u_{J,I}+u_{K,I}u_{K,J}\bigr).
\label{eq:n_Eu}
\end{align}
Symmetric ($E_{IJ}=E_{JI}$) \checkmark. The first two terms are the
infinitesimal strain; the last term $\tfrac12u_{K,I}u_{K,J}$ is the
geometric (initial-displacement) nonlinearity that later yields the
initial-stress stiffness.

\textbf{Selection among strain measures} (citable precedent
\cite{Holzapfel2000,Hill1968}): \emph{Green--Lagrange}
$\bE=\tfrac12(\bC-\bI)$ -- material, TL, exact quadratic recovery of
$\bep$ in the linear limit \textbf{(adopted)}. \emph{Almansi}
$\tfrac12(\bI-\bb^{-1})$ -- spatial, ill-suited to a TL $C^{1}$ BFS
element \textbf{(rejected)}. \emph{Biot} $\tfrac12(\bU-\bI)$ -- awkward
gradient \textbf{(rejected)}. \emph{Hencky/log} $\tfrac12\ln\bC$ --
best for very large stretches but its pull-back does not collapse
cleanly to the BFS spatial-gradient operator, threatening exact benchmark
recovery \textbf{(rejected)}. \emph{Infinitesimal} $\bep$ --
the 2022 benchmark measure; it is the linear limit of $\bE$, not a
finite-strain competitor.

\subsection{Small-strain limit verification (item 12)}

With $\bF=\bI+\bH$, $\bH=\gradz\bu$:
$\bC=\bI+\bH^{T}+\bH+\bH^{T}\bH$, hence
\begin{align}
\bE &= \tfrac12(\bH+\bH^{T})+\tfrac12\bH^{T}\bH
    = \bep + \tfrac12\bH^{T}\bH.
\label{eq:n_limit1}
\end{align}

As $\|\bH\|\to0$, $\|\bH^{T}\bH\|=O(\|\bH\|^{2})\to0$, so
\begin{align}
\bE \to \bep,\qquad
E_{IJ}\to\tfrac12(u_{I,J}+u_{J,I})=\varepsilon_{IJ}.
\label{eq:n_limit2}
\end{align}

The limit is exact (discarded term strictly $O(\|\bH\|^{2})$).

\subsection{Material gradient of the Green--Lagrange strain (item 13)}
\label{sec:Gdef}

\begin{align}
\bm{\mathcal{G}} &= \gradz\bE,\qquad \mathcal{G}_{IJK}=\frac{\partial E_{IJ}}{\partial X_{K}}.
\label{eq:n_Gdef}
\end{align}

In terms of $\bF$ (product rule):
\begin{align}
\mathcal{G}_{IJK} &= \tfrac12\bigl(F_{kI,K}F_{kJ}+F_{kI}F_{kJ,K}\bigr).
\label{eq:n_GcompF}
\end{align}

In terms of displacement (product rule, mixed partials commute):
\begin{align}
\mathcal{G}_{IJK} &= \tfrac12\bigl(
      u_{J,IK}+u_{I,JK}+u_{k,IK}u_{k,J}+u_{k,I}u_{k,JK}
   \bigr).
\label{eq:n_Gcompu}
\end{align}

Symmetry $\mathcal{G}_{IJK}=\mathcal{G}_{JIK}$ (strain indices only) \checkmark. Equation
\eqref{eq:n_Gcompu} contains second displacement derivatives
$u_{i,IK}$; this is the origin of the $C^{1}$ requirement (Section~\ref{sec:bfs}).

\subsection{Mindlin Form-II finite-strain kinematic assumptions (item 14)}
\label{sec:form2}

Mindlin's micro-structured elasticity depends on strain, strain
gradient, rotation, and rotation-gradient \cite{Mindlin1964,
Mindlin1965,MindlinEshel1968}. The finite-strain generalisation must
preserve this structure and reduce to it linearly.

\textbf{Polar decomposition} \cite{Holzapfel2000}: every invertible
$\bF$ admits $\bF=\bR\bU=\bV\bR$, with
$\bR\in SO(3)$ proper orthogonal and $\bU=\sqrt{\bC}$,
$\bV=\sqrt{\bb}$ SPD. The rotation part is $\bR$ (two-point tensor).

\textbf{Finite rotation variable.} $\bR$ is the Mindlin Form-II
rotation measure (exact finite-rotation generalisation of the
small-strain rotation). Its axial vector $\bth$ is used when a vector
form is convenient.

\textbf{Rotation-gradient variable (kinematic).}
\begin{align}
\gradz\bR \quad\text{or}\quad \bK=\gradz\bth,\qquad
K_{iK}=\frac{\partial\theta_{i}}{\partial X_{K}}.
\label{eq:n_rotgrad}
\end{align}

In the linear limit $\bR\to\bI$ and $\bK\to$ the classical Form-II
rotation gradient. \emph{The rotation gradient is a kinematic variable
only; it is not an argument of the free-energy density (see Section~\ref{sec:constitutive}).}

\textbf{Complete Mindlin Form-II kinematic set (finite strain):}
strain $\bE$, strain gradient $\bm{\mathcal{G}}=\gradz\bE$, rotation
$\bR$ (or $\bth$), rotation gradient $\gradz\bR$ (or
$\bK=\gradz\bth$).

\textbf{Assumption ledger.}
\begin{itemize}
\item[A1] $\bp$ smooth orientation-preserving diffeomorphism, $J>0$
\cite{Holzapfel2000,BonetWood2008}.

\item[A2] Displacement-based primal field $\bu$; no independent
microrotation DOF -- preserves the validated $C^{1}$ BFS structure
(programme constraint).

\item[A3] $\bp,\bu\in C^{1}(\mathcal{B}_{0})$. Second material
derivatives exist piecewise in $L^{2}$ (i.e.\ $H^{2}$ regularity),
exactly what the BFS element provides (Section~\ref{sec:bfs}).

\item[A4] Strain measure $\bE$ (Total Lagrangian)
\cite{Holzapfel2000,SimoHughes1998,BonetWood2008}.

\item[A5] Rotation measure $\bR$ (polar) \cite{Holzapfel2000}.

\item[A6] Strain-gradient variable $\gradz\bE$ (material)
\cite{Beheshti2016,Mindlin1964}.
\end{itemize}

% =====================================================================
\section{Literature verification of the finite-strain strain-gradient
measure $\bm{\mathcal{G}}=\mathrm{Grad}(\mathbf{E})$}
\label{sec:Gverify}
% =====================================================================

\textbf{Question.} Is $\bm{\mathcal{G}}=\gradz\mathbf{E}$ the correct
finite-strain analogue of the classical Mindlin Form-II strain-gradient
measure, or should one use $\mathrm{Grad}(\mathbf{F})$,
$\mathrm{Grad}(\mathbf{C})$, or another measure?

\subsection{Has $\mathrm{Grad}(\mathbf{E})$ been used at finite strain?}

Yes. Beheshti \cite{Beheshti2016} explicitly generalises linear
strain-gradient theory to finite elastic deformation by adopting the
Green--Lagrange strain tensor, with a generalised
Saint-Venant--Kirchhoff strain-gradient material in both material and
spatial configurations. More generally, geometrically nonlinear
higher-gradient elasticity with the strain energy depending on a
finite-strain measure and its material gradient is treated by Javili,
dell'Isola \& Steinmann
\cite{Javili2013} and in the strain-gradient-thermoelasticity literature
built on Mindlin Form-II \cite{Aouadi2020}. The displacement-based
finite-strain strain-gradient formulation of Rubin \cite{Rubin2008}
likewise rests on a strain measure and its gradient. Hence
$\mathrm{Grad}(\mathbf{E})$ is a literature-supported, not ad-hoc,
choice.

\subsection{Alternative finite-strain strain-gradient measures}

\begin{itemize}
\item \textbf{$\mathrm{Grad}(\mathbf{F})$} -- the third-order tensor
$F_{iJ,K}$ (gradient of the deformation gradient). This is the most
primitive variable; it underlies the classical higher-grade theories of
Toupin \cite{Toupin1962}, Green \& Rivlin \cite{GreenRivlin1964}, and
the finite-deformation second-gradient theories. It is the
natural primitive for a fully general higher-grade hyperelastic solid
with energy $W(\mathbf{F},\mathrm{Grad}\mathbf{F})$.

\item \textbf{$\mathrm{Grad}(\mathbf{C})$} -- since
$\mathbf{E}=\tfrac12(\mathbf{C}-\bI)$, $\mathrm{Grad}(\mathbf{E})
=\tfrac12\,\mathrm{Grad}(\mathbf{C})$; the two are equivalent up to a
constant factor.

\item \textbf{$\mathrm{Grad}(\mathbf{E})$} -- adopted here.

\item \textbf{Decomposed form} -- Fleck \& Hutchinson
\cite{FleckHutchinson1997} and the Form-II tradition decompose the
strain gradient into a stretch-gradient and a rotation-gradient part;
this is exactly the structure $\gradz\bE$ (strain gradient) $+$
$\gradz\bR$ (rotation gradient) used in Section~\ref{sec:form2}.
\end{itemize}

\subsection{Why $\mathrm{Grad}(\mathbf{E})$ is preferable here}

For a \emph{displacement-based}, \emph{Total-Lagrangian},
\emph{$C^{1}$ BFS} formulation that must recover the 2022 small-strain
Mindlin Form-II benchmark, $\gradz\bE$ is the unique best choice:

\begin{enumerate}
\item \textbf{Exact benchmark recovery (mandatory rule).} $\gradz\bE\to
\gradz\bep$ in the linear limit (Section~\ref{sec:recovery}), so the
finite-strain variable reduces exactly to the infinitesimal
strain-gradient of the 2022 theory.

\item \textbf{Total-Lagrangian consistency.} $\gradz\bE$ is a material
tensor. By the exact chain rule \eqref{eq:id_gradpull}, a spatial
gradient would be related to the material one only through the inverse
deformation gradient $\bF^{-T}$; defining the constitutive gradient in
material coordinates keeps every later operator reference-based and the
2022 reduction direct.

\item \textbf{Form-II structure.} Mindlin Form-II separates strain
gradient and rotation gradient \cite{MindlinEshel1968,Aouadi2020};
$\gradz\bE$ pairs naturally with $\gradz\bR$, whereas
$\mathrm{Grad}(\mathbf{F})$ entangles stretch and rotation in a single
third-order object.

\item \textbf{$C^{1}$ BFS compatibility.} Both $\gradz\bE$ and
$\mathrm{Grad}(\mathbf{F})$ require the same second displacement
derivatives $u_{i,IK}$ (compare \eqref{eq:n_Gcompu} with
$F_{iJ,K}=u_{i,JK}$), so the element requirement is identical;
$\gradz\bE$ is preferred because it is the literal generalisation of
the 2022 small-strain gradient variable.
\end{enumerate}

\emph{Conclusion:} $\bm{\mathcal{G}}=\gradz\mathbf{E}$ is correct and
preferable for this formulation; $\mathrm{Grad}(\mathbf{F})$ is the
more general primitive but entangles stretch and rotation and is
therefore not adopted.

% =====================================================================
\section{Mindlin Form-II finite-strain compatibility}
\label{sec:compat}
% =====================================================================

\textbf{Is extending classical Mindlin Form-II strain-gradient
elasticity to finite strains mathematically justified?} Yes, provided
the kinematic structure (strain + strain gradient + rotation + rotation
gradient) is preserved, which the present framework does.

\begin{itemize}
\item \textbf{Lineage.} Mindlin \cite{Mindlin1964,Mindlin1965} founded
strain-gradient (micro-structured) elasticity; Mindlin \& Eshel
\cite{MindlinEshel1968} formalised the first strain-gradient theories,
and Mindlin \& Tiersten \cite{MindlinTiersten1962} the couple-stress
version. Toupin \cite{Toupin1962} and Green \& Rivlin
\cite{GreenRivlin1964} provided the higher-grade / multipolar
foundations; Germain \cite{Germain1973} derived the equilibrium
relations via virtual power. These are the recognised antecedents of
any finite-strain gradient theory.

\item \textbf{Form II identification.} In Mindlin Form-II the strain
energy density is a function of the strain and its first-order gradient
\cite{Mindlin1964,MindlinEshel1968,Aouadi2020}; the rotation-gradient content is
recovered when the strain gradient is decomposed into stretch- and
rotation-gradient parts \cite{FleckHutchinson1997}. The present
kinematics $(\bE,\gradz\bE,\bR,\gradz\bR)$ realise exactly this
structure at finite strain. The \emph{constitutive} state, however, is
$\{\bE,\bm{\mathcal{G}}\}$ (strain-gradient Form-II); the rotation
gradient $\bK$ is kinematic only and does not enter the energy.

\item \textbf{Displacement-based (compatible) reduction.} The present
programme is \emph{displacement-based}: there is no independent
microrotation degree of freedom. The microrotation is therefore
identified with the macroscopic (polar) rotation $\bR$ -- the
``compatible'' or ``reduced'' Mindlin Form-II, as opposed to the
full Cosserat / micropolar version with an independent rotation field
\cite{Koiter1964,Toupin1962}. This is consistent with the validated
2022 small-strain formulation (which likewise derived rotation from the
displacement gradient) and is the standard displacement-based
realisation of Form-II \cite{AskesAifantis2011,LazarMaugin2005}.

\item \textbf{Relation to modern finite-strain gradient theories.}
Beheshti \cite{Beheshti2016} and Javili et al.\ \cite{Javili2013}
treat geometrically nonlinear higher-gradient elasticity with energy
depending on a finite-strain measure and its gradient; Rubin
\cite{Rubin2008} and the Gurtin--Anand--Lele gradient framework
\cite{GurtinAnandLele2007} provide finite-strain gradient
formulations in related contexts. The present kinematics are fully
consistent with this body of work: they supply the finite-strain
measure ($\bE$) and its material gradient ($\gradz\bE$) that those
theories require, within a displacement-based TL description.

\item \textbf{Consistency verdict.} The proposed kinematics are
\emph{fully consistent} with the existing literature; no new
kinematic principle is introduced. The only programme-specific choice
is the displacement-based reduction of the microrotation to the polar
rotation, which is the standard compatible Form-II realisation.
\end{itemize}

% =====================================================================
\section{BFS $C^{1}$ continuity requirement}
\label{sec:bfs}
% =====================================================================

The document's statement that $C^{1}$ continuity is required is
necessary but was previously under-justified. The rigorous
justification is as follows.

\subsection{Source of the high-order derivative}

From \eqref{eq:n_Gcompu}, the strain-gradient variable contains second
material derivatives of displacement, $u_{i,IK}$. Consequently any
energy density that depends on $\gradz\bE$ depends on
$\nabla_{0}\nabla_{0}\bu$. For a \emph{displacement-based} (primal)
Galerkin formulation, the double stress is work-conjugate to
$\gradz\bE$, and the virtual-work term contains
$\int_{\Omega_{0}}\boldsymbol{\tau}:\delta(\gradz\bE)\,\mathrm{d}V$.
Integrating by parts once on the reference domain yields a domain term
involving products of \emph{second} derivatives of both test and trial
functions, e.g.\ terms of the form $\int (\cdot)_{,IJ}\,(\delta
u)_{,KL}\,\mathrm{d}V$.

\subsection{Regularity required}

Such a term is finite only if both trial and test functions have
square-integrable second derivatives, i.e.\ $\bu,\delta\bu\in
H^{2}(\Omega_{0})$. A $C^{0}$ element guarantees only
$H^{1}$ regularity (first derivatives in $L^{2}$); its second
derivatives are distributions, not in $L^{2}$, so the gradient-energy
contribution is ill-defined for a $C^{0}$ primal field. A $C^{0}$
approach would therefore \emph{force} a mixed formulation in which the
strain gradient / double stress is introduced as an independent
variable -- incompatible with the displacement-based BFS architecture.

\subsection{What $C^{1}$ provides}

The Bogner--Fox--Schmit element is a bicubic-Hermite rectangular
element that is $C^{1}$ across element boundaries, so its shape
functions and their first derivatives are continuous and the second
derivatives are \emph{piecewise} continuous and in $L^{2}$
\cite{Bathe2006,BonetWood2008}. Thus the BFS field delivers exactly the
$H^{2}$ regularity required: second derivatives exist and are
square-integrable. Equivalently, the BFS element's $C^{1}$ continuity is precisely the minimal regularity that renders the displacement-based strain-gradient weak form well posed; no stronger ($C^{2}$) continuity is needed. \emph{Full $C^{2}$ is not required} -- $C^{1}$ (i.e.
$H^{2}$) suffices, because the second derivatives need only belong to
$L^{2}$, not be continuous across elements.

\subsection{Summary of continuity requirements}

\begin{itemize}
\item Continuity of displacement $\bu$: required (as in any FE
displacement method).

\item Continuity of first derivatives $\nabla_{0}\bu$: required -- this
is the defining property of a $C^{1}$ element and is what a $C^{0}$
element lacks.

\item Element-wise second derivatives $\nabla_{0}\nabla_{0}\bu$: must
exist and be in $L^{2}$; provided piecewise by the BFS bicubic-Hermite
basis. They need not be continuous across elements.

\item Hence a $C^{0}$ element is insufficient for a
displacement-based primal formulation of this gradient theory; the
$C^{1}$ BFS element is the minimal-sufficient choice.
\end{itemize}

% =====================================================================
\section{Two-dimensional specialisation}
\label{sec:2d}
% =====================================================================

Because the computational implementation will be two-dimensional, the
3D kinematics are reduced to a plane. Two standard reductions are
plane strain and plane stress; both are stated, together with the
assumptions used in the present FEM framework.

\subsection{Plane strain}

Assumptions: the out-of-plane displacement vanishes and there is no
dependence on the out-of-plane coordinate,
\begin{align}
u_{3}=0,\qquad \frac{\partial(\cdot)}{\partial X_{3}}=0.
\label{eq:ps_u3}
\end{align}

The deformation gradient reduces to the block form
\begin{align}
\bF =
\begin{bmatrix}
F_{11} & F_{12} & 0\\
F_{21} & F_{22} & 0\\
0      & 0      & 1
\end{bmatrix},
\qquad F_{\alpha\beta}=\delta_{\alpha\beta}+u_{\alpha,\beta},
\label{eq:ps_F}
\end{align}
with Greek indices $\alpha,\beta\in\{1,2\}$. Consequently
$C_{33}=F_{3K}F_{3K}=1$ and $E_{33}=0$; the in-plane strain is
$E_{\alpha\beta}=\tfrac12(F_{k\alpha}F_{k\beta}-\delta_{\alpha\beta})$.
The strain gradient reduces to the in-plane components
$\mathcal{G}_{\alpha\beta\gamma}$ ($\alpha,\beta,\gamma\in\{1,2\}$); all
components with an index equal to $3$ vanish by
\eqref{eq:ps_u3}. Plane
strain is the natural 2D model for the BFS element when the
through-thickness variation is ignored (the BFS element is a
$C^{1}$ rectangular plate/mid-surface element).

\subsection{Plane stress}

Plane stress is a \emph{constitutive} constraint, not a kinematic one:
the out-of-plane (Cauchy) stress vanishes, $\sigma_{33}=0$. In a
Total-Lagrangian description with the Green--Lagrange strain, this
constraint determines the out-of-plane normal stretch $F_{33}$ (or
$E_{33}$) as a function of the in-plane deformation through the (future)
constitutive law. Kinematically, the in-plane displacement field
$u_{\alpha}(X_{1},X_{2})$ and its gradient -- and therefore the
in-plane strain gradient $\mathcal{G}_{\alpha\beta\gamma}$ -- are identical to
the plane-strain case; only the out-of-plane normal stretch is slaved
to the in-plane state via the constitutive relation. We therefore do
\emph{not} set $E_{33}=0$ kinematically under plane stress; we state
the $\sigma_{33}=0$ condition and defer its satisfaction to the
constitutive step. This distinction is important and was missing in
Revision~1.

\subsection{Reduction summary and framework assumptions}

\begin{itemize}
\item[A7] The 2D model uses in-plane indices $\alpha,\beta,\gamma\in
\{1,2\}$; the out-of-plane direction is $X_{3}$.

\item[A8] Plane strain: $u_{3}=0$, $\partial/\partial X_{3}=0$
(kinematic). Adopted as the default for the BFS plate/mid-surface
implementation.

\item[A9] Plane stress: $\sigma_{33}=0$ (constitutive); out-of-plane
normal stretch slaved to in-plane state via the constitutive law, not
set to zero kinematically.

\item The strain-gradient variable in 2D is the in-plane third-order
tensor $\mathcal{G}_{\alpha\beta\gamma}$; the BFS $C^{1}$ requirement
(Section~\ref{sec:bfs}) applies identically in 2D.
\end{itemize}

\noindent\textbf{Reduction to the two-dimensional BFS finite-element
formulation.} The three-dimensional kinematic fields
$\bp,\bF,\bE,\bm{\mathcal{G}}=\gradz\bE,\bR,\gradz\bR$ are restricted to the
two-dimensional mid-surface $\Omega_{0}^{2D}\subset\mathcal{B}_{0}$ with
every out-of-plane derivative set to zero, so the in-plane components
$F_{\alpha\beta}$, $E_{\alpha\beta}$, the in-plane strain gradient
$\mathcal{G}_{\alpha\beta\gamma}$, and the in-plane rotation-gradient reduce
exactly to the variables the BFS element will interpolate. The
bicubic-Hermite BFS basis then supplies the required $C^{1}$ (i.e.\
$H^{2}$) continuity for the in-plane displacement and its derivatives
across element boundaries, so the 3D strain-gradient operator collapses
without approximation to the 2D operator used in the implementation;
under plane stress the out-of-plane normal stretch is recovered from the
in-plane state through the (future) constitutive law.

% =====================================================================
\section{Consolidated verification suite}
\label{sec:recovery}
% =====================================================================

\subsection{Small-strain limit / recovery of the 2022 benchmark}

Established in equations~\eqref{eq:n_limit1}--\eqref{eq:n_limit2} and
Section~\ref{sec:Gdef}: as $\|\gradz\bu\|\to0$,
\begin{align}
\bE &\to \bep, &
\bm{\mathcal{G}}=\gradz\bE &\to \gradz\bep, &
\bR &\to \bI, &
\bK=\gradz\bth &\to \text{(infinitesimal rotation gradient)}.
\end{align}

The finite-strain Mindlin Form-II kinematic set collapses exactly to the
2022 small-strain set; the recovery is exact (discarded terms strictly
$O(\|\gradz\bu\|^{2})$). \checkmark \textbf{Mandatory recovery rule
satisfied at the kinematic level.}

\subsection{Objectivity verification}
\label{sec:obj}

Under a superposed rigid-body motion
$\bar{\bx}=\mathbf{Q}(t)\bx+\mathbf{c}(t)$,
$\mathbf{Q}\in SO(3)$, $\bar{\bF}=\mathbf{Q}\bF$
\cite{Holzapfel2000,MarsdenHughes1983}:

\begin{itemize}
\item $\bar{\bC}=\bar{\bF}^{T}\bar{\bF}
=\bF^{T}\mathbf{Q}^{T}\mathbf{Q}\bF=\bC$ -- invariant, objective
(material) \checkmark.

\item $\bar{\bE}=\tfrac12(\bar{\bC}-\bI)=\bE$ -- objective \checkmark.

\item $\bar{\bb}=\mathbf{Q}\bb\mathbf{Q}^{T}$ -- proper spatial
tensor, objective \checkmark.

\item $\bar{\bm{\mathcal{G}}}=\gradz\bar{\bE}=\gradz\bE=\bm{\mathcal{G}}$ -- material
third-order tensor, objective \checkmark.

\item $\bar{\bR}=\mathbf{Q}\bR$; $\gradz\bar{\bR}=\mathbf{Q}\gradz\bR$ --
the spatial index transforms by $\mathbf{Q}$, so $\gradz\bR$ is
objective (proper spatial-tensor-valued material gradient), but it is
\emph{not} invariant ($\gradz\bar{\bR}\neq\gradz\bR$ for
$\mathbf{Q}\neq\mathbf{I}$). \checkmark
\end{itemize}

All kinematic variables are objective \checkmark
\cite{MarsdenHughes1983,Holzapfel2000}.

\subsection{Tensor symmetry checks}

\begin{itemize}
\item $\bC^{T}=\bC$, $\bb^{T}=\bb$, $\bE^{T}=\bE$ -- symmetric
\checkmark.

\item $\bC,\bb$ SPD (proved Sections~9--10) \checkmark.

\item $\bm{\mathcal{G}}_{IJK}=\bm{\mathcal{G}}_{JIK}$ (symmetry in strain indices only)
\checkmark.

\item $\bR^{T}\bR=\bI$, $\det\bR=1$ (proper orthogonal) \checkmark.
\end{itemize}

% =====================================================================
\section{Common mistakes found in the literature}
% =====================================================================

\begin{enumerate}
\item $\bE=\tfrac12(\bF+\bF^{T})-\bI$ -- \textbf{wrong}; the strain is
built from $\bF^{T}\bF$, not the unsquared $\bF+\bF^{T}$. Correct:
$\bE=\tfrac12(\bF^{T}\bF-\bI)$.

\item Confusing Green--Lagrange $\tfrac12(\bC-\bI)$ with Almansi
$\tfrac12(\bI-\bb^{-1})$ -- different frames; do not interchange.

\item Calling $\bb=\bF\bF^{T}$ the ``Finger tensor'' -- the Finger
tensor is $\bb^{-1}$; avoid Ogden's $\mathbf{B}$ for $\bb$ (clashes
with the FE strain-displacement matrix).

\item Stating $J=\det\bF$ without the $J>0$ requirement -- invertibility
is mandatory before any pull-back/push-forward.

\item In a TL strain-gradient theory, defining the strain gradient as the
\emph{spatial} $\nabla\bE$ instead of $\gradz\bE$ -- introduces an
implicit $\bF^{-T}$ pull-back and can break exact recovery of the 2022
benchmark. Use $\gradz\bE$ (Section~\ref{sec:Gverify}).

\item In Mindlin Form-II at finite strain, reusing the infinitesimal
rotation $\tfrac12(\gradz\bu-\gradz\bu^{T})$ for large rotations --
invalid; the correct finite-rotation measure is the polar rotation
$\bR$ (Section~\ref{sec:form2}).

\item Notation collision from Bathe: ${}^{t}_{0}\boldsymbol{\epsilon}$
for the Green strain collides with $\boldsymbol{\varepsilon}$
(infinitesimal strain). We use $\bE$ and $\bep$ distinctly.

\item Index slip: writing the displacement gradient as $u_{I,J}$ and
mixing the component/strain index with the derivative index. We keep
$u_{i,I}$ (component $i$, derivative $I$) unambiguously.

\item \textbf{Plane-stress vs plane-strain confusion}: under plane
stress one must not set $E_{33}=0$ kinematically; $\sigma_{33}=0$ is a
constitutive constraint that determines $F_{33}$ (Section~\ref{sec:2d}).
\end{enumerate}

% =====================================================================
\section{Constitutive Framework}
\label{sec:constitutive}
% =====================================================================

\subsection{Purpose and restriction of this section}

This section establishes only the constitutive assumptions required for a geometrically nonlinear, displacement-based, $C^1$ Bogner--Fox--Schmit finite element formulation based on finite-strain Mindlin Form-II strain-gradient elasticity. The kinematic variables are those frozen in Section~\ref{sec:kinematics}:

\begin{equation}
\mathbf{F}=\Grad\bm{\varphi},
\qquad
\mathbf{C}=\mathbf{F}^{T}\mathbf{F},
\qquad
\mathbf{E}=\frac{1}{2}(\mathbf{C}-\mathbf{I}),
\qquad
\bm{\mathcal{G}}=\Grad\mathbf{E},
\qquad
\mathcal{G}_{ABC}=E_{AB,C}.
\label{eq:sec2_kinematic_input}
\end{equation}

No governing equations, stress tensors, tangent operators, weak forms, Hamilton principle, or finite element matrices are derived in this section.

The objective is to define a Helmholtz free-energy density that is compatible with:

\begin{enumerate}
\item total-Lagrangian finite-strain kinematics;
\item Mindlin Form-II strain-gradient elasticity;
\item displacement-only kinematics;
\item $C^1$ BFS interpolation;
\item exact recovery of the validated small-strain 2022 benchmark formulation.
\end{enumerate}

\subsection{Literature background: constitutive frameworks for finite-strain strain-gradient elasticity}

Classical finite elasticity describes the material response through a Helmholtz free-energy density depending on finite-deformation measures such as $\mathbf{F}$, $\mathbf{C}$, or $\mathbf{E}$, subject to frame indifference and material symmetry restrictions \cite{Rivlin1948,Ogden1997,Holzapfel2000,BonetWood2008,GurtinFriedAnand2010}. In a total-Lagrangian formulation, it is standard to use material strain measures such as $\mathbf{C}$ or $\mathbf{E}$.

Generalized continuum theories introduce additional kinematic measures to account for microstructural length-scale effects. The foundational higher-gradient and multipolar theories were developed by Toupin, Green and Rivlin, Germain, Mindlin, and Mindlin--Eshel \cite{Toupin1962,Toupin1964,GreenRivlin1964Simple,GreenRivlin1964Multipolar,Germain1973,Mindlin1964,MindlinEshel1968}. In Mindlin Form-II first strain-gradient elasticity, the small-strain free energy depends on the infinitesimal strain tensor and its first gradient:

\begin{equation}
\Psi=\Psi(\bm{\varepsilon},\Grad\bm{\varepsilon}).
\label{eq:sec2_small_strain_formII_literature}
\end{equation}

The finite-strain analogue is not unique. The literature contains several related possibilities:

\begin{enumerate}
\item classical hyperelasticity: $\Psi=\Psi(\mathbf{C})$ or $\Psi=\Psi(\mathbf{E})$;
\item Toupin-type second-gradient elasticity: $\Psi=\Psi(\mathbf{F},\Grad\mathbf{F})$;
\item Green--Lagrange Form-II strain-gradient elasticity: $\Psi=\Psi(\mathbf{E},\Grad\mathbf{E})$;
\item more general finite-gradient elasticity frameworks involving first and second deformation gradients \cite{Bertram2015,JaviliDellIsolaSteinmann2013};
\item reduced finite-strain gradient theories based on selected scalar gradient measures, such as gradients of dilatation \cite{DellIsolaEremeyev2021}.
\end{enumerate}

The present work selects

\begin{equation}
\boxed{
\Psi=\Psi(\mathbf{E},\bm{\mathcal{G}}),
\qquad
\bm{\mathcal{G}}=\Grad\mathbf{E}.
}
\label{eq:sec2_selected_state_dependence}
\end{equation}

This choice is not claimed to be the only possible finite-strain higher-gradient constitutive framework. Rather, it is the most direct total-Lagrangian extension of the validated small-strain Mindlin Form-II formulation because

\begin{equation}
\mathbf{E}\rightarrow\bm{\varepsilon},
\qquad
\Grad\mathbf{E}\rightarrow\Grad\bm{\varepsilon}
\label{eq:sec2_formII_recovery_kinematic}
\end{equation}

in the small-displacement-gradient limit.

\subsection{Comparison of possible Helmholtz free-energy choices}

\subsubsection{Classical hyperelastic free energy}

A purely classical finite-strain hyperelastic material uses

\begin{equation}
\Psi=\Psi(\mathbf{C})
\quad
\text{or equivalently}
\quad
\Psi=\Psi(\mathbf{E}),
\label{eq:sec2_classical_hyperelastic}
\end{equation}

with frame indifference enforced by dependence on objective material measures \cite{Rivlin1948,Ogden1997,Holzapfel2000}. This class is appropriate for ordinary finite elasticity but contains no material length scale and cannot recover the validated strain-gradient benchmark when gradient effects are retained.

\paragraph{Advantage.}
It is standard, objective, and well developed.

\paragraph{Disadvantage for the present work.}
It omits $\Grad\bm{\varepsilon}$ in the linear limit and therefore cannot reproduce Mindlin Form-II strain-gradient elasticity.

\subsubsection{Toupin-type second-gradient free energy}

A general second-gradient elastic solid may be described by

\begin{equation}
\Psi=\Psi(\mathbf{F},\Grad\mathbf{F}),
\label{eq:sec2_toupin_type}
\end{equation}

or by equivalent objective combinations of these variables \cite{Toupin1962,Toupin1964,JaviliDellIsolaSteinmann2013,Bertram2015}. This is a broad and mathematically powerful finite-deformation framework.

\paragraph{Advantage.}
It is general and directly incorporates second gradients of the motion.

\paragraph{Disadvantage for the present work.}
The tensor $\Grad\mathbf{F}$ is not a symmetric strain-gradient tensor. Its linearization is related to second displacement gradients, whereas the validated benchmark is based on the Form-II measure $\Grad\bm{\varepsilon}$. Additional symmetrization or constitutive restrictions would be required to recover the exact Mindlin Form-II structure.

\subsubsection{Green--Lagrange Form-II free energy}

The selected framework uses

\begin{equation}
\Psi=\Psi(\mathbf{E},\Grad\mathbf{E}).
\label{eq:sec2_green_lagrange_formII}
\end{equation}

This is the finite-strain counterpart of the small-strain Form-II structure. Modern finite-deformation strain-gradient formulations use Green--Lagrange strain or nonlinear Green--Lagrange strain relations when extending strain-gradient elasticity to geometrically nonlinear regimes \cite{Beheshti2016,Beheshti2017,TorabiNiiranenAnsari2021,AbaliMullerEremeyev2015}.

\paragraph{Advantage.}
It is objective, total-Lagrangian, displacement-based, and recovers $\Psi(\bm{\varepsilon},\Grad\bm{\varepsilon})$ exactly in the linearized limit.

\paragraph{Disadvantage.}
It is less general than a full $\Psi(\mathbf{F},\Grad\mathbf{F})$ second-gradient theory. This is intentional because the present goal is not to construct the most general higher-gradient solid, but to extend the validated Mindlin Form-II formulation with maximum theoretical continuity.

\subsection{Thermodynamic setting}

The material is assumed to be isothermal and elastic. The Helmholtz free-energy density is taken per unit reference volume:

\begin{equation}
[\Psi]=\mathrm{J\,m^{-3}}=\mathrm{Pa}.
\label{eq:sec2_energy_unit}
\end{equation}

The use of a Helmholtz potential as the constitutive generator follows the standard Coleman--Noll thermodynamic framework for elastic materials \cite{ColemanNoll1963,GurtinFriedAnand2010}.

No dissipative internal variables are introduced. Consequently, the constitutive model is conservative. In a later section, the energetic quantities conjugate to $\mathbf{E}$ and $\bm{\mathcal{G}}$ will be obtained from the free energy. That later step is not performed here.

The thermodynamic assumptions are:

\begin{enumerate}
\item isothermal response;
\item no viscosity;
\item no plasticity;
\item no damage;
\item no independent microdeformation;
\item no dissipative internal variables;
\item homogeneous material parameters in the reference configuration.
\end{enumerate}

\subsection{Independent constitutive state variables}

The local constitutive state is chosen as

\begin{equation}
\mathcal{S}
=
\left\{
\mathbf{E},
\bm{\mathcal{G}}
\right\},
\qquad
\bm{\mathcal{G}}=\Grad\mathbf{E}.
\label{eq:sec2_state_set}
\end{equation}

Although $\bm{\mathcal{G}}$ is kinematically derived from $\mathbf{E}$, it is treated as an independent local argument of the free-energy density in the constitutive function:

\begin{equation}
\Psi=\Psi(\mathbf{E},\bm{\mathcal{G}}).
\label{eq:sec2_local_constitutive_arguments}
\end{equation}

This is standard in gradient theories: the strain and its gradient are locally listed as constitutive arguments, while their compatibility is enforced by the displacement field \cite{Mindlin1964,MindlinEshel1968,AskesAifantis2011}.

The physical meanings are:

\begin{enumerate}
\item $\mathbf{E}$: finite Lagrangian strain; it measures stretch and change of angle in the reference description.

\item $\bm{\mathcal{G}}=\Grad\mathbf{E}$: material gradient of finite strain; it measures spatial inhomogeneity of strain over the reference body and introduces size-dependent effects.
\end{enumerate}

The tensor symmetries inherited from Section~\ref{sec:kinematics} are

\begin{equation}
E_{AB}=E_{BA},
\qquad
\mathcal{G}_{ABC}=\mathcal{G}_{BAC}.
\label{eq:sec2_state_symmetries}
\end{equation}

\subsection{Step-by-step construction of the selected free energy}

\subsubsection{Step 1: General dependence}

The most general local total-Lagrangian Form-II dependence adopted here is

\begin{equation}
\Psi=\Psi(\mathbf{E},\bm{\mathcal{G}}).
\label{eq:sec2_step1_general_energy}
\end{equation}

\subsubsection{Step 2: Expansion about the reference state}

The reference state is

\begin{equation}
\mathbf{E}=\mathbf{0},
\qquad
\bm{\mathcal{G}}=\mathbf{0}.
\label{eq:sec2_reference_state}
\end{equation}

A formal Taylor expansion of $\Psi$ about this state gives

\begin{align}
\Psi(\mathbf{E},\bm{\mathcal{G}})
&=
\Psi_0
+
\mathbb{A}_{AB}E_{AB}
+
\mathbb{B}_{ABC}\mathcal{G}_{ABC}
\nonumber\\
&\quad
+
\frac{1}{2}\mathbb{C}_{ABCD}E_{AB}E_{CD}
+
\mathbb{H}_{ABCPQ}E_{AB}\mathcal{G}_{CPQ}
\nonumber\\
&\quad
+
\frac{1}{2}\mathbb{D}_{ABCPQR}\mathcal{G}_{ABC}\mathcal{G}_{PQR}
+
\text{higher-order terms}.
\label{eq:sec2_taylor_energy}
\end{align}

Here $\mathbb{C}_{ABCD}$ is a fourth-order classical elastic modulus tensor, $\mathbb{D}_{ABCPQR}$ is a sixth-order gradient-elastic modulus tensor, and $\mathbb{H}_{ABCPQ}$ is a possible fifth-order coupling tensor.

\subsubsection{Step 3: Reference normalization}

The arbitrary zero of energy is selected such that

\begin{equation}
\Psi_0=0.
\label{eq:sec2_zero_energy}
\end{equation}

The reference configuration is assumed to be initially free of energetic bias; therefore the terms linear in $\mathbf{E}$ and $\bm{\mathcal{G}}$ are omitted:

\begin{equation}
\mathbb{A}_{AB}=0,
\qquad
\mathbb{B}_{ABC}=0.
\label{eq:sec2_no_linear_terms}
\end{equation}

This is the constitutive analogue of selecting an undeformed, homogeneous reference state. (The vanishing of $\mathbb{B}_{ABC}$, linear in the strain gradient, is also required by centrosymmetry; see Step~4.)

\subsubsection{Step 4: Centrosymmetry}

For a centrosymmetric material, the energy must be invariant under inversion of material directions. The strain $\mathbf{E}$ is even under this inversion, while the material gradient $\bm{\mathcal{G}}=\Grad\mathbf{E}$ is odd in the gradient direction. Hence terms containing an odd number of strain-gradient factors are excluded. The mixed quadratic term

\begin{equation}
\mathbb{H}_{ABCPQ}E_{AB}\mathcal{G}_{CPQ}
\label{eq:sec2_mixed_coupling}
\end{equation}

is therefore omitted for an isotropic centrosymmetric material. This is consistent with the usual centrosymmetric reduction in Mindlin-type strain-gradient elasticity \cite{Mindlin1964,MindlinEshel1968,LazarPo2018}.

\subsubsection{Step 5: Quadratic truncation}

The validated 2022 benchmark is a small-strain linear Mindlin Form-II formulation. To recover it exactly in the small-strain limit, the finite-strain extension must retain the same quadratic structure in the strain and strain-gradient measures. Therefore, the selected energy is

\begin{equation}
\boxed{
\Psi(\mathbf{E},\bm{\mathcal{G}})
=
\frac{1}{2}\mathbb{C}_{ABCD}E_{AB}E_{CD}
+
\frac{1}{2}\mathbb{D}_{ABCPQR}\mathcal{G}_{ABC}\mathcal{G}_{PQR}.
}
\label{eq:sec2_selected_tensor_energy}
\end{equation}

This is a geometrically nonlinear, materially linear-in-$\mathbf{E}$ Saint-Venant--Kirchhoff-type strain-gradient model. The nonlinearity enters through the kinematic dependence of $\mathbf{E}$ and $\bm{\mathcal{G}}$ on the displacement field.

\subsection{Classical finite-strain part}

For an isotropic material, the fourth-order elasticity tensor is

\begin{equation}
\mathbb{C}_{ABCD}
=
\lambda\,\delta_{AB}\delta_{CD}
+
\mu\left(
\delta_{AC}\delta_{BD}
+
\delta_{AD}\delta_{BC}
\right).
\label{eq:sec2_isotropic_C_tensor}
\end{equation}

Substitution into the first term of \eqref{eq:sec2_selected_tensor_energy} gives

\begin{align}
\Psi_{\mathrm{SVK}}(\mathbf{E})
&=
\frac{1}{2}
\left[
\lambda\,\delta_{AB}\delta_{CD}
+
\mu\left(
\delta_{AC}\delta_{BD}
+
\delta_{AD}\delta_{BC}
\right)
\right]
E_{AB}E_{CD}
\nonumber\\
&=
\frac{\lambda}{2}E_{AA}E_{CC}
+
\frac{\mu}{2}
\left(
E_{AB}E_{AB}
+
E_{AB}E_{BA}
\right).
\label{eq:sec2_svk_intermediate}
\end{align}

Since $E_{AB}=E_{BA}$,

\begin{equation}
E_{AB}E_{BA}=E_{AB}E_{AB}.
\label{eq:sec2_E_sym_contraction}
\end{equation}

Therefore,

\begin{equation}
\boxed{
\Psi_{\mathrm{SVK}}(\mathbf{E})
=
\frac{\lambda}{2}\left(\tr\mathbf{E}\right)^2
+
\mu\,\mathbf{E}:\mathbf{E}.
}
\label{eq:sec2_svk_energy}
\end{equation}

This is the Saint-Venant--Kirchhoff finite-strain elastic energy. It is objective because it depends on $\mathbf{E}$, which is objective. It is selected because it recovers linear isotropic elasticity exactly under small-strain linearization.

\subsection{Mindlin Form-II gradient contribution}

The gradient contribution is

\begin{equation}
\boxed{
\Psi_{\mathrm{g}}(\bm{\mathcal{G}})
=
\frac{1}{2}\mathbb{D}_{ABCPQR}\mathcal{G}_{ABC}\mathcal{G}_{PQR}.
}
\label{eq:sec2_gradient_energy_tensor}
\end{equation}

The sixth-order tensor $\mathbb{D}_{ABCPQR}$ satisfies the index symmetries

\begin{equation}
\mathbb{D}_{ABCPQR}
=
\mathbb{D}_{BACPQR}
=
\mathbb{D}_{ABCQPR}
=
\mathbb{D}_{PQRABC}.
\label{eq:sec2_D_symmetries}
\end{equation}

The first two symmetries are inherited from the symmetry of $\mathcal{G}_{ABC}$ in its first two indices; the last symmetry follows from the scalar quadratic form.

For an isotropic centrosymmetric Mindlin Form-II material, the gradient contribution may be expressed using five independent scalar invariants of $\bm{\mathcal{G}}$:

\begin{align}
\Psi_{\mathrm{g}}(\bm{\mathcal{G}})
&=
a_1\,\mathcal{G}_{ABB}\mathcal{G}_{ACC}
+
a_2\,\mathcal{G}_{AAC}\mathcal{G}_{CBB}
+
a_3\,\mathcal{G}_{AAC}\mathcal{G}_{BBC}
\nonumber\\
&\quad
+
a_4\,\mathcal{G}_{ABC}\mathcal{G}_{ABC}
+
a_5\,\mathcal{G}_{ABC}\mathcal{G}_{CBA}.
\label{eq:sec2_gradient_energy_invariants}
\end{align}

This invariant representation is the finite-strain counterpart of the common isotropic Mindlin Form-II quadratic energy written in terms of $\varepsilon_{ij,k}$ \cite{Mindlin1964,MindlinEshel1968,LazarPo2018}. The ordering and naming of the five coefficients may vary across authors; therefore, for the finite element formulation it is safest to regard $\mathbb{D}_{ABCPQR}$ as the primary modulus tensor and treat $a_1,\ldots,a_5$ as its isotropic invariant parameters.

\subsection{Complete selected Helmholtz free-energy density}

Combining \eqref{eq:sec2_svk_energy} and \eqref{eq:sec2_gradient_energy_tensor}, the selected free-energy density is

\begin{equation}
\boxed{
\Psi(\mathbf{E},\bm{\mathcal{G}})
=
\frac{\lambda}{2}\left(\tr\mathbf{E}\right)^2
+
\mu\,\mathbf{E}:\mathbf{E}
+
\frac{1}{2}
\mathbb{D}_{ABCPQR}
\mathcal{G}_{ABC}\mathcal{G}_{PQR}.
}
\label{eq:sec2_complete_energy_tensor}
\end{equation}

Equivalently, in the five-parameter isotropic invariant basis,

\begin{align}
\Psi(\mathbf{E},\bm{\mathcal{G}})
&=
\frac{\lambda}{2}\left(\tr\mathbf{E}\right)^2
+
\mu\,E_{AB}E_{AB}
\nonumber\\
&\quad
+
a_1\,\mathcal{G}_{ABB}\mathcal{G}_{ACC}
+
a_2\,\mathcal{G}_{AAC}\mathcal{G}_{CBB}
+
a_3\,\mathcal{G}_{AAC}\mathcal{G}_{BBC}
\nonumber\\
&\quad
+
a_4\,\mathcal{G}_{ABC}\mathcal{G}_{ABC}
+
a_5\,\mathcal{G}_{ABC}\mathcal{G}_{CBA}.
\label{eq:sec2_complete_energy_invariants}
\end{align}

\subsection{Physical meaning of the constitutive constants}

\begin{longtable}[c]{>{\raggedright\arraybackslash}p{0.15\linewidth}
>{\raggedright\arraybackslash}p{0.18\linewidth}
>{\raggedright\arraybackslash}p{0.18\linewidth}
>{\raggedright\arraybackslash}p{0.39\linewidth}}
\caption{Constitutive constants in the selected finite-strain Mindlin Form-II model.}
\label{tab:sec2_constitutive_constants}\\
\toprule
Parameter & SI unit & Type & Physical role \\
\midrule
\endfirsthead
\toprule
Parameter & SI unit & Type & Physical role \\
\midrule
\endhead
$\lambda$ & Pa & Lam\'{e} modulus & Controls volumetric part of the classical finite-strain response. \\
$\mu$ & Pa & Lam\'{e} modulus & Controls distortional/shear part of the classical finite-strain response. \\
$a_1,\ldots,a_5$ & N, equivalently Pa\,m$^2$ & Gradient-elastic constants & Control the energetic penalty associated with strain inhomogeneity. They introduce intrinsic length-scale effects and determine how strongly the material resists gradients of finite strain. \\
$\mathbb{D}_{ABCPQR}$ & Pa\,m$^2$ & Sixth-order gradient modulus tensor & Tensorial representation of the five isotropic gradient constants; acts on tensors symmetric in the first two indices. \\
\bottomrule
\end{longtable}

Because $\mathbf{E}$ is dimensionless and $\bm{\mathcal{G}}$ has dimension $L^{-1}$,

\begin{equation}
[\mathbf{E}]=1,
\qquad
[\bm{\mathcal{G}}]=L^{-1}.
\label{eq:sec2_state_dimensions}
\end{equation}

The classical energy term has units

\begin{equation}
[\lambda E^2]=[\mu E^2]=\mathrm{Pa}.
\label{eq:sec2_classical_units}
\end{equation}

The gradient energy term has units

\begin{equation}
[a_i\mathcal{G}^2]
=
(\mathrm{Pa}\,m^2)(m^{-2})
=
\mathrm{Pa}.
\label{eq:sec2_gradient_units}
\end{equation}

Thus every term in $\Psi$ has the correct unit of energy per reference volume.

\subsection{Objectivity and frame indifference}

Under a superposed rigid-body motion,

\begin{equation}
\widehat{\mathbf{x}}=\mathbf{Q}(t)\mathbf{x}+\mathbf{c}(t),
\qquad
\mathbf{Q}^{T}\mathbf{Q}=\mathbf{I},
\qquad
\det\mathbf{Q}=1,
\label{eq:sec2_superposed_motion}
\end{equation}

Section~\ref{sec:recovery} established that

\begin{equation}
\widehat{\mathbf{E}}=\mathbf{E},
\qquad
\widehat{\bm{\mathcal{G}}}=\bm{\mathcal{G}}.
\label{eq:sec2_E_G_objective}
\end{equation}

Therefore,

\begin{equation}
\widehat{\Psi}
=
\Psi(\widehat{\mathbf{E}},\widehat{\bm{\mathcal{G}}})
=
\Psi(\mathbf{E},\bm{\mathcal{G}}).
\label{eq:sec2_energy_objective}
\end{equation}

The selected free energy is frame-indifferent because it is constructed solely from objective material kinematic measures. This is consistent with the standard finite-elasticity requirement of material frame indifference \cite{Holzapfel2000,BonetWood2008,GurtinFriedAnand2010}.

\subsection{Material isotropy}

Material isotropy is imposed with respect to rotations of the reference basis. For every proper orthogonal material tensor $\mathbf{Q}_0\in SO(3)$, isotropy requires

\begin{equation}
\Psi(\mathbf{E},\bm{\mathcal{G}})
=
\Psi(\mathbf{Q}_0^{T}\mathbf{E}\mathbf{Q}_0,\bm{\mathcal{G}}^{\mathbf{Q}_0}),
\label{eq:sec2_material_isotropy}
\end{equation}

where

\begin{equation}
\mathcal{G}^{\mathbf{Q}_0}_{ABC}
=
Q_{0PA}Q_{0QB}Q_{0RC}\mathcal{G}_{PQR}.
\label{eq:sec2_G_material_rotation}
\end{equation}

The classical Saint-Venant--Kirchhoff part satisfies this requirement because it depends only on the invariants $\tr\mathbf{E}$ and $\mathbf{E}:\mathbf{E}$. The gradient contribution satisfies it when $\mathbb{D}_{ABCPQR}$ is an isotropic sixth-order tensor, equivalently when the invariant representation \eqref{eq:sec2_gradient_energy_invariants} is used.

\subsection{Thermodynamic admissibility}

The reduced isothermal thermodynamic setting requires that the stored energy be a state function of the selected variables. Since no dissipative internal variables are introduced, the constitutive response is elastic. In a later section, the generalized energetic quantities conjugate to $\mathbf{E}$ and $\bm{\mathcal{G}}$ will be obtained from $\Psi$. With that later energetic construction, the purely elastic part of the dissipation is identically zero, which is the standard Coleman--Noll admissibility mechanism for elastic materials \cite{ColemanNoll1963,GurtinFriedAnand2010}.

At the present stage, the thermodynamic admissibility requirements imposed on $\Psi$ are:

\begin{enumerate}
\item $\Psi$ is a scalar state function;
\item $\Psi$ is objective;
\item $\Psi(\mathbf{0},\mathbf{0})=0$;
\item $\Psi(\mathbf{E},\bm{\mathcal{G}})\ge 0$ locally near the reference configuration;
\item $\Psi$ contains no dissipative history dependence.
\end{enumerate}

\subsection{Positive definiteness, stability, and convexity}

The classical quadratic form

\begin{equation}
\Psi_{\mathrm{SVK}}
=
\frac{\lambda}{2}(\tr\mathbf{E})^2+\mu\,\mathbf{E}:\mathbf{E}
\label{eq:sec2_classical_positive_form}
\end{equation}

is positive definite in $\mathbf{E}$ if

\begin{equation}
\mu>0,
\qquad
3\lambda+2\mu>0.
\label{eq:sec2_lame_positive_conditions}
\end{equation}

The gradient quadratic form is positive definite if

\begin{equation}
\mathbb{D}_{ABCPQR}\mathcal{G}_{ABC}\mathcal{G}_{PQR}>0
\quad
\text{for all nonzero }
\bm{\mathcal{G}}
\text{ satisfying }
\mathcal{G}_{ABC}=\mathcal{G}_{BAC}.
\label{eq:sec2_D_positive_definite}
\end{equation}

In the five-constant invariant representation, this requirement imposes restrictions on the admissible combinations of $a_1,\ldots,a_5$. These restrictions are not equivalent to the simple condition $a_i>0$ for all $i$. Strong ellipticity and stability conditions for Toupin--Mindlin first strain-gradient elasticity have been studied in the modern literature \cite{EremeyevLazar2022}. Therefore, admissible numerical choices of $a_1,\ldots,a_5$ must be checked at the level of the full quadratic form or through known ellipticity conditions.

The selected Saint-Venant--Kirchhoff part is quadratic and convex in $\mathbf{E}$, but it is not globally convex as a function of $\mathbf{F}$. This is a known limitation of Saint-Venant--Kirchhoff finite elasticity and is related to the broader distinction between convexity, rank-one convexity, quasiconvexity, and polyconvexity in nonlinear elasticity \cite{Ball1976}. For the present research program this is acceptable because the objective is a geometrically nonlinear extension that exactly recovers the validated small-strain benchmark, not a globally stable model for extreme finite stretches.

Thus, the stability statement adopted here is local:

\begin{equation}
\Psi
\text{ is required to be locally positive definite about }
(\mathbf{E},\bm{\mathcal{G}})=(\mathbf{0},\mathbf{0}).
\label{eq:sec2_local_stability_statement}
\end{equation}

\subsection{Small-strain recovery}

From Section~\ref{sec:kinematics},

\begin{equation}
\mathbf{E}
=
\bm{\varepsilon}
+
O(\|\Grad\mathbf{u}\|^2),
\label{eq:sec2_E_small_recovery}
\end{equation}

and

\begin{equation}
\bm{\mathcal{G}}
=
\Grad\mathbf{E}
=
\Grad\bm{\varepsilon}
+
O(\|\Grad\mathbf{u}\|\,\|\Grad^2\mathbf{u}\|).
\label{eq:sec2_G_small_recovery}
\end{equation}

Substituting these into \eqref{eq:sec2_complete_energy_tensor} gives, to quadratic order in the small displacement gradient,

\begin{equation}
\Psi
\rightarrow
\frac{\lambda}{2}(\tr\bm{\varepsilon})^2
+
\mu\,\bm{\varepsilon}:\bm{\varepsilon}
+
\frac{1}{2}\mathbb{D}_{ABCPQR}
\varepsilon_{AB,C}\varepsilon_{PQ,R}.
\label{eq:sec2_small_strain_energy_tensor}
\end{equation}

Equivalently, in the isotropic invariant representation,

\begin{align}
\Psi
\rightarrow
&
\frac{\lambda}{2}(\varepsilon_{AA})^2
+
\mu\,\varepsilon_{AB}\varepsilon_{AB}
\nonumber\\
&
+ a_1\,\varepsilon_{AB,B}\varepsilon_{AC,C}
+ a_2\,\varepsilon_{AA,C}\varepsilon_{CB,B}
+ a_3\,\varepsilon_{AA,C}\varepsilon_{BB,C}
\nonumber\\
&
+ a_4\,\varepsilon_{AB,C}\varepsilon_{AB,C}
+ a_5\,\varepsilon_{AB,C}\varepsilon_{CB,A}.
\label{eq:sec2_small_strain_energy_invariants}
\end{align}

This is the small-strain isotropic Mindlin Form-II energy, provided the same invariant convention and the same constants as the validated 2022 benchmark are used.

\subsection{Recovery of the validated 2022 benchmark formulation}

The validated 2022 benchmark is recovered under the following limiting assumptions:

\begin{enumerate}
\item displacement gradients are small;
\item the Green--Lagrange strain is replaced by its linear part, $\mathbf{E}\rightarrow\bm{\varepsilon}$;
\item the finite strain-gradient tensor is replaced by $\bm{\mathcal{G}}\rightarrow\Grad\bm{\varepsilon}$;
\item the same Lam\'{e} constants $\lambda,\mu$ are used;
\item the same Mindlin Form-II gradient constants or the same sixth-order tensor $\mathbb{D}$ are used;
\item the same boundary conditions and benchmark geometry are retained.
\end{enumerate}

Then the finite-strain free energy \eqref{eq:sec2_complete_energy_tensor} reduces exactly to the energy density underlying the validated small-strain Mindlin Form-II formulation:

\begin{equation}
\boxed{
\Psi(\mathbf{E},\Grad\mathbf{E})
\longrightarrow
\Psi_{\mathrm{2022}}(\bm{\varepsilon},\Grad\bm{\varepsilon}).
}
\label{eq:sec2_2022_benchmark_recovery}
\end{equation}

This recovery is the principal reason for selecting the Saint-Venant--Kirchhoff--Mindlin Form-II free-energy structure rather than a more general hyperelastic or full second-gradient constitutive model.

\subsection{Reviewer-level discussion}

\subsubsection*{Objection 1: Why use Saint-Venant--Kirchhoff elasticity instead of a modern hyperelastic model?}

A neo-Hookean, Mooney--Rivlin, Ogden, or polyconvex hyperelastic model would be more appropriate for extreme elastic stretches \cite{Rivlin1948,Ogden1997,Ball1976}. However, the present research objective is not to introduce a new nonlinear material law. The goal is to extend the already validated displacement-based Mindlin Form-II strain-gradient formulation to finite kinematics while preserving the benchmark. The Saint-Venant--Kirchhoff choice is the minimal objective finite-strain extension that recovers linear elasticity exactly.

\subsubsection*{Objection 2: Is $\Grad\mathbf{E}$ the only possible finite-strain strain-gradient measure?}

No. General finite second-gradient theories may use $\Grad\mathbf{F}$ or related objective variables \cite{Toupin1962,Toupin1964,JaviliDellIsolaSteinmann2013,Bertram2015}. The present selection is not based on uniqueness. It is based on compatibility with Mindlin Form II and exact recovery of $\Grad\bm{\varepsilon}$ in the small-strain limit.

\subsubsection*{Objection 3: Are five gradient constants practical?}

The full isotropic Mindlin Form-II theory contains five gradient-elastic constants \cite{Mindlin1964,MindlinEshel1968,LazarPo2018}. Their experimental determination is difficult. However, using the full five-constant tensor preserves the most general isotropic Form-II structure. Reduced one- or two-length-scale models can later be obtained as special cases, but should not be imposed at the foundational level unless required by the benchmark.

\subsubsection*{Objection 4: Is the model globally stable?}

The Saint-Venant--Kirchhoff energy is not globally convex in $\mathbf{F}$. Therefore, the present framework should be interpreted as locally stable around the reference configuration and suitable for geometrically nonlinear finite-rotation or moderate-stretch analysis. Global polyconvexity is not claimed.

\subsubsection*{Objection 5: Does the framework introduce new physics beyond the validated benchmark?}

Yes, but only through finite-strain kinematics. The constitutive structure is deliberately chosen so that the benchmark is recovered exactly in the small-strain limit. Therefore, any deviation from the validated benchmark at finite deformation is attributable to geometrical nonlinearity and finite strain-gradient effects, not to an unrelated constitutive modification.

\subsection{Final constitutive statement}

The constitutive framework adopted for the remainder of the formulation is:

\begin{equation}
\boxed{
\Psi(\mathbf{E},\bm{\mathcal{G}})
=
\frac{\lambda}{2}\left(\tr\mathbf{E}\right)^2
+
\mu\,\mathbf{E}:\mathbf{E}
+
\frac{1}{2}
\mathbb{D}_{ABCPQR}
\mathcal{G}_{ABC}\mathcal{G}_{PQR},
\qquad
\bm{\mathcal{G}}=\Grad\mathbf{E}.
}
\label{eq:sec2_final_constitutive_framework}
\end{equation}

The model is objective, isothermal, elastic, total-Lagrangian, isotropic, centrosymmetric, displacement-based, locally stable under admissible material parameters, and exactly reduces to the validated small-strain Mindlin Form-II benchmark formulation.

% =====================================================================
\section{Variational Formulation -- Research Notes}
\label{sec:variational}
% =====================================================================

\paragraph{Scope and status.}

This section introduces \emph{only} the variational (action) principle: Hamilton's principle, kinetic energy, the total Helmholtz free-energy functional, external virtual work, inertia virtual work, the total variational (action) functional, and the final unreduced variational statement, together with the complete derivation of the strong form by integration by parts. Integration by parts, the strong form, natural/essential boundary conditions, the weak form, finite-element matrices, and tangent operators are developed here to closure of the strong form. Every algebraic step, identity, substitution, variation, and check is recorded below.

\paragraph{Notation retained from Sections~\ref{sec:kinematics} and~\ref{sec:constitutive}.}

Motion $\boldsymbol{\varphi}(\mathbf{X},t)$; displacement $\mathbf{u}=\boldsymbol{\varphi}-\mathbf{X}$; deformation gradient $\mathbf{F}=\nabla_{0}\boldsymbol{\varphi}=\mathbf{I}+\nabla_{0}\mathbf{u}$; $J=\det\mathbf{F}>0$; Green--Lagrange strain $\mathbf{E}=\tfrac12(\mathbf{F}^{T}\mathbf{F}-\mathbf{I})$ with components $E_{IJ}=\tfrac12(F_{kI}F_{kJ}-\delta_{IJ})$; strain gradient $\bm{\mathcal{G}}=\mathrm{Grad}\,\mathbf{E}$ with $\mathcal{G}_{ABC}=E_{AB,C}=\mathcal{G}_{BAC}$; polar rotation $\mathbf{R}=\mathrm{polar}(\mathbf{F})$ (kinematic variable, listed in Section~\ref{sec:kinematics} item~14, not a constitutive argument). Helmholtz free-energy density (eq:psi of Section~\ref{sec:constitutive}):

\begin{align}
\psi(\mathbf{E},\bm{\mathcal{G}})
&= \tfrac12\lambda(\mathrm{tr}\,\mathbf{E})^{2}
   +\mu\,\mathbf{E}:\mathbf{E}
   +\tfrac12\mathbb{D}_{ABCPQR}\,\mathcal{G}_{ABC}\mathcal{G}_{PQR},
\qquad \mathcal{G}_{ABC}=E_{AB,C}.
\label{rn:psi}
\end{align}

All arguments are material tensors. Material indices $I,J,K,A,B,C\in\{1,2,3\}$; spatial indices $i,j,k\in\{1,2,3\}$.

% =====================================================================
\subsection{A. Preliminaries, scope, and compatibility audit}
% =====================================================================

\textbf{A.1 New symbols introduced in this section only.}

\begin{itemize}
\item $\rho_{0}(\mathbf{X})$: reference mass density $[\mathrm{kg\,m^{-3}}]$ (constant for the homogeneous solid).

\item $\mathbf{f}_{0}(\mathbf{X},t)$: prescribed body force per unit reference volume $[\mathrm{N\,m^{-3}}]$.

\item $\bar{\mathbf{T}}_{0}(\mathbf{X},t)$: prescribed nominal traction per unit reference area $[\mathrm{N\,m^{-2}}]$ on $\partial\mathcal{B}_{0t}$.

\item $\mathbf{N}$: outward reference unit normal on $\partial\mathcal{B}_{0}$.

\item $t_{1},t_{2}$: time endpoints of the observation interval; $\delta\boldsymbol{\varphi}(t_{1})=\delta\boldsymbol{\varphi}(t_{2})=\mathbf{0}$.

\item $\mathcal{T},\Pi_{\mathrm{int}},\delta\mathcal{W}_{\mathrm{ext}},\delta\mathcal{W}_{\mathrm{iner}},\mathcal{L},\mathcal{A}$: scalar functionals/actions.
\end{itemize}

\textbf{A.2 Compatibility audit.} Every kinematic and constitutive object equals that of Sections~\ref{sec:kinematics}--\ref{sec:constitutive}: $\mathbf{E},\bm{\mathcal{G}}$ are reused verbatim (eq:n\_E, eq:n\_Gdef / sec:Gdef; rn:psi, eq:sec2\_final\_constitutive\_framework). No new constitutive variables or material parameters are introduced: the modulus object is the sixth-order tensor $\mathbb{D}_{ABCPQR}$ (equivalently the five constants $a_{1},\ldots,a_{5}$ of Section~\ref{sec:constitutive}, eq:sec2\_gradient\_energy\_invariants), and the energy depends only on $\{\mathbf{E},\bm{\mathcal{G}}\}$. Therefore the mandatory recovery rule (every new formulation must reduce exactly to the validated 2022 benchmark) can be enforced at the variational level using the already-verified small-strain limits of Sections~\ref{sec:kinematics}--\ref{sec:constitutive}.

% =====================================================================
\subsection{B. Hamilton's principle: statement and literature}
% =====================================================================

\textbf{B.1 General principle.} For a holonomic system with Lagrangian $\mathcal{L}=T-V$ (kinetic minus potential energy) the actual path renders the action stationary:

\begin{equation}
\delta\int_{t_{1}}^{t_{2}}\mathcal{L}\,\mathrm{d}t = 0,
\qquad
\delta\boldsymbol{\varphi}(t_{1})=\delta\boldsymbol{\varphi}(t_{2})=\mathbf{0}.
\label{rn:Ham_conservative}
\end{equation}

This is Hamilton's principle \cite{Hamilton1834,Goldstein2002}. With prescribed (non-conservative) external loads the principle is augmented by their virtual work, giving the Lagrange--d'Alembert form \cite{Goldstein2002,Gurtin1981}:

\begin{equation}
\delta\int_{t_{1}}^{t_{2}}\mathcal{L}\,\mathrm{d}t
+ \int_{t_{1}}^{t_{2}}\delta\mathcal{W}_{\mathrm{ext}}\,\mathrm{d}t = 0.
\label{rn:Ham_noncons}
\end{equation}

\textbf{B.2 Continuum specialisation.} In finite-strain continuum mechanics the potential $V$ is the internally stored Helmholtz free energy $\Pi_{\mathrm{int}}$, and the admissible variations are virtual displacement fields satisfying the kinematic (essential) boundary conditions \cite{MarsdenHughes1983,Gurtin1981}. For micro-structured / strain-gradient media the same principle applies with the internal energy augmented by gradient terms; the virtual-power / virtual-work foundation is due to Germain \cite{Germain1973} and underlies Mindlin Form-II \cite{MindlinEshel1968,Toupin1962}. Maugin's method of virtual power \cite{Maugin1980} explicitly extends the power principle to media with higher-grade (gradient) variables, supporting the present use of a single variational statement that later accommodates the $\delta\bm{\mathcal{G}}$ conjugate stress.

% =====================================================================
\subsection{C. Kinetic energy -- definition, derivation, dimensions, objectivity}
% =====================================================================

\textbf{C.1 Material velocity.} Hold $\mathbf{X}$ fixed and differentiate the motion:
$\dot{\boldsymbol{\varphi}} = \partial\boldsymbol{\varphi}(\mathbf{X},t)/\partial t|_{\mathbf{X}}$.
Since $\mathbf{u}=\boldsymbol{\varphi}-\mathbf{X}$ and $\mathbf{X}$ is time-independent,
$\dot{\mathbf{u}} = \dot{\boldsymbol{\varphi}}$. Hence the velocity field is $\dot{\boldsymbol{\varphi}}=\dot{\mathbf{u}}$.

\textbf{C.2 Definition (reference-volume integral).} Consistency with the Total-Lagrangian description (assumption A4 of Section~\ref{sec:kinematics}) requires all volume integrals to be taken over the fixed reference domain $\mathcal{B}_{0}$ with the reference density $\rho_{0}$:

\begin{equation}
\mathcal{T}(t)
= \int_{\mathcal{B}_{0}}
  \tfrac12\,\rho_{0}\,\dot{\boldsymbol{\varphi}}\cdot\dot{\boldsymbol{\varphi}}\,
  \mathrm{d}V_{0}
= \int_{\mathcal{B}_{0}}
  \tfrac12\,\rho_{0}\,\dot{\mathbf{u}}\cdot\dot{\mathbf{u}}\,
  \mathrm{d}V_{0}.
\label{rn:T}
\end{equation}

This is the standard TL kinetic-energy functional \cite{BonetWood2008,Bathe2006,Holzapfel2000}.

\textbf{C.3 Dimensional check.}
$[\rho_{0}]=\mathrm{kg\,m^{-3}}$, $[\dot{\boldsymbol{\varphi}}^{2}]=\mathrm{m^{2}\,s^{-2}}$, $[\mathrm{d}V_{0}]=\mathrm{m^{3}}$.
$[\mathcal{T}]=\mathrm{kg\,m^{-3}}\cdot\mathrm{m^{2}\,s^{-2}}\cdot\mathrm{m^{3}}=\mathrm{kg\,m^{2}\,s^{-2}}=\mathrm{J}$. \checkmark

\textbf{C.4 Positivity.} $\mathcal{T}\ge 0$ with equality iff $\dot{\boldsymbol{\varphi}}=\mathbf{0}$ a.e. \checkmark

\textbf{C.5 Objectivity / frame indifference.} Under a superposed rigid-body motion $\bar{\mathbf{x}}=\mathbf{Q}(t)\mathbf{x}+\mathbf{c}(t)$, $\mathbf{Q}\in SO(3)$, the material velocity transforms as $\dot{\bar{\boldsymbol{\varphi}}}=\mathbf{Q}\dot{\boldsymbol{\varphi}}+\dot{\mathbf{Q}}\boldsymbol{\varphi}+\dot{\mathbf{c}}$. The kinetic energy is therefore not invariant under a \emph{time-dependent} rotation $\mathbf{Q}(t)$. However, Hamilton's principle is stated in an inertial frame, where the standard kinetic energy is used. Under a Galilean boost $\mathbf{c}(t)=\mathbf{v}_{0}t+\mathbf{c}_{0}$ the Lagrangian changes by a total time derivative, so the stationarity condition (and hence the Euler--Lagrange equations) is unchanged \cite{Goldstein2002,MarsdenHughes1983}. The constitutive part of the action is fully objective (material tensors, Section~K); the combination yields objective field equations. Non-inertial frames would require additional inertial body forces, which are outside the present (inertial-frame) scope.

% =====================================================================
\subsection{D. Total Helmholtz free-energy functional}
% =====================================================================

\textbf{D.1 Definition.} Using the density $\psi$ of eq:psi,

\begin{equation}
\Pi_{\mathrm{int}}[\mathbf{u}]
= \int_{\mathcal{B}_{0}}
  \psi\bigl(\mathbf{E}(\mathbf{u}),\,
           \bm{\mathcal{G}}(\mathbf{u})\bigr)\,
  \mathrm{d}V_{0}.
\label{rn:Pi}
\end{equation}

\textbf{D.2 Explicit dependence chain (every substitution).}

\begin{enumerate}
\item $\boldsymbol{\varphi}=\mathbf{X}+\mathbf{u}$ (eq:n\_disp, Section~\ref{sec:kinematics}).

\item $\mathbf{F}=\mathbf{I}+\nabla_{0}\mathbf{u}$ (eq:n\_Fdec\_t).

\item $\mathbf{E}=\tfrac12(\mathbf{F}^{T}\mathbf{F}-\mathbf{I})$ (eq:n\_E).

\item $\bm{\mathcal{G}}=\nabla_{0}\mathbf{E}$, i.e.\ $\mathcal{G}_{ABC}=\partial E_{AB}/\partial X_{C}$ (eq:n\_Gdef).
\end{enumerate}

Thus $\Pi_{\mathrm{int}}$ is a functional of the single primal field $\mathbf{u}$ (displacement-based, assumption A2 of Section~\ref{sec:kinematics}). No independent microrotation degree of freedom is introduced.

\textbf{D.3 Fixed domain.} Because the description is Total Lagrangian, $\mathcal{B}_{0}$ and $\mathrm{d}V_{0}$ are independent of $\mathbf{u}$ and $t$; the variation of the integral is the integral of the variation (no Reynolds/domain term). This is essential for exact benchmark recovery and is the reason the TL description was mandated in Section~\ref{sec:kinematics}.

\textbf{D.4 Dimensional / sign checks.} $[\psi]=\mathrm{J\,m^{-3}}=[\mathrm{Pa}]$, $[\mathrm{d}V_{0}]=\mathrm{m^{3}}\Rightarrow[\Pi_{\mathrm{int}}]=\mathrm{J}$. By the construction of Section~\ref{sec:constitutive}, $\psi\ge 0$ and $\psi=0$ at $\mathbf{E}=\bm{\mathcal{G}}=\mathbf{0}$ (stress-free reference), so $\Pi_{\mathrm{int}}\ge 0$ with $\Pi_{\mathrm{int}}=0$ at the reference state. \checkmark

% =====================================================================
\subsection{E. External virtual work}
% =====================================================================

\textbf{E.1 Definition.} For a virtual displacement $\delta\mathbf{u}=\delta\boldsymbol{\varphi}$,

\begin{equation}
\delta\mathcal{W}_{\mathrm{ext}}
= \int_{\mathcal{B}_{0}}
    \mathbf{f}_{0}\cdot\delta\mathbf{u}\,\mathrm{d}V_{0}
+ \int_{\partial\mathcal{B}_{0t}}
    \bar{\mathbf{T}}_{0}\cdot\delta\mathbf{u}\,\mathrm{d}A_{0}.
\label{rn:Wext}
\end{equation}

$\mathbf{f}_{0}$ is the body force per unit reference volume; $\bar{\mathbf{T}}_{0}$ is the nominal (first Piola--Kirchhoff) traction per unit reference area on the Neumann part $\partial\mathcal{B}_{0t}$.

\textbf{E.2 Nanson's relation (area-element pull-back).} Let $\mathbf{n}$ be the current outward normal and $\mathrm{d}a$ the current area element; $\mathbf{N},\mathrm{d}A_{0}$ the reference counterparts. Nanson's formula \cite{Holzapfel2000,BonetWood2008} is

\begin{equation}
\mathrm{d}\mathbf{a}=J\,\mathbf{F}^{-T}\mathbf{N}\,\mathrm{d}A_{0},
\qquad\text{i.e.}\qquad
\mathrm{d}a\,\mathbf{n}=J\,\mathbf{F}^{-T}\mathbf{N}\,\mathrm{d}A_{0}.
\label{rn:nanson}
\end{equation}

The nominal traction is defined so that the same force acts on corresponding area elements:
$\bar{\mathbf{T}}_{0}\,\mathrm{d}A_{0}=\bar{\mathbf{t}}\,\mathrm{d}a$,
where $\bar{\mathbf{t}}$ is the Cauchy (current) traction. (Equivalently, the natural boundary condition that will later emerge from the principle is $\bar{\mathbf{T}}_{0}=\mathbf{P}\mathbf{N}$, with $\mathbf{P}$ the first Piola--Kirchhoff stress; we do \emph{not} derive this here.) If a Cauchy traction $\bar{\mathbf{t}}$ is prescribed instead, its contribution is $\int_{\partial\mathcal{B}_{t}}\bar{\mathbf{t}}\cdot\delta\mathbf{u}\,\mathrm{d}a$, which is converted to the reference integral via \eqref{rn:nanson}. We deliberately avoid the scalar-vector identity $\bar{\mathbf{T}}_{0}=J|\mathbf{F}^{-T}\mathbf{N}|\bar{\mathbf{t}}$ in the main statement and keep the exact Nanson vector-area form \eqref{rn:nanson}.

\textbf{E.3 Dimensional check.} Body term: $[\mathbf{f}_{0}]=\mathrm{N\,m^{-3}}$, $[\delta\mathbf{u}]=\mathrm{m}$, $[\mathrm{d}V_{0}]=\mathrm{m^{3}}\Rightarrow[\cdot]=\mathrm{N\,m}=\mathrm{J}$. Surface term: $[\bar{\mathbf{T}}_{0}]=\mathrm{N\,m^{-2}}$, $[\delta\mathbf{u}]=\mathrm{m}$, $[\mathrm{d}A_{0}]=\mathrm{m^{2}}\Rightarrow[\cdot]=\mathrm{J}$. \checkmark

\textbf{E.4 Small-strain limit of the load terms.} As $\|\nabla_{0}\mathbf{u}\|\to0$: $J\to1$; the nominal traction $\bar{\mathbf{T}}_{0}\to\bar{\mathbf{t}}$ (current traction per current area); $\mathbf{f}_{0}\to\mathbf{b}$ (body force per current volume). Hence $\delta\mathcal{W}_{\mathrm{ext}}\to\int_{\Omega}\mathbf{b}\cdot\delta\mathbf{u}\,\mathrm{d}v+\int_{\partial\Omega_{t}}\bar{\mathbf{t}}\cdot\delta\mathbf{u}\,\mathrm{d}a$, the standard small-strain external virtual work.

% =====================================================================
\subsection{F. Inertia virtual work}
% =====================================================================

\textbf{F.1 Definition (d'Alembert).} The inertia body force per unit reference volume is $-\rho_{0}\ddot{\boldsymbol{\varphi}}$. Its virtual work defines

\begin{equation}
\delta\mathcal{W}_{\mathrm{iner}}
:= \int_{\mathcal{B}_{0}}
     \rho_{0}\,\ddot{\boldsymbol{\varphi}}\cdot\delta\boldsymbol{\varphi}\,
     \mathrm{d}V_{0}
 = \int_{\mathcal{B}_{0}}
     \rho_{0}\,\ddot{\mathbf{u}}\cdot\delta\mathbf{u}\,
     \mathrm{d}V_{0},
\label{rn:Winer}
\end{equation}

with $\ddot{\boldsymbol{\varphi}}=\partial^{2}\boldsymbol{\varphi}/\partial t^{2}=\ddot{\mathbf{u}}$ (since $\mathbf{X}$ is time-independent in the second material derivative as well).

\textbf{F.2 Sign convention.} The d'Alembert inertia force $-\rho_{0}\ddot{\boldsymbol{\varphi}}$ performs virtual work $-\delta\mathcal{W}_{\mathrm{iner}}$. We define $\delta\mathcal{W}_{\mathrm{iner}}$ with the positive sign above so that the instantaneous equilibrium reads $\delta\Pi_{\mathrm{int}}-\delta\mathcal{W}_{\mathrm{ext}}+\delta\mathcal{W}_{\mathrm{iner}}=0$ (derived in the next section; recorded here only as a forward reference, not derived).

\textbf{F.3 Link to $\delta\mathcal{T}$ (stated, not derived).} Because $\delta\mathcal{T}=\int_{\mathcal{B}_{0}}\rho_{0}\dot{\boldsymbol{\varphi}}\cdot\delta\dot{\boldsymbol{\varphi}}\,\mathrm{d}V_{0}$, integration by parts in time using $\delta\boldsymbol{\varphi}(t_{1})=\delta\boldsymbol{\varphi}(t_{2})=\mathbf{0}$ gives $\delta\int\mathcal{T}\,\mathrm{d}t=-\int\delta\mathcal{W}_{\mathrm{iner}}\,\mathrm{d}t$. This step is \textbf{deferred} to the next section (no integration by parts is performed here).

\textbf{F.4 Dimensional check.} $[\rho_{0}\ddot{\boldsymbol{\varphi}}]=\mathrm{kg\,m^{-3}}\cdot\mathrm{m\,s^{-2}}=\mathrm{N\,m^{-3}}$; times $[\delta\boldsymbol{\varphi}]=\mathrm{m}$ times $[\mathrm{d}V_{0}]=\mathrm{m^{3}}$ gives $\mathrm{N\,m}=\mathrm{J}$. \checkmark

% =====================================================================
\subsection{G. Total variational functional (Lagrangian + action)}
% =====================================================================

\textbf{G.1 Lagrangian functional.} With $\psi$ acting as the potential,

\begin{equation}
\mathcal{L}[\boldsymbol{\varphi},\dot{\boldsymbol{\varphi}}]
= \mathcal{T}(\dot{\boldsymbol{\varphi}})
- \Pi_{\mathrm{int}}[\boldsymbol{\varphi}].
\label{rn:Lagrangian}
\end{equation}

\textbf{G.2 Action functional.} Over $[t_{1},t_{2}]$,

\begin{equation}
\mathcal{A}[\boldsymbol{\varphi}]
= \int_{t_{1}}^{t_{2}}
  \mathcal{L}[\boldsymbol{\varphi},\dot{\boldsymbol{\varphi}}]\,\mathrm{d}t
= \int_{t_{1}}^{t_{2}}
  \Bigl[\,\mathcal{T}(\dot{\boldsymbol{\varphi}})
        -\Pi_{\mathrm{int}}[\boldsymbol{\varphi}]\Bigr]\,\mathrm{d}t.
\label{rn:action}
\end{equation}

Explicitly expanding both functionals,

\begin{equation}
\mathcal{A}[\boldsymbol{\varphi}]
= \int_{t_{1}}^{t_{2}}
  \left[
    \int_{\mathcal{B}_{0}}\tfrac12\rho_{0}\dot{\boldsymbol{\varphi}}\cdot\dot{\boldsymbol{\varphi}}\,\mathrm{d}V_{0}
  - \int_{\mathcal{B}_{0}}\psi(\mathbf{E},\bm{\mathcal{G}})\,\mathrm{d}V_{0}
  \right]\mathrm{d}t.
\label{rn:action_expand}
\end{equation}

\textbf{G.3 Unit check.} $[\mathcal{L}]=[\mathrm{J}]$, so $[\mathcal{A}]=[\mathrm{J\,s}]$ (action). \checkmark

% =====================================================================
\subsection{H. Final unreduced variational statement}
% =====================================================================

\textbf{H.1 Statement.} Hamilton's principle for the present finite-strain Mindlin Form-II strain-gradient solid is

\begin{equation}
\boxed{
\delta\mathcal{A}
+ \int_{t_{1}}^{t_{2}}\delta\mathcal{W}_{\mathrm{ext}}\,\mathrm{d}t = 0
\quad\Longleftrightarrow\quad
\delta\int_{t_{1}}^{t_{2}}
\Bigl[\,\mathcal{T}-\Pi_{\mathrm{int}}\Bigr]\,\mathrm{d}t
+ \int_{t_{1}}^{t_{2}}\delta\mathcal{W}_{\mathrm{ext}}\,\mathrm{d}t = 0,
}
\label{rn:Hamiltoneq}
\end{equation}

for every admissible virtual displacement $\delta\boldsymbol{\varphi}$ satisfying the homogeneous essential (kinematic) boundary conditions and the time-endpoint conditions $\delta\boldsymbol{\varphi}(t_{1})=\delta\boldsymbol{\varphi}(t_{2})=\mathbf{0}$.

\textbf{H.2 Admissibility / boundary conditions (forward reference only).} We require only that $\delta\boldsymbol{\varphi}$ vanish on the essential (Dirichlet) part $\partial\mathcal{B}_{0u}$ and at $t_{1},t_{2}$. The precise partition of the boundary into displacement-essential vs.\ traction/couple-stress-natural conditions for a $C^{1}$ (second-gradient) theory is obtained by integration by parts in the next section and is \textbf{not} derived here. Stating the admissibility set suffices for the unreduced statement.

\textbf{H.3 Static (quasi-static) reduction.} Setting $\mathcal{T}\equiv 0$ gives

\begin{equation}
\delta\Pi_{\mathrm{int}}-\delta\mathcal{W}_{\mathrm{ext}}=0,
\label{rn:static}
\end{equation}

the principle of virtual work. This is the form used by the quasi-static 2022 benchmark. The dynamic part is retained so that the formulation also covers inertial (buckling / vibration) problems, consistent with the programme's Q1 novelty target.

% =====================================================================
\subsection{I. Work-conjugacy verification}
% =====================================================================

\textbf{I.1 Internal virtual work as a first variation.} Because $\mathcal{B}_{0}$ is fixed, $\delta\Pi_{\mathrm{int}}=\int_{\mathcal{B}_{0}}\delta\psi\,\mathrm{d}V_{0}$. The density $\psi$ has declared arguments $(\mathbf{E},\bm{\mathcal{G}})$ only, so the multivariable chain rule gives

\begin{equation}
\delta\psi
= \frac{\partial\psi}{\partial\mathbf{E}}:\delta\mathbf{E}
+ \frac{\partial\psi}{\partial\bm{\mathcal{G}}}:\delta\bm{\mathcal{G}}.
\label{rn:dpsi}
\end{equation}

Define the work-conjugate stresses (explicit closed forms are derived in the Stress-Tensors section; here only their definitions and conjugacy are established, respecting scope):

\begin{itemize}
\item $\mathbf{S}:=\partial\psi/\partial\mathbf{E}$ -- second Piola--Kirchhoff stress (2nd-order material).

\item $\boldsymbol{\Sigma}:=\partial\psi/\partial\bm{\mathcal{G}}$ -- double stress (3rd-order material, with $\Sigma_{ABC}=\Sigma_{BAC}$ to match the symmetry of $\bm{\mathcal{G}}$).
\end{itemize}

Then

\begin{equation}
\delta\Pi_{\mathrm{int}}
= \int_{\mathcal{B}_{0}}
  \Bigl[\,\mathbf{S}:\delta\mathbf{E}
        +\boldsymbol{\Sigma}:\delta\bm{\mathcal{G}}\Bigr]\,
  \mathrm{d}V_{0}.
\label{rn:Wint}
\end{equation}

\textbf{I.2 Conjugacy contractions (every term yields a scalar energy density).}

\begin{itemize}
\item $\mathbf{S}:\delta\mathbf{E}=S_{IJ}\,\delta E_{IJ}$; $\mathbf{S},\mathbf{E}$ both 2nd-order material $\Rightarrow$ scalar. $\int[\cdot]\,\mathrm{d}V_{0}\to$ energy. \checkmark

\item $\boldsymbol{\Sigma}:\delta\bm{\mathcal{G}}=\Sigma_{ABC}\,\delta\mathcal{G}_{ABC}$; $\boldsymbol{\Sigma},\bm{\mathcal{G}}$ both 3rd-order with matching $A\leftrightarrow B$ symmetry $\Rightarrow$ scalar. \checkmark
\end{itemize}

Each stress is therefore the unique work-conjugate of its kinematic variable, and the internal virtual work is a sum of frame-indifferent scalar contributions. \checkmark

\textbf{I.3 Variation chain (symbolic, every link).}

\begin{align}
\delta\mathbf{F} &= \nabla_{0}\delta\mathbf{u}
  \quad(\text{exact; } \mathbf{F}=\mathbf{I}+\nabla_{0}\mathbf{u},\ \delta\perp\nabla_{0}),
  \label{rn:chain_F}\\
\delta\mathbf{E}
  &= \tfrac12\bigl(\delta\mathbf{F}^{T}\mathbf{F}+\mathbf{F}^{T}\delta\mathbf{F}\bigr)
  \quad(\text{exact; from } \mathbf{E}=\tfrac12(\mathbf{F}^{T}\mathbf{F}-\mathbf{I})),
  \label{rn:chain_E}\\
\delta\bm{\mathcal{G}} &= \nabla_{0}\delta\mathbf{E}.
  \label{rn:chain_G}
\end{align}

The chain $\delta\mathbf{u}\to\delta\mathbf{F}\to\delta\mathbf{E}\to\delta\bm{\mathcal{G}}$ is exact and well defined.

\textbf{I.4 External and inertia conjugacy.} In \eqref{rn:Wext} the body force $\mathbf{f}_{0}$ is work-conjugate to $\delta\mathbf{u}$ (volume term) and the nominal traction $\bar{\mathbf{T}}_{0}$ to $\delta\mathbf{u}$ (surface term); in \eqref{rn:Winer} the inertia density $\rho_{0}\ddot{\boldsymbol{\varphi}}$ is work-conjugate to $\delta\boldsymbol{\varphi}$. All virtual-work terms are thus scalar-energy contributions. \checkmark

% =====================================================================
\subsection{J. Dimensional consistency table}
% =====================================================================

\begin{table}[h]
\centering
\caption{Dimensional consistency of the Section~3A functionals.}
\label{tab:dim3a}
\begin{tabular}{llll}
\toprule
Quantity & Definition & Units of integrand & Total units \\
\midrule
$\mathcal{T}$ & $\int\tfrac12\rho_{0}\dot{\boldsymbol{\varphi}}^{2}\,\mathrm{d}V_{0}$ & $(\mathrm{kg\,m^{-3}})(\mathrm{m^{2}\,s^{-2}})(\mathrm{m^{3}})$ & $\mathrm{J}$ \\
$\Pi_{\mathrm{int}}$ & $\int\psi\,\mathrm{d}V_{0}$ & $(\mathrm{J\,m^{-3}})(\mathrm{m^{3}})$ & $\mathrm{J}$ \\
$\delta\mathcal{W}_{\mathrm{ext}}$ & body + surface & $(\mathrm{N\,m^{-3}})(\mathrm{m})(\mathrm{m^{3}})$, $(\mathrm{N\,m^{-2}})(\mathrm{m})(\mathrm{m^{2}})$ & $\mathrm{J}$ \\
$\delta\mathcal{W}_{\mathrm{iner}}$ & $\int\rho_{0}\ddot{\boldsymbol{\varphi}}\cdot\delta\boldsymbol{\varphi}\,\mathrm{d}V_{0}$ & $(\mathrm{N\,m^{-3}})(\mathrm{m})(\mathrm{m^{3}})$ & $\mathrm{J}$ \\
$\mathcal{L}$ & $\mathcal{T}-\Pi_{\mathrm{int}}$ & -- & $\mathrm{J}$ \\
$\mathcal{A}$ & $\int\mathcal{L}\,\mathrm{d}t$ & -- & $\mathrm{J\,s}$ \\
\bottomrule
\end{tabular}
\end{table}

All terms share the energy unit $[\mathrm{J}]$, so the variational statement \eqref{rn:Hamiltoneq} is dimensionally homogeneous. \checkmark

% =====================================================================
\subsection{K. Objectivity verification}
% =====================================================================

Under a superposed rigid-body motion $\bar{\mathbf{x}}=\mathbf{Q}(t)\mathbf{x}+\mathbf{c}(t)$, $\mathbf{Q}\in SO(3)$:

\begin{itemize}
\item \textbf{Internal energy.} $\bar{\mathbf{E}}=\mathbf{E}$, $\bar{\bm{\mathcal{G}}}=\bm{\mathcal{G}}$ (objective material tensors, verified in Section~\ref{sec:recovery}, sec:obj). Hence $\psi(\bar{\mathbf{E}},\bar{\bm{\mathcal{G}}})=\psi(\mathbf{E},\bm{\mathcal{G}})$ and $\Pi_{\mathrm{int}}[\bar{\boldsymbol{\varphi}}]=\Pi_{\mathrm{int}}[\boldsymbol{\varphi}]$. Objective. \checkmark

\item \textbf{Kinetic energy.} As discussed in C.5, the standard inertial-frame kinetic energy is used; the action is Galilean-invariant up to a total time derivative, so the Euler--Lagrange equations are frame-indifferent \cite{Goldstein2002,MarsdenHughes1983}. \checkmark

\item \textbf{External virtual work.} Prescribed loads are material vectors; under a rigid superposed motion they transform consistently and the virtual work is frame-indifferent in an inertial frame. \checkmark

\item \textbf{Conclusion.} The reference (Total-Lagrangian) description guarantees that the emerging field equations are objective; this is the standard advantage of the TL action \cite{BonetWood2008,MarsdenHughes1983}. \checkmark
\end{itemize}

% =====================================================================
\subsection{L. Compatibility with Sections~\ref{sec:kinematics} and~\ref{sec:constitutive}}
% =====================================================================

\begin{table}[h]
\centering
\caption{Audit of objects reused from frozen sections.}
\label{tab:compat3a}
\begin{tabular}{lll}
\toprule
Object & Definition used & Source \\
\midrule
$\boldsymbol{\varphi},\mathbf{u},\mathbf{F},J$ & eq:n\_motion, eq:n\_disp, eq:n\_F, eq:n\_J & Section~\ref{sec:kinematics} \\
$\mathbf{E},\bm{\mathcal{G}}$ & eq:n\_E, eq:n\_Gdef & Section~\ref{sec:kinematics} (sec:Gdef) \\
$\psi(\mathbf{E},\bm{\mathcal{G}})$ & eq:psi & Section~\ref{sec:constitutive} (eq:sec2\_final\_constitutive\_framework) \\
Parameters $\lambda,\mu,\mathbb{D}_{ABCPQR}$ (equivalently $a_{1},\ldots,a_{5}$) & -- & Section~\ref{sec:constitutive} (eq:sec2\_final\_constitutive\_framework, tab:sec2\_constitutive\_constants) \\
Index convention (material $I,J,K$; spatial $i,j,k$) & -- & Section~\ref{sec:kinematics} (sec:ledger) \\
\bottomrule
\end{tabular}
\end{table}

No notation conflict; all Section~\ref{sec:kinematics}/\ref{sec:constitutive} symbols are reused exactly; $\psi$ is cited as eq:psi without alteration. \checkmark

% =====================================================================
\subsection{M. Small-strain reduction to the validated 2022 benchmark}
% =====================================================================

\textbf{M.1 The 2022 small-strain benchmark functionals.} Reconstructed from the small-strain limits established in Sections~\ref{sec:kinematics}--\ref{sec:constitutive} (eq:n\_limit1--eq:n\_limit2 and sec:recovery; eq:psi of Section~\ref{sec:constitutive}), the validated 2022 benchmark energy density is

\begin{align}
\psi_{2022}(\boldsymbol{\varepsilon},\nabla\boldsymbol{\varepsilon})
&= \tfrac12\lambda(\varepsilon_{AA})^{2}
   +\mu\,\varepsilon_{AB}\varepsilon_{AB}
   + a_{1}\,\varepsilon_{AB,B}\varepsilon_{AC,C}
   + a_{2}\,\varepsilon_{AA,C}\varepsilon_{CB,B}
   + a_{3}\,\varepsilon_{AA,C}\varepsilon_{BB,C}
\nonumber\\
&\quad
   + a_{4}\,\varepsilon_{AB,C}\varepsilon_{AB,C}
   + a_{5}\,\varepsilon_{AB,C}\varepsilon_{CB,A},
\label{rn:psi2022}
\end{align}

with the \emph{same} $\lambda,\mu,a_1,\ldots,a_5$ as eq:psi.

\textbf{M.2 Limit chain (every step).} As $\|\nabla_{0}\mathbf{u}\|\to0$:

\begin{enumerate}
\item $J=\det\mathbf{F}\to 1$, so $\mathrm{d}V_{0}\to\mathrm{d}v$ (to leading order).

\item $\mathbf{E}\to\boldsymbol{\varepsilon}$ (eq:n\_limit2); $\bm{\mathcal{G}}=\nabla_{0}\mathbf{E}\to\nabla_{0}\boldsymbol{\varepsilon}$ (sec:recovery). The rotation-gradient kinematic variable $\nabla_{0}\boldsymbol{\theta}$ reduces to the infinitesimal rotation gradient (sec:form2, linear limit) but does not enter the energy.

\item $\rho_{0}$ is constant (homogeneous): $\rho_{0}\to\rho$ at leading order, with $\rho=\rho_{0}/J\to\rho_{0}$.

\item $\mathbf{f}_{0}\to\mathbf{b}$ (body force per current volume); $\bar{\mathbf{T}}_{0}\to\bar{\mathbf{t}}$ (sec:E.4).
\end{enumerate}

\textbf{M.3 Functional reduction.}

\begin{align}
\Pi_{\mathrm{int}} &\to \int_{\Omega}\psi_{2022}(\boldsymbol{\varepsilon},\nabla\boldsymbol{\varepsilon})\,\mathrm{d}v,
\label{rn:rec_Pi}\\
\mathcal{T} &\to \int_{\Omega}\tfrac12\rho\,\dot{\mathbf{u}}\cdot\dot{\mathbf{u}}\,\mathrm{d}v,
\label{rn:rec_T}\\
\delta\mathcal{W}_{\mathrm{ext}} &\to \int_{\Omega}\mathbf{b}\cdot\delta\mathbf{u}\,\mathrm{d}v
                                  + \int_{\partial\Omega_{t}}\bar{\mathbf{t}}\cdot\delta\mathbf{u}\,\mathrm{d}a,
\label{rn:rec_Wext}\\
\mathcal{A} &\to \int_{t_{1}}^{t_{2}}\Bigl[\mathcal{T}_{2022}-\Pi_{\mathrm{int},2022}\Bigr]\,\mathrm{d}t.
\label{rn:rec_A}
\end{align}

Because eq:psi was \emph{constructed} (Section~\ref{sec:constitutive}) to reduce exactly to $\psi_{2022}$ in the linear limit, the reduction of $\Pi_{\mathrm{int}}$ is exact (discarded terms strictly $O(\|\nabla_{0}\mathbf{u}\|^{2})$). The static part of \eqref{rn:Hamiltoneq} therefore reduces exactly to the 2022 benchmark principle of virtual work.

\textbf{M.4 Mandatory-recovery-rule verdict.} \checkmark Satisfied at the variational/functional level: every functional of the present finite-strain formulation collapses exactly to the corresponding functional of the validated 2022 small-strain Mindlin Form-II benchmark. The benchmark is unchanged.

% =====================================================================
\subsection{N. Literature discussion}
% =====================================================================

\textbf{New references introduced in this section} (full entries in \texttt{references\_calculation.bib}):

\begin{itemize}
\item \textbf{Hamilton (1834)} \cite{Hamilton1834} -- original statement of Hamilton's principle (pre-DOI historical source).

\item \textbf{Goldstein, Poole \& Safko (2002)} \cite{Goldstein2002} -- textbook derivation of Hamilton's principle, the Lagrange--d'Alembert augmentation, kinetic energy, and Galilean invariance of the action.

\item \textbf{Gurtin (1981)} \cite{Gurtin1981} -- continuum energy/action principles and virtual work for finite-deformation solids.

\item \textbf{Maugin (1980)} \cite{Maugin1980} -- method of virtual power extended to higher-grade (gradient) continua; supports the single variational statement that later accommodates $\delta\bm{\mathcal{G}}$ conjugates. \emph{DOI flagged for verification before final submission.}
\end{itemize}

\textbf{Reused (do not duplicate in the master .bib):} Germain1973 (virtual power), MindlinEshel1968, Toupin1962 (gradient/virtual work), MarsdenHughes1983 (nonlinear action, objectivity), BonetWood2008, Bathe2006, Holzapfel2000 (TL kinetic energy, Nanson), Beheshti2016 (finite-strain gradient reference formulation), Mindlin1964 (Form II). All DOIs for reused entries were verified in prior sections.

% =====================================================================
\subsection{O. Reviewer comments (Q\&A)}
% =====================================================================

\begin{enumerate}
\item \textbf{R: Why include kinetic energy / inertia in a benchmark that is (quasi-)static?}
A: The variational formulation must be general (dynamic buckling/vibration is part of the Q1 novelty target). The static 2022 benchmark is recovered exactly by setting $\mathcal{T}\equiv 0$ in \eqref{rn:Hamiltoneq}, which leaves the static potential $\Pi_{\mathrm{int}}$ and $\delta\mathcal{W}_{\mathrm{ext}}$ untouched. The mandatory recovery rule is therefore preserved (Section~M).

\item \textbf{R: Is the sign convention for $\delta\mathcal{W}_{\mathrm{iner}}$ consistent?}
A: Yes. We define $\delta\mathcal{W}_{\mathrm{iner}}:=\int\rho_{0}\ddot{\boldsymbol{\varphi}}\cdot\delta\boldsymbol{\varphi}\,\mathrm{d}V_{0}$; the d'Alembert inertia force $-\rho_{0}\ddot{\boldsymbol{\varphi}}$ thus contributes $-\delta\mathcal{W}_{\mathrm{iner}}$. The equivalence $\delta\int\mathcal{T}\,\mathrm{d}t=-\int\delta\mathcal{W}_{\mathrm{iner}}\,\mathrm{d}t$ is established by time integration by parts in the next section; we only state it here (no IBP performed).

\item \textbf{R: Why are all integrals over the reference volume, including kinetic energy?}
A: Total-Lagrangian consistency (assumption A4). Keeping every quantity material guarantees exact benchmark recovery and an objective formulation; the reference kinetic energy $\int\tfrac12\rho_{0}\dot{\boldsymbol{\varphi}}^{2}\,\mathrm{d}V_{0}$ is the standard TL form \cite{BonetWood2008,Bathe2006,Holzapfel2000}.

\item \textbf{R: Is the nominal-traction boundary expression correct?}
A: We define $\bar{\mathbf{T}}_{0}=\mathbf{P}\mathbf{N}$ (with $P_{iJ}$ convention matching $\mathbf{F}_{iJ}$ of Section~\ref{sec:kinematics}) as the natural condition that \emph{will} emerge from integration by parts; we state only the exact Nanson vector-area relation $\mathrm{d}\mathbf{a}=J\mathbf{F}^{-T}\mathbf{N}\,\mathrm{d}A_{0}$ and avoid any questionable scalar-vector pull-back formula. The explicit derivation is deferred.

\item \textbf{R: Stresses appear in the work-conjugacy subsection -- did you derive them?}
A: No. $\mathbf{S},\boldsymbol{\Sigma}$ are introduced only as the definitions $\partial\psi/\partial\mathbf{E},\partial\psi/\partial\bm{\mathcal{G}}$ to verify the \emph{structure} of work conjugacy. Their explicit closed forms are the content of the next (Stress-Tensors) section, per the strict scope.

\item \textbf{R: Is the action principle frame-indifferent?}
A: Yes (Section~K). The constitutive part is objective because all arguments are material tensors; the kinetic-energy term is the standard inertial-frame form and is Galilean-invariant up to a total derivative. Non-inertial frames (requiring extra inertial body forces) are outside scope.

\item \textbf{R: Have you split the boundary conditions?}
A: No. We require only that admissible variations satisfy the homogeneous essential conditions and vanish at $t_{1},t_{2}$. The $C^{1}$ (second-gradient) theory's partition into displacement/slope essential vs.\ traction/double-traction natural conditions is resolved by integration by parts in the next section; stating the admissibility set here is sufficient for the unreduced statement.
\end{enumerate}

\noindent\textbf{Closing statement for Section~3A (part 1).} The variational (Hamilton) principle for the geometrically nonlinear, displacement-based, $C^{1}$ BFS finite-strain Mindlin Form-II strain-gradient formulation is fully specified as the stationary-action statement \eqref{rn:Hamiltoneq}, built from the kinetic-energy functional $\mathcal{T}$ (C), the total Helmholtz free-energy functional $\Pi_{\mathrm{int}}$ (D, via eq:psi of Section~\ref{sec:constitutive}), the external virtual work $\delta\mathcal{W}_{\mathrm{ext}}$ (E), and the inertia virtual work $\delta\mathcal{W}_{\mathrm{iner}}$ (F). All terms are dimensionally consistent (J), objective (K), work-conjugate (I), and reduce exactly to the validated 2022 benchmark in the small-strain limit (M), satisfying the mandatory recovery rule. Integration by parts and the strong form are developed below.

% =====================================================================
\subsection{Integration by parts -- complete derivation}
\label{subsec:ibp}
% =====================================================================

\paragraph{P.1 Starting point.}

The unreduced statement is \eqref{rn:Hamiltoneq}:
$\delta\int_{t_1}^{t_2}[\mathcal{T}-\Pi_{\mathrm{int}}]\,\mathrm{d}t+\int_{t_1}^{t_2}\delta\mathcal{W}_{\mathrm{ext}}\,\mathrm{d}t=0$.

Take the first variation of every term. The domain $\mathcal{B}_0$ and the time interval $[t_1,t_2]$ are fixed, so variations commute with the integrals and with $\nabla_0$ and $\partial/\partial t$.

\paragraph{P.2 Time integration by parts (kinetic energy).}

$\mathcal{T}(t)=\int_{\mathcal{B}_0}\tfrac12\rho_0\dot{\varphi}_i\dot{\varphi}_i\,\mathrm{d}V_0$ (eq:T). Its variation at fixed $t$ is
$\delta\mathcal{T}=\int_{\mathcal{B}_0}\rho_0\dot{\varphi}_i\,\delta\dot{\varphi}_i\,\mathrm{d}V_0$,
because $\delta(\dot{\varphi}_i\dot{\varphi}_i)=2\dot{\varphi}_i\delta\dot{\varphi}_i$ and $\rho_0$ is fixed. The variation commutes with the material time derivative:
$\delta\dot{\varphi}_i=\delta(\partial\varphi_i/\partial t|_{\mathbf{X}})=\partial(\delta\varphi_i)/\partial t|_{\mathbf{X}}=\delta\dot{\varphi}_i$ (i.e.\ $\delta\dot{\boldsymbol{\varphi}}=\delta\dot{\mathbf{u}}$). Hence

\begin{align}
\delta\int_{t_1}^{t_2}\mathcal{T}\,\mathrm{d}t
&= \int_{t_1}^{t_2}\int_{\mathcal{B}_0}\rho_0\dot{\varphi}_i\,\delta\dot{\varphi}_i\,\mathrm{d}V_0\,\mathrm{d}t \notag\\
&= \int_{\mathcal{B}_0}\rho_0\left[\int_{t_1}^{t_2}\dot{\varphi}_i\,\frac{\partial(\delta\varphi_i)}{\partial t}\,\mathrm{d}t\right]\mathrm{d}V_0 \notag\\
&= \int_{\mathcal{B}_0}\rho_0\Bigl[\dot{\varphi}_i\delta\varphi_i\Bigr]_{t_1}^{t_2}\,\mathrm{d}V_0
 - \int_{t_1}^{t_2}\int_{\mathcal{B}_0}\rho_0\ddot{\varphi}_i\,\delta\varphi_i\,\mathrm{d}V_0\,\mathrm{d}t.
\label{rn2:timeibp}
\end{align}

The endpoint term vanishes by the admissible-condition $\delta\boldsymbol{\varphi}(t_1)=\delta\boldsymbol{\varphi}(t_2)=\mathbf{0}$. Since $\ddot{\varphi}_i=\ddot{u}_i$ and $\delta\varphi_i=\delta u_i$,

\begin{equation}
\delta\int_{t_1}^{t_2}\mathcal{T}\,\mathrm{d}t
= -\int_{t_1}^{t_2}\underbrace{\left(\int_{\mathcal{B}_0}\rho_0\ddot{u}_i\,\delta u_i\,\mathrm{d}V_0\right)}_{\displaystyle\delta\mathcal{W}_{\mathrm{iner}}}\mathrm{d}t.
\label{rn2:timeibp2}
\end{equation}

This is exactly the identity deferred from subsec:inertia. \checkmark

\paragraph{P.3 Internal virtual work (restated).}

From the directional-derivative form (rn:Wint, with stresses defined as $\mathbf{S}=\partial\psi/\partial\mathbf{E}$, $\boldsymbol{\Sigma}=\partial\psi/\partial\bm{\mathcal{G}}$),

\begin{equation}
\delta\Pi_{\mathrm{int}}
= \int_{\mathcal{B}_0}
  \bigl[\,S_{IJ}\,\delta E_{IJ}
        +\Sigma_{IJK}\,\delta\mathcal{G}_{IJK}\bigr]\,
  \mathrm{d}V_0.
\label{rn2:intvw}
\end{equation}

\paragraph{P.4 Kinematic variations in terms of $\delta\mathbf{u}$.}

Using $\delta F_{iI}=\delta u_{i,I}$ (since $F_{iI}=\delta_{iI}+u_{i,I}$):

\begin{align}
\delta E_{IJ}
 &= \tfrac12\bigl(\delta F_{iI}F_{iJ}+F_{iI}\delta F_{iJ}\bigr)
  = \tfrac12\bigl(F_{iI}\delta u_{i,J}+F_{iJ}\delta u_{i,I}\bigr),
  \label{rn2:dE}\\
\delta\mathcal{G}_{IJK}
 &= \delta(\partial E_{IJ}/\partial X_K)
  = \partial(\delta E_{IJ})/\partial X_K
  = \delta E_{IJ,K}.
  \label{rn2:dG}
\end{align}

Because $\mathbf{R}=\mathrm{polar}(\mathbf{F})$ is a rotation, the rotation gradient variation $\delta\mathbf{K}=\nabla_0\delta\boldsymbol{\theta}$ is slaved to $\delta\mathbf{F}$ through the polar-decomposition sensitivity; the rotation is kinematic only and does not enter the energy, so no independent $\delta\mathbf{K}$ variation appears in $\delta\Pi_{\mathrm{int}}$.

\paragraph{P.5 First Piola--Kirchhoff stress from the $\mathbf{S}:\delta\mathbf{E}$ branch.}

Substitute \eqref{rn2:dE}:
$S_{IJ}\delta E_{IJ}
= \tfrac12 S_{IJ}F_{iI}\delta u_{i,J}+\tfrac12 S_{IJ}F_{iJ}\delta u_{i,I}$.

In the second term swap dummy indices $I\leftrightarrow J$ and use $S_{JI}=S_{IJ}$:
$= \tfrac12 S_{IJ}F_{iI}\delta u_{i,J}+\tfrac12 S_{IJ}F_{iI}\delta u_{i,J}
= S_{IJ}F_{iI}\delta u_{i,J}$.

Define the first Piola--Kirchhoff stress $P_{iJ}=S_{IJ}F_{iI}$ (summation over $I$). Then $S_{IJ}\delta E_{IJ}=P_{iJ}\delta u_{i,J}=P_{iJ}\delta F_{iJ}$. Integrate by parts over $X_J$:

\begin{align}
\int_{\mathcal{B}_0}P_{iJ}\delta u_{i,J}\,\mathrm{d}V_0
&= \int_{\mathcal{B}_0}(P_{iJ}\delta u_i)_{,J}\,\mathrm{d}V_0
 - \int_{\mathcal{B}_0}P_{iJ,J}\,\delta u_i\,\mathrm{d}V_0 \notag\\
&= \int_{\partial\mathcal{B}_0}P_{iJ}N_J\,\delta u_i\,\mathrm{d}A_0
 - \int_{\mathcal{B}_0}P_{iJ,J}\,\delta u_i\,\mathrm{d}V_0.
\label{rn2:ibpS}
\end{align}

Domain contribution: $-P_{iJ,J}\delta u_i$. Boundary contribution: $P_{iJ}N_J\delta u_i$. \checkmark

\paragraph{P.6 Integration by parts on the $\boldsymbol{\Sigma}:\delta\bm{\mathcal{G}}$ branch.}

$\Sigma_{IJK}\delta\mathcal{G}_{IJK}=\Sigma_{IJK}\delta E_{IJ,K}$. \textbf{First IBP} over $X_K$:

\begin{align}
\int_{\mathcal{B}_0}\Sigma_{IJK}\delta E_{IJ,K}\,\mathrm{d}V_0
&= \int_{\partial\mathcal{B}_0}\Sigma_{IJK}N_K\,\delta E_{IJ}\,\mathrm{d}A_0
 - \int_{\mathcal{B}_0}\Sigma_{IJK,K}\,\delta E_{IJ}\,\mathrm{d}V_0.
\label{rn2:ibpSig1}
\end{align}

The remaining domain term still contains $\delta E_{IJ}$ (a first derivative of $\delta\mathbf{u}$). Substitute \eqref{rn2:dE}:
$-\int\Sigma_{IJK,K}\tfrac12(F_{iI}\delta u_{i,J}+F_{iJ}\delta u_{i,I})\,\mathrm{d}V_0$.

In the second half swap $I\leftrightarrow J$ (using $\Sigma_{JIK}=\Sigma_{IJK}$ and $F_{iJ}\to F_{iI}$): it equals the first half, so the sum is
$-\int_{\mathcal{B}_0}\Sigma_{IJK,K}F_{iI}\,\delta u_{i,J}\,\mathrm{d}V_0$.

\textbf{Second IBP} over $X_J$:

\begin{align}
-\int_{\mathcal{B}_0}\Sigma_{IJK,K}F_{iI}\,\delta u_{i,J}\,\mathrm{d}V_0
&= -\int_{\partial\mathcal{B}_0}(\Sigma_{IJK,K}F_{iI})N_J\,\delta u_i\,\mathrm{d}A_0
 + \int_{\mathcal{B}_0}(\Sigma_{IJK,K}F_{iI})_{,J}\,\delta u_i\,\mathrm{d}V_0.
\label{rn2:ibpSig2}
\end{align}

Define the hyperstress $\mathcal{H}_{iJK}=\Sigma_{IJK}F_{iI}$ (rank three; free indices $i$ spatial, $J,K$ material). The domain contribution of the $\boldsymbol{\Sigma}$ branch is $(\Sigma_{IJK,K}F_{iI})_{,J}\delta u_i$. By the product rule,

\begin{equation}
(\Sigma_{IJK,K}F_{iI})_{,J}=\Sigma_{IJK,KJ}F_{iI}+\Sigma_{IJK,K}F_{iI,J},
\label{rn2:hypdiv}
\end{equation}

which equals the double divergence $\mathcal{H}_{iJK,KJ}$ only after adding back the missing $-\Sigma_{IJK,J}F_{iI,K}-\Sigma_{IJK}F_{iI,JK}$. The strong form therefore uses $(\Sigma_{IJK,K}F_{iI})_{,J}$.

\textbf{Important note on boundary-term identifications.} At finite strain, the product rule gives $\mathcal{H}_{iJK,K}=\Sigma_{IJK,K}F_{iI}+\Sigma_{IJK}F_{iI,K}$, so $\mathcal{H}_{iJK,K}\neq\Sigma_{IJK,K}F_{iI}$ in general. The boundary terms are therefore expressed in terms of $\Sigma_{IJK,K}F_{iI}$ (not $\mathcal{H}_{iJK,K}$) and the higher-order surface term is correctly rewritten via $\delta E_{IJ}=\tfrac12(F_{lI}\delta u_{l,J}+F_{lJ}\delta u_{l,I})$ and the symmetry $\Sigma_{IJK}=\Sigma_{JIK}$ as:

\begin{equation}
\Sigma_{IJK}N_K\,\delta E_{IJ}
= \Sigma_{IJK}F_{lI}\,N_K\,\delta u_{l,J}
= \mathcal{H}_{lJK}\,N_K\,\delta u_{l,J},
\label{rn2:higher_order_bnd}
\end{equation}

which is a scalar (all indices $l,J,K$ contracted). This expression is used to extract the double traction $\mathcal{D}_l=\mathcal{H}_{lJK}N_KN_J$ conjugate to the normal-slope variation $\delta u_{l,N}=N_J\delta u_{l,J}$ on the free boundary.

Collecting the $\boldsymbol{\Sigma}$ branch:

\begin{itemize}
\item Domain: $+(\Sigma_{IJK,K}F_{iI})_{,J}\,\delta u_i$ (the fourth-order operator).

\item Boundary piece (a): $+\Sigma_{IJK}N_K\,\delta E_{IJ} = \mathcal{H}_{lJK}N_K\,\delta u_{l,J}$ (the higher-order / double-traction term; the correct $\mathcal{H}$-based form follows from eq:rn2:higher\_order\_bnd).

\item Boundary piece (b): $-(\Sigma_{IJK,K}F_{iI})N_J\,\delta u_i$ (the effective higher-order traction correction; at finite strain this is \emph{not} equal to $-\mathcal{H}_{iJK,K}N_J\,\delta u_i$ due to the product-rule term $\Sigma_{IJK}F_{iI,K}$).
\end{itemize}

\checkmark

\paragraph{P.7 Assembly of the domain (strong) form.}

Summing the volume contributions:
$-P_{iJ,J}\delta u_i$ (P.5) $+(\Sigma_{IJK,K}F_{iI})_{,J}\delta u_i$ (P.6) $-f_{0i}\delta u_i$ (external body force, from $-\delta\mathcal{W}_{\mathrm{ext}}$) $+\rho_0\ddot{u}_i\delta u_i$ (kinetic, from $+\delta\mathcal{W}_{\mathrm{iner}}$, P.2). Since $\delta\mathbf{u}$ is arbitrary in $\mathcal{B}_0$, the coefficient vanishes:

\begin{equation}
-P_{iJ,J}+(\Sigma_{IJK,K}F_{iI})_{,J}-f_{0i}+\rho_0\ddot{u}_i=0,
\qquad\text{in }\mathcal{B}_0.
\label{rn2:strong}
\end{equation}

Equivalently, bringing the source and inertia terms to the right-hand side:

\begin{equation}
P_{iJ,J}-(\Sigma_{IJK,K}F_{iI})_{,J}+f_{0i}=\rho_0\ddot{u}_i.
\label{rn2:strong_balance}
\end{equation}

Defining the effective first P-K stress $\mathbf{T}$ by $\mathrm{Div}_0\mathbf{T}=P_{iJ,J}-(\Sigma_{IJK,K}F_{iI})_{,J}$, this is the standard balance of linear momentum:

\begin{equation}
\mathrm{Div}_0\mathbf{T}+\mathbf{f}_0=\rho_0\ddot{\mathbf{u}}.
\label{rn2:strong_balance_short}
\end{equation}

\checkmark

\paragraph{P.8 Assembly of the boundary form.}

Summing the surface contributions:
$+P_{iJ}N_J\delta u_i$ (P.5) $-(\Sigma_{IJK,K}F_{iI})N_J\delta u_i$ (P.6b) $+\Sigma_{IJK}N_K\delta E_{IJ}$ (P.6a) $+\bar{T}_{0i}\delta u_i$ (external surface, moved to RHS as natural data). The coefficient of $\delta u_i$ (terms not involving derivatives of $\delta u_i$) is $(P_{iJ}-\Sigma_{IJK,K}F_{iI})N_J$, giving the effective single traction. The higher-order boundary term $\Sigma_{IJK}N_K\delta E_{IJ}=\mathcal{H}_{lJK}N_K\delta u_{l,J}$ is the double-traction virtual work (see eq:rn2:higher\_order\_bnd). \checkmark

% =====================================================================
\subsection{Strong governing equations -- interpretation and checks}
\label{subsec:strongnotes}
% =====================================================================

\textbf{Q.1 Physical meaning.} Equation~\eqref{rn2:strong} is a single vector balance (force balance) -- there is no independent moment equation because the rotation is slaved to the displacement (assumption A2, displacement-based Form-II). The term $-P_{iJ,J}$ is the classical TL divergence of the first P-K stress; $+(\Sigma_{IJK,K}F_{iI})_{,J}$ is the fourth-order gradient-elasticity term (double divergence of the Piola-transformed double stress, with the product-rule terms retained); $-f_{0i}$ the body force; $+\rho_0\ddot{u}_i$ the inertia.

\textbf{Q.2 Dimensional consistency.} Units table (material length $[L]$, force $[N]$, energy density $[\mathrm{Pa}]=[\mathrm{N\,L^{-2}}]$):

\begin{table}[h]
\centering
\caption{Dimensions of the strong-form terms (Section~3A).}
\label{rn2:dimtab}
\begin{tabular}{ll}
\toprule
Term & Unit \\
\midrule
$E_{IJ}$ & dimensionless; $\mathcal{G}_{IJK}=E_{IJ,K}$: $[L^{-1}]$ \\
$\Sigma_{IJK}=\partial\psi/\partial\mathcal{G}_{IJK}$ & $[\mathrm{Pa}]/[L^{-1}]=[\mathrm{N\,L^{-1}}]$ \\
$\mathcal{H}_{iJK}=\Sigma_{IJK}F_{iI}$ & $[\mathrm{N\,L^{-1}}]$ (third-order) \\
$(\Sigma_{IJK,K}F_{iI})_{,J}$ & $[\mathrm{N\,L^{-1}}]/[L^{2}]=[\mathrm{N\,L^{-3}}]$ (force density) \\
$P_{iJ}=S_{IJ}F_{iI}$, $P_{iJ,J}$ & $[\mathrm{Pa}]/[L]=[\mathrm{N\,L^{-3}}]$ \\
$f_{0i},\ \rho_0\ddot{u}_i$ & $[\mathrm{N\,L^{-3}}]$ \\
\bottomrule
\end{tabular}
\end{table}

All domain terms share $[\mathrm{N\,L^{-3}}]$; all boundary terms share $[\mathrm{N}]$ after multiplication by $\delta u_i$ ($[L]$) and $\mathrm{d}A_0$ ($[L^2]$). Homogeneous. \checkmark

\textbf{Q.3 Work-conjugacy verification (final).} The internal virtual work after IBP is

$\int_{\mathcal{B}_0}[P_{iJ}\delta F_{iJ}+\Sigma_{IJK}\delta\mathcal{G}_{IJK}]\,\mathrm{d}V_0+\text{boundary}$,

with $\delta F_{iJ}\leftrightarrow P_{iJ}$ (both rank-two, material/spatial mixed) and $\delta\mathcal{G}_{IJK}\leftrightarrow\Sigma_{IJK}$ (both rank-three, $I,J,K$ material) -- exact work-conjugate pairs; after the two IBP steps the gradient branch yields the hyperstress boundary fluxes $\mathcal{H}_{lJK}N_K\delta u_{l,J}-(\Sigma_{IJK,K}F_{iI})N_J\delta u_i$ and the domain operator $(\Sigma_{IJK,K}F_{iI})_{,J}$. Every contraction yields a scalar energy density. \checkmark

% =====================================================================
\subsection{Essential boundary conditions -- interpretation}
\label{subsec:essentialnotes}
% =====================================================================

For a second-gradient (strain-gradient) theory the energy depends on second displacement derivatives $\nabla_0\nabla_0\mathbf{u}$; the natural (weak) form therefore requires the test functions to have square-integrable second derivatives, i.e.\ $\mathbf{u}\in H^2(\mathcal{B}_0)$ -- exactly the $C^1$ regularity furnished by the BFS element (Section~\ref{sec:bfs}). Consequently the essential (Dirichlet) data specify \emph{both} the displacement and its normal derivative (the slope) on $\partial\mathcal{B}_{0u}$ \cite{AmanatidouAravas2002,EngelGarikipati2002,Toupin1964}. This is eq:essBC. The slope prescription also fixes the admissible rotation variation $\delta\boldsymbol{\theta}$ on the essential boundary, consistent with the displacement-based (no-independent-microrotation) assumption. \checkmark

% =====================================================================
\subsection{Natural boundary conditions -- interpretation}
\label{subsec:naturalnotes}
% =====================================================================

On $\partial\mathcal{B}_{0t}$ the variations are free, so their conjugates are prescribed as data. The effective single traction $(P_{iJ}-\Sigma_{IJK,K}F_{iI})N_J$ is the classical nominal traction corrected by the divergence of the double stress pulled back to the reference configuration -- the finite-strain generalisation of the Mindlin Form-II "stress minus double-stress divergence" traction \cite{Mindlin1964,MindlinEshel1968,Toupin1964}. This reduces to the classical TL traction when the gradient modulus $\mathbb{D}$ vanishes. \checkmark

% =====================================================================
\subsection{Higher-order boundary terms -- interpretation}
\label{subsec:higherordernotes}
% =====================================================================

The surface term $\Sigma_{IJK}N_K\delta E_{IJ}=\mathcal{H}_{lJK}N_K\delta u_{l,J}$ (eq:rn2:higher\_order\_bnd) is the \emph{double (higher-order) traction} virtual work. Substituting $\delta u_{l,J}=N_J\delta u_{l,N}+(\text{tangential part})$, the part proportional to the free normal-slope variation $\delta(\partial u_l/\partial N)=\delta u_{l,N}$ is the double-traction conjugate; its coefficient is $\mathcal{D}_l=\mathcal{H}_{lJK}N_KN_J$. This is the finite-strain generalisation of the classical strain-gradient "double force" / "surface double stress" of Toupin and Germain \cite{Toupin1964,Germain1973}. On $\partial\mathcal{B}_{0u}$ the slope is prescribed (essential), so $\delta(\partial u_i/\partial N)=0$ and this term vanishes -- consistent with the $C^1$ BFS essential DOFs ($u_i$ and nodal slopes). The natural DOFs of the BFS element are thus the single traction and the double traction \cite{AmanatidouAravas2002,EngelGarikipati2002}. \checkmark

% =====================================================================
\subsection{Final governing boundary-value problem -- and small-strain recovery}
\label{subsec:bvpnotes}
% =====================================================================

\paragraph{R.1 Assembled BVP.} Equations~\eqref{rn2:strong} (strong form), essential BCs (displacement + slope), natural BCs (effective single traction $(P_{iJ}-\Sigma_{IJK,K}F_{iI})N_J$ and double traction $\mathcal{D}_i=\mathcal{H}_{iJK}N_KN_J$), and initial conditions collect the complete strong-form IBVP of the geometrically nonlinear, displacement-based, $C^1$ BFS finite-strain Mindlin Form-II strain-gradient solid.

\paragraph{R.2 Small-strain reduction to the validated 2022 benchmark.}

As $\|\nabla_0\mathbf{u}\|\to0$: $F_{iI}\to\delta_{iI}$; $P_{iJ}=S_{IJ}F_{iI}\to S_{IJ}\to\sigma_{ij}$ (2nd P-K $\to$ Cauchy, both reduce to the symmetric Cauchy stress in the linear limit); $\Sigma_{IJK,K}F_{iI}\to\Sigma_{iJK,K}\to h_{ijk,k}$ (the double stress / hyperstress of the 2022 theory, carrying the $\mathbb{D}/a_i$ content of eq:psi). The strong form \eqref{rn2:strong} therefore becomes

\begin{equation}
-P_{iJ,J}+(\Sigma_{IJK,K}F_{iI})_{,J}-f_{0i}+\rho_0\ddot{u}_i=0
\;\;\xrightarrow{\text{linear limit}}\;\;
-\sigma_{ij,j}+h_{ijk,jk}-b_i+\rho\ddot{u}_i=0,
\end{equation}

or equivalently,

\begin{equation}
\sigma_{ij,j}-h_{ijk,jk}+b_i=\rho\ddot{u}_i,
\label{rn2:strong_small}
\end{equation}

the 2022 benchmark fourth-order Mindlin Form-II equilibrium. Setting $h_{ijk}=0$ yields $\sigma_{ij,j}+b_i=\rho\ddot{u}_i$, the classical Cauchy/Navier equation of linear elastodynamics. \checkmark

The boundary conditions reduce to: essential $u_i,\partial u_i/\partial n$ prescribed; natural single traction $t_i=\bar{t}_i$ and double traction $d_i=\bar{d}_i$ prescribed -- precisely the 2022 benchmark boundary-value problem. The static sub-case ($\ddot{\mathbf{u}}=\mathbf{0}$) is the exact 2022 quasi-static benchmark. \textbf{Mandatory recovery rule SATISFIED at the strong-form / boundary-value level.} \checkmark

\paragraph{R.3 Reviewer comments (Q\&A).}

\begin{enumerate}
\item \textbf{R: Why is there only one balance equation and no separate moment equation?}
A: The microrotation is slaved to the displacement through $\mathbf{R}=\mathrm{polar}(\mathbf{F})$ (assumption A2, displacement-based / compatible Form-II). The energy depends only on $\{\mathbf{E},\bm{\mathcal{G}}\}$ (finite-strain strain-gradient Form-II, with no independent couple-stress energy), so there is no separate moment balance; the single force balance contains the double-stress divergence term $(\Sigma_{IJK,K}F_{iI})_{,J}$. (A full Cosserat version with independent $\boldsymbol{\theta}$ would add one -- intentionally excluded here.)

\item \textbf{R: Sign of the inertia term in \eqref{rn2:strong}.}
A: From $\delta\Pi_{\mathrm{int}}=\delta\mathcal{W}_{\mathrm{ext}}-\delta\mathcal{W}_{\mathrm{iner}}$ (P.2) the inertia appears as $-\rho_0\ddot{u}_i$ in the "internal + body -- inertia = 0" arrangement, equivalent to $\mathrm{Div}_0\mathbf{T}+\mathbf{f}_0=\rho_0\ddot{\mathbf{u}}$. Standard. \checkmark

\item \textbf{R: Is $\mathcal{H}_{iJK}=\Sigma_{IJK}F_{iI}$ the correct hyperstress?}
A: Yes. It is the unique rank-three tensor that makes $\Sigma_{IJK}\delta\mathcal{G}_{IJK}$ reduce (after the two IBP steps) to the boundary fluxes $\mathcal{H}_{lJK}N_K\delta u_{l,J}-(\Sigma_{IJK,K}F_{iI})N_J\delta u_i$ plus the domain operator $(\Sigma_{IJK,K}F_{iI})_{,J}$. It carries the $\mathbb{D}/a_i$ content of eq:psi through $\boldsymbol{\Sigma}=\partial\psi/\partial\bm{\mathcal{G}}$. \checkmark

\item \textbf{R: Boundary-condition split for a $C^1$ element.}
A: Essential = displacement + slope (the BFS nodal DOFs); natural = single traction + double traction. This is the standard partition for second-gradient $C^1$ elements \cite{AmanatidouAravas2002,EngelGarikipati2002}. \checkmark

\item \textbf{R: Finite-strain form of the double traction.}
A: At finite strain $\delta\bm{\mathcal{G}}$ follows from $\delta\mathbf{E}$ and $\delta\mathbf{F}$; the double-traction boundary term is $\Sigma_{IJK}N_K\delta E_{IJ}=\mathcal{H}_{lJK}N_K\delta u_{l,J}$ (the hyperstress flux conjugate to the free slope variation $\delta u_{l,J}$), whose normal-slope coefficient is $\mathcal{D}_l=\mathcal{H}_{lJK}N_KN_J$. This is the finite-strain generalisation of the classical strain-gradient "double force" / "surface double stress" of Toupin and Germain \cite{Toupin1964,Germain1973}. \checkmark

\item \textbf{R: Does the strong form reduce exactly to the 2022 benchmark?}
A: Yes (R.2) -- every term maps to the 2022 Mindlin Form-II fourth-order equilibrium and its boundary conditions; the benchmark is unchanged. \checkmark
\end{enumerate}

\noindent\textbf{Closing statement for Section~3A.} Integration by parts of the unreduced Hamilton statement yields the strong-form balance \eqref{rn2:strong} (equivalently $\mathrm{Div}_0\mathbf{T}+\mathbf{f}_0=\rho_0\ddot{\mathbf{u}}$), the essential BCs (displacement + slope), the natural BCs (effective single traction $(P_{iJ}-\Sigma_{IJK,K}F_{iI})N_J$) and the higher-order double-traction condition $\mathcal{D}_i=\mathcal{H}_{iJK}N_KN_J$, and the assembled IBVP. All steps are dimensionally consistent, work-conjugate, and reduce exactly to the validated 2022 small-strain Mindlin Form-II boundary-value problem in the linear limit, satisfying the mandatory recovery rule. The explicit closed forms of $\mathbf{S},\boldsymbol{\Sigma}$ from eq:psi and the finite-element (BFS $C^1$) discretisation remain for the next section.

% =====================================================================
% APPENDIX: Complete Mathematical Derivations from the Verification Audit
% =====================================================================
\appendix

\section{Complete Mathematical Derivations from the Verification Audit}
\label{app:proofs}

This appendix records the complete step-by-step mathematical derivations performed during the independent verification audit of Version~1.0. These derivations were constructed to prove or disprove each reported issue, starting from Hamilton's principle and proceeding through every algebraic step, tensor contraction, integration by parts, and product rule. All four issues reported in the audit are proven genuine; no false positives were found.

\subsection{Common starting point: Hamilton's principle}

The variational statement is:

\begin{equation}
\delta\int_{t_1}^{t_2}\bigl[\mathcal{T}-\Pi_{\mathrm{int}}\bigr]\,\mathrm{d}t
\;+\;\int_{t_1}^{t_2}\delta\mathcal{W}_{\mathrm{ext}}\,\mathrm{d}t
\;=\;0,
\qquad
\delta\boldsymbol{\varphi}(t_1)=\delta\boldsymbol{\varphi}(t_2)=\mathbf{0}.
\label{app:hamilton}
\end{equation}

\subsection{Step A: Time integration by parts on $\mathcal{T}$}

\begin{equation}
\mathcal{T}
= \int_{\mathcal{B}_0}\tfrac{1}{2}\rho_0\dot{\varphi}_i\dot{\varphi}_i\,\mathrm{d}V_0,
\qquad
\delta\mathcal{T}
= \int_{\mathcal{B}_0}\rho_0\dot{\varphi}_i\delta\dot{\varphi}_i\,\mathrm{d}V_0.
\end{equation}

Since $\mathbf{u}=\boldsymbol{\varphi}-\mathbf{X}$ and $\mathbf{X}$ is time-independent: $\dot{\varphi}_i=\dot{u}_i$, $\delta\dot{\varphi}_i=\partial(\delta\varphi_i)/\partial t$.

\begin{align}
\int_{t_1}^{t_2}\delta\mathcal{T}\,\mathrm{d}t
&= \int_{t_1}^{t_2}\!\int_{\mathcal{B}_0}\rho_0\dot{\varphi}_i\frac{\partial(\delta\varphi_i)}{\partial t}\,\mathrm{d}V_0\,\mathrm{d}t \notag\\
&= \int_{\mathcal{B}_0}\rho_0\Bigl[\dot{\varphi}_i\delta\varphi_i\Bigr]_{t_1}^{t_2}\mathrm{d}V_0
   - \int_{t_1}^{t_2}\!\int_{\mathcal{B}_0}\rho_0\ddot{\varphi}_i\delta\varphi_i\,\mathrm{d}V_0\,\mathrm{d}t.
\end{align}

Endpoint conditions kill the first term:

\begin{equation}
= -\int_{t_1}^{t_2}\!\int_{\mathcal{B}_0}\rho_0\ddot{u}_i\delta u_i\,\mathrm{d}V_0\,\mathrm{d}t
= -\int_{t_1}^{t_2}\delta\mathcal{W}_{\mathrm{iner}}\,\mathrm{d}t.
\label{app:timeibp}
\end{equation}

\subsection{Step B: Expanded Hamilton principle}

Substituting into \eqref{app:hamilton}:

\begin{equation}
-\int_{t_1}^{t_2}\delta\mathcal{W}_{\mathrm{iner}}\,\mathrm{d}t
- \int_{t_1}^{t_2}\delta\Pi_{\mathrm{int}}\,\mathrm{d}t
+ \int_{t_1}^{t_2}\delta\mathcal{W}_{\mathrm{ext}}\,\mathrm{d}t
= 0.
\end{equation}

Since this holds for arbitrary time intervals, the integrand vanishes at each instant:

\begin{equation}
\boxed{\delta\Pi_{\mathrm{int}} = \delta\mathcal{W}_{\mathrm{ext}} - \delta\mathcal{W}_{\mathrm{iner}}}
\label{app:equilibrium}
\end{equation}

\subsection{Step C: Variation of internal energy}

\begin{equation}
\delta\Pi_{\mathrm{int}}
= \int_{\mathcal{B}_0}\delta\psi\,\mathrm{d}V_0
= \int_{\mathcal{B}_0}\bigl[S_{IJ}\delta E_{IJ} + \Sigma_{IJK}\delta\mathcal{G}_{IJK}\bigr]\,\mathrm{d}V_0,
\end{equation}

where $\mathbf{S}=\partial\psi/\partial\mathbf{E}$ and $\boldsymbol{\Sigma}=\partial\psi/\partial\boldsymbol{\mathcal{G}}$.

\subsection{Step D: Variation chain}

\begin{align}
\delta F_{iI} &= \delta u_{i,I} \qquad(\text{since }F_{iI}=\delta_{iI}+u_{i,I}),\\
\delta E_{IJ} &= \tfrac{1}{2}\bigl(F_{kI}\delta u_{k,J}+F_{kJ}\delta u_{k,I}\bigr),\\
\delta\mathcal{G}_{IJK} &= (\delta E_{IJ})_{,K}.
\end{align}

\subsection{Step E: Classical branch $S_{IJ}\delta E_{IJ}$}

\begin{align}
S_{IJ}\delta E_{IJ}
&= \tfrac{1}{2}S_{IJ}\bigl(F_{kI}\delta u_{k,J}+F_{kJ}\delta u_{k,I}\bigr).
\end{align}

Relabel dummy indices $I\leftrightarrow J$ in the second term:
\begin{equation}
\tfrac{1}{2}S_{IJ}F_{kJ}\delta u_{k,I}
\;\xrightarrow{I\to J',\,J\to I'}\;
\tfrac{1}{2}S_{J'I'}F_{kI'}\delta u_{k,J'}
= \tfrac{1}{2}S_{I'J'}F_{kI'}\delta u_{k,J'},
\end{equation}
using $S_{J'I'}=S_{I'J'}$. Rename $I'\to I$, $J'\to J$:
\begin{equation}
= \tfrac{1}{2}S_{IJ}F_{kI}\delta u_{k,J}.
\end{equation}

Both terms equal, therefore:
\begin{equation}
S_{IJ}\delta E_{IJ} = S_{IJ}F_{kI}\delta u_{k,J} = P_{kJ}\delta u_{k,J}
= P_{iJ}\delta u_{i,J},
\end{equation}
where $P_{iJ}:=S_{IJ}F_{iI}$. Integration by parts:

\begin{equation}
\int_{\mathcal{B}_0}P_{iJ}\delta u_{i,J}\,\mathrm{d}V_0
= \int_{\partial\mathcal{B}_0}P_{iJ}N_J\delta u_i\,\mathrm{d}A_0
- \int_{\mathcal{B}_0}P_{iJ,J}\delta u_i\,\mathrm{d}V_0.
\label{app:classical_ibp}
\end{equation}

\subsection{Step F: Gradient branch $\Sigma_{IJK}\delta\mathcal{G}_{IJK}$}

\begin{equation}
\int_{\mathcal{B}_0}\Sigma_{IJK}\delta\mathcal{G}_{IJK}\,\mathrm{d}V_0
= \int_{\mathcal{B}_0}\Sigma_{IJK}(\delta E_{IJ})_{,K}\,\mathrm{d}V_0.
\end{equation}

\textbf{First IBP over $X_K$:}

\begin{align}
&= \int_{\partial\mathcal{B}_0}\Sigma_{IJK}N_K\delta E_{IJ}\,\mathrm{d}A_0
   - \int_{\mathcal{B}_0}\Sigma_{IJK,K}\delta E_{IJ}\,\mathrm{d}V_0.
\label{app:grad_ibp1}
\end{align}

Substitute $\delta E_{IJ}$:
\begin{equation}
-\int_{\mathcal{B}_0}\Sigma_{IJK,K}\cdot\tfrac{1}{2}\bigl(F_{iI}\delta u_{i,J}+F_{iJ}\delta u_{i,I}\bigr)\,\mathrm{d}V_0.
\end{equation}

In the second half, relabel $I\leftrightarrow J$:
\begin{equation}
-\tfrac{1}{2}\int_{\mathcal{B}_0}\Sigma_{IJK,K}F_{iJ}\delta u_{i,I}\,\mathrm{d}V_0
\;\xrightarrow{I\leftrightarrow J}\;
-\tfrac{1}{2}\int_{\mathcal{B}_0}\Sigma_{JIK,K}F_{iI}\delta u_{i,J}\,\mathrm{d}V_0
= -\tfrac{1}{2}\int_{\mathcal{B}_0}\Sigma_{IJK,K}F_{iI}\delta u_{i,J}\,\mathrm{d}V_0,
\end{equation}
using $\Sigma_{JIK}=\Sigma_{IJK}$. Both halves equal:

\begin{equation}
-\int_{\mathcal{B}_0}\Sigma_{IJK,K}\delta E_{IJ}\,\mathrm{d}V_0
= -\int_{\mathcal{B}_0}\Sigma_{IJK,K}F_{iI}\delta u_{i,J}\,\mathrm{d}V_0.
\end{equation}

\textbf{Second IBP over $X_J$:}

Define $A_{iJ}:=\Sigma_{IJK,K}F_{iI}$:

\begin{align}
-\int_{\mathcal{B}_0}A_{iJ}\delta u_{i,J}\,\mathrm{d}V_0
&= -\int_{\partial\mathcal{B}_0}A_{iJ}N_J\delta u_i\,\mathrm{d}A_0
   + \int_{\mathcal{B}_0}A_{iJ,J}\delta u_i\,\mathrm{d}V_0 \notag\\
&= -\int_{\partial\mathcal{B}_0}(\Sigma_{IJK,K}F_{iI})N_J\delta u_i\,\mathrm{d}A_0
   + \int_{\mathcal{B}_0}(\Sigma_{IJK,K}F_{iI})_{,J}\delta u_i\,\mathrm{d}V_0.
\label{app:grad_ibp2}
\end{align}

\textbf{Summary of gradient branch:}

\begin{align}
\int_{\mathcal{B}_0}\Sigma_{IJK}\delta\mathcal{G}_{IJK}\,\mathrm{d}V_0
&= \underbrace{\int_{\partial\mathcal{B}_0}\Sigma_{IJK}N_K\delta E_{IJ}\,\mathrm{d}A_0}_{\text{(a)}}
   - \underbrace{\int_{\partial\mathcal{B}_0}(\Sigma_{IJK,K}F_{iI})N_J\delta u_i\,\mathrm{d}A_0}_{\text{(b)}}
   + \underbrace{\int_{\mathcal{B}_0}(\Sigma_{IJK,K}F_{iI})_{,J}\delta u_i\,\mathrm{d}V_0}_{\text{domain}}.
\label{app:grad_summary}
\end{align}

\subsection{Step G: Assembly of the domain (strong-form) equation}

Collect all domain terms from \eqref{app:classical_ibp} and \eqref{app:grad_summary}:

\begin{equation}
\delta\Pi_{\mathrm{int}}\big|_{\mathrm{domain}}
= \int_{\mathcal{B}_0}\bigl[-P_{iJ,J}+(\Sigma_{IJK,K}F_{iI})_{,J}\bigr]\delta u_i\,\mathrm{d}V_0.
\end{equation}

Domain parts of the external and inertia terms:
\begin{equation}
\delta\mathcal{W}_{\mathrm{ext}}\big|_{\mathrm{domain}} = \int_{\mathcal{B}_0}f_{0i}\delta u_i\,\mathrm{d}V_0,
\qquad
\delta\mathcal{W}_{\mathrm{iner}} = \int_{\mathcal{B}_0}\rho_0\ddot{u}_i\delta u_i\,\mathrm{d}V_0.
\end{equation}

Substitute into \eqref{app:equilibrium}:

\begin{equation}
\int_{\mathcal{B}_0}\bigl[-P_{iJ,J}+(\Sigma_{IJK,K}F_{iI})_{,J}\bigr]\delta u_i\,\mathrm{d}V_0
= \int_{\mathcal{B}_0}f_{0i}\delta u_i\,\mathrm{d}V_0
  - \int_{\mathcal{B}_0}\rho_0\ddot{u}_i\delta u_i\,\mathrm{d}V_0.
\end{equation}

Since $\delta u_i$ is arbitrary in $\mathcal{B}_0$:

\begin{equation}
-P_{iJ,J}+(\Sigma_{IJK,K}F_{iI})_{,J} = f_{0i}-\rho_0\ddot{u}_i.
\end{equation}

Rearranging to $=0$:

\begin{equation}
\boxed{-P_{iJ,J}+(\Sigma_{IJK,K}F_{iI})_{,J}-f_{0i}+\rho_0\ddot{u}_i=0}
\label{app:strong_correct}
\end{equation}

Multiply by $-1$:

\begin{equation}
P_{iJ,J}-(\Sigma_{IJK,K}F_{iI})_{,J}+f_{0i} = \rho_0\ddot{u}_i
\quad\Longleftrightarrow\quad
\mathrm{Div}_0\mathbf{T}+\mathbf{f}_0=\rho_0\ddot{\mathbf{u}}.
\label{app:balance}
\end{equation}

This is the standard balance of linear momentum \cite{Holzapfel2000,BonetWood2008,SimoHughes1998}. $\checkmark$

%-------------------------------------------------------------------
\subsection{Proof of Issue 1: Sign error in (rn2:strong)}
\label{app:proof1}

The manuscript's equation (rn2:strong) in Version~1.0 stated:

\begin{equation}
-P_{iJ,J}+(\Sigma_{IJK,K}F_{iI})_{,J}+f_{0i}-\rho_0\ddot{u}_i=0.
\label{app:manuscript_strong}
\end{equation}

Comparing term-by-term with the derivation \eqref{app:strong_correct}:

\begin{itemize}
\item Stress divergence: $-P_{iJ,J}$ matches. $\checkmark$
\item Gradient divergence: $+(\Sigma_{IJK,K}F_{iI})_{,J}$ matches. $\checkmark$
\item Body force: derivation gives $-f_{0i}$; manuscript has $+f_{0i}$. \textbf{Reversed.}
\item Inertia: derivation gives $+\rho_0\ddot{u}_i$; manuscript has $-\rho_0\ddot{u}_i$. \textbf{Reversed.}
\end{itemize}

Physical verification: multiply \eqref{app:manuscript_strong} by $-1$:
\begin{equation}
P_{iJ,J}-(\Sigma_{IJK,K}F_{iI})_{,J}-f_{0i}=-\rho_0\ddot{u}_i
\;\;\Longrightarrow\;\;
\mathrm{Div}_0\mathbf{T}-\mathbf{f}_0=-\rho_0\ddot{\mathbf{u}}.
\end{equation}

Static limit ($\ddot{\mathbf{u}}=\mathbf{0}$): $\mathrm{Div}_0\mathbf{T}=\mathbf{f}_0$. This says stress divergence \emph{equals} body force in the same direction, rather than \emph{opposing} it. For a body under gravity ($\mathbf{f}_0=\rho_0\mathbf{g}$), resting on a support ($\ddot{\mathbf{u}}=\mathbf{0}$, $\mathbf{P}=\mathbf{0}$): $0=\rho_0\mathbf{g}$, which is nonsensical.

\textbf{Verdict:} Issue~1 is a \textbf{genuine CRITICAL error}. The signs of $f_{0i}$ and $\rho_0\ddot{u}_i$ in the manuscript's (rn2:strong) are reversed. This has been corrected in \eqref{rn2:strong} of the present version.

%-------------------------------------------------------------------
\subsection{Proof of Issue 2: Sign error in (rn2:strong\_small)}
\label{app:proof2}

The manuscript's (rn2:strong\_small) in Version~1.0 stated:

\begin{equation}
-\sigma_{ij,j}+h_{ijk,jk}+b_i=\rho\ddot{u}_i.
\label{app:manuscript_small}
\end{equation}

Apply the small-strain limit to the correct derivation \eqref{app:strong_correct}:

\textbf{Limit of each term:}
\begin{itemize}
\item $F_{iI}=\delta_{iI}+u_{i,I}\to\delta_{iI}$.
\item $P_{iJ}=S_{IJ}F_{iI}\to S_{IJ}\delta_{iI}=S_{iJ}\to\sigma_{ij}$. Therefore $P_{iJ,J}\to\sigma_{ij,j}$.
\item $\Sigma_{IJK,K}F_{iI}\to\Sigma_{IJK,K}\delta_{iI}=\Sigma_{iJK,K}\to h_{ijk,k}$. Therefore $(\Sigma_{IJK,K}F_{iI})_{,J}\to h_{ijk,jk}$.
\item $f_{0i}\to b_i$.
\item $\rho_0\ddot{u}_i\to\rho\ddot{u}_i$.
\end{itemize}

Substitute into \eqref{app:strong_correct}:

\begin{equation}
-\sigma_{ij,j}+h_{ijk,jk}-b_i+\rho\ddot{u}_i=0
\quad\Longleftrightarrow\quad
\boxed{\sigma_{ij,j}-h_{ijk,jk}+b_i=\rho\ddot{u}_i}.
\label{app:correct_small}
\end{equation}

Setting $h_{ijk}=0$: $\sigma_{ij,j}+b_i=\rho\ddot{u}_i$, the classical Cauchy/Navier equation. $\checkmark$

The manuscript's \eqref{app:manuscript_small} with $h_{ijk}=0$ gives $-\sigma_{ij,j}+b_i=\rho\ddot{u}_i$, i.e., $\sigma_{ij,j}=b_i-\rho\ddot{u}_i$, which has wrong signs on both $\sigma_{ij,j}$ and $b_i$.

\textbf{Verdict:} Issue~2 is a \textbf{genuine CRITICAL error}. The signs of $\sigma_{ij,j}$ and $h_{ijk,jk}$ are reversed in the manuscript's (rn2:strong\_small). This has been corrected to \eqref{rn2:strong_small} in the present version.

%-------------------------------------------------------------------
\subsection{Proof of Issue 3: Incorrect boundary-term identification (b)}
\label{app:proof3}

The document defines $\mathcal{H}_{iJK}:=\Sigma_{IJK}F_{iI}$ and claimed (Version~1.0):

\begin{equation}
(\Sigma_{IJK,K}F_{iI})N_J\delta u_i
\stackrel{?}{=}
\mathcal{H}_{iJK,K}N_J\delta u_i.
\label{app:claim3}
\end{equation}

Compute $\mathcal{H}_{iJK,K}$ by the product rule:

\begin{align}
\mathcal{H}_{iJK,K}
&= \frac{\partial}{\partial X_K}(\Sigma_{IJK}F_{iI}) \notag\\
&= \frac{\partial\Sigma_{IJK}}{\partial X_K}F_{iI}
   + \Sigma_{IJK}\frac{\partial F_{iI}}{\partial X_K} \notag\\
&= \Sigma_{IJK,K}F_{iI} + \Sigma_{IJK}F_{iI,K}.
\label{app:product_rule_H}
\end{align}

Since $F_{iI}=\delta_{iI}+u_{i,I}$:

\begin{equation}
F_{iI,K}=u_{i,I,K}.
\end{equation}

Therefore:

\begin{equation}
\boxed{\mathcal{H}_{iJK,K} = \Sigma_{IJK,K}F_{iI} + \Sigma_{IJK}u_{i,I,K}}
\label{app:exact_H_div}
\end{equation}

\textbf{Order analysis of the extra term:}
\begin{itemize}
\item $\Sigma_{IJK}=\partial\psi/\partial\mathcal{G}_{IJK}=O(\|\mathbb{D}\|\cdot\|\nabla_0^2\mathbf{u}\|)$
\item $u_{i,I,K}=O(\|\nabla_0^2\mathbf{u}\|)$
\item Product: $\Sigma_{IJK}u_{i,I,K}=O(\|\mathbb{D}\|\cdot\|\nabla_0^2\mathbf{u}\|^2)$, which is \textbf{non-zero} at finite strain whenever $\mathbb{D}\neq\mathbf{0}$.
\end{itemize}

\textbf{Difference:}

\begin{equation}
\mathcal{H}_{iJK,K}N_J\delta u_i - \Sigma_{IJK,K}F_{iI}N_J\delta u_i
= \Sigma_{IJK}u_{i,I,K}N_J\delta u_i \neq 0.
\end{equation}

\textbf{Small-strain check:} As $\|\nabla_0\mathbf{u}\|\to0$: $u_{i,I,K}\to 0$ and $\Sigma_{IJK}\to 0$ simultaneously, so the product $\Sigma_{IJK}u_{i,I,K}=O(\|\nabla_0\mathbf{u}\|^2)$ vanishes. The error is invisible in the benchmark recovery.

\textbf{Verdict:} Issue~3 is a \textbf{genuine MAJOR error}. The product rule generates an extra term $\Sigma_{IJK}F_{iI,K}$ that the manuscript dropped. This makes the finite-strain natural boundary condition incorrect. The corrected boundary term (b) is $-(\Sigma_{IJK,K}F_{iI})N_J\delta u_i$ (not $-\mathcal{H}_{iJK,K}N_J\delta u_i$). This has been corrected in Section~P.6 of the present version.

%-------------------------------------------------------------------
\subsection{Proof of Issue 4: Scalar $\neq$ tensor in boundary-term identification (a)}
\label{app:proof4}

The manuscript claimed (Version~1.0):

\begin{equation}
\Sigma_{IJK}N_K\delta E_{IJ}
\stackrel{?}{=}
\mathcal{H}_{iJK}N_K\delta E_{IJ}.
\label{app:claim4}
\end{equation}

\textbf{Index analysis of the LHS:} $\Sigma_{IJK}N_K\delta E_{IJ}$.

\begin{itemize}
\item $I$: appears in $\Sigma_{IJK}$ and $\delta E_{IJ}$ $\Rightarrow$ 2 times $\Rightarrow$ summed.
\item $J$: appears in $\Sigma_{IJK}$ and $\delta E_{IJ}$ $\Rightarrow$ 2 times $\Rightarrow$ summed.
\item $K$: appears in $\Sigma_{IJK}$ and $N_K$ $\Rightarrow$ 2 times $\Rightarrow$ summed.
\end{itemize}

All indices contracted. \textbf{The LHS is a scalar.} $\checkmark$

\textbf{Index analysis of the RHS:} Substitute $\mathcal{H}_{iJK}=\Sigma_{LJK}F_{iL}$:

\begin{equation}
\mathcal{H}_{iJK}N_K\delta E_{IJ} = \Sigma_{LJK}F_{iL}N_K\delta E_{IJ}.
\end{equation}

\begin{itemize}
\item $L$: appears in $\Sigma_{LJK}$ and $F_{iL}$ $\Rightarrow$ 2 times $\Rightarrow$ summed.
\item $K$: appears in $\Sigma_{LJK}$ (via $\mathcal{H}$) and $N_K$ $\Rightarrow$ 2 times $\Rightarrow$ summed.
\item $J$: appears in $\mathcal{H}_{iJK}$ and $\delta E_{IJ}$ $\Rightarrow$ 2 times $\Rightarrow$ summed.
\item $i$: appears only in $\mathcal{H}_{iJK}$ $\Rightarrow$ \textbf{1 time $\Rightarrow$ FREE}.
\item $I$: appears only in $\delta E_{IJ}$ $\Rightarrow$ \textbf{1 time $\Rightarrow$ FREE}.
\end{itemize}

The RHS has two free indices $(i,I)$. \textbf{It is a second-order tensor, not a scalar.}

A scalar cannot equal a second-order tensor. The claimed equality \eqref{app:claim4} is \textbf{algebraically impossible}.

\textbf{The correct rewriting:}

\begin{align}
\Sigma_{IJK}N_K\delta E_{IJ}
&= \Sigma_{IJK}N_K\cdot\tfrac{1}{2}(F_{lI}\delta u_{l,J}+F_{lJ}\delta u_{l,I}) \notag\\
&= \tfrac{1}{2}\Sigma_{IJK}F_{lI}N_K\delta u_{l,J}
   + \tfrac{1}{2}\Sigma_{IJK}F_{lJ}N_K\delta u_{l,I}.
\end{align}

In the second term, relabel $I\leftrightarrow J$:

\begin{equation}
\tfrac{1}{2}\Sigma_{IJK}F_{lJ}N_K\delta u_{l,I}
\;\xrightarrow{I\to J',\,J\to I'}\;
\tfrac{1}{2}\Sigma_{J'I'K}F_{lI'}N_K\delta u_{l,J'}
= \tfrac{1}{2}\Sigma_{I'J'K}F_{lI'}N_K\delta u_{l,J'},
\end{equation}

using $\Sigma_{J'I'K}=\Sigma_{I'J'K}$. Rename:

\begin{equation}
= \tfrac{1}{2}\Sigma_{IJK}F_{lI}N_K\delta u_{l,J}.
\end{equation}

Both terms equal:

\begin{equation}
\Sigma_{IJK}N_K\delta E_{IJ}
= \Sigma_{IJK}F_{lI}N_K\delta u_{l,J}
= \mathcal{H}_{lJK}N_K\delta u_{l,J}.
\label{app:correct_bnd_a}
\end{equation}

Index check on $\mathcal{H}_{lJK}N_K\delta u_{l,J}$:

\begin{itemize}
\item $l$: $\mathcal{H}_{lJK}$, $\delta u_{l,J}$ $\Rightarrow$ 2 times $\Rightarrow$ summed.
\item $J$: $\mathcal{H}_{lJK}$, $\delta u_{l,J}$ $\Rightarrow$ 2 times $\Rightarrow$ summed.
\item $K$: $\mathcal{H}_{lJK}$, $N_K$ $\Rightarrow$ 2 times $\Rightarrow$ summed.
\end{itemize}

All indices contracted. \textbf{Scalar.} $\checkmark$

\textbf{Double traction:} Decompose $\delta u_{l,J}$ on the boundary:

\begin{equation}
\delta u_{l,J} = N_J\delta u_{l,N} + (\delta_{JM}-N_JN_M)\delta u_{l,M}.
\end{equation}

The normal-slope part:

\begin{equation}
\mathcal{H}_{lJK}N_KN_J\delta u_{l,N}
= \mathcal{D}_l\,\delta u_{l,N},
\qquad
\mathcal{D}_l = \mathcal{H}_{lJK}N_KN_J.
\end{equation}

The double traction $\mathcal{D}_i=\mathcal{H}_{iJK}N_KN_J$ is correctly obtained from the correct intermediate form \eqref{app:correct_bnd_a}.

\textbf{Verdict:} Issue~4 is a \textbf{genuine MAJOR error}. The expression $\mathcal{H}_{iJK}N_K\delta E_{IJ}$ has two free indices $(i,I)$ and is a second-order tensor, while $\Sigma_{IJK}N_K\delta E_{IJ}$ is a scalar. The equality is algebraically impossible. The correct substitution yields $\mathcal{H}_{lJK}N_K\delta u_{l,J}$ (eq:\eqref{rn2:higher_order_bnd} of the present version).

%-------------------------------------------------------------------
\subsection{Summary of verification verdicts}

\begin{table}[h]
\centering
\caption{Verification results for Issues 1--4.}
\label{tab:verdicts}
\begin{tabular}{lll}
\toprule
Issue & Verdict & Proof basis \\
\midrule
1 (strong form signs) & Genuine CRITICAL & Derivation Steps A--G \\
2 (small-strain signs) & Genuine CRITICAL & Small-strain limit of Steps A--G \\
3 (boundary term b) & Genuine MAJOR & Product-rule expansion \eqref{app:product_rule_H} \\
4 (boundary term a) & Genuine MAJOR & Free-index analysis on both sides \\
\bottomrule
\end{tabular}
\end{table}

All four issues are mathematically proven. No false positives were found in Issues 1--4.

% =====================================================================
% SECTION 4 — FULL MATHEMATICAL DERIVATION
% Self-contained derivation suitable for independent verification.
% Every algebraic step, tensor identity, IBP, and boundary term
% is shown explicitly.
% Uses ONLY frozen notation from Sections 1-3 (Master Version 1.2).
% Compile with pdflatex + bibtex (or Overleaf).
% =====================================================================


% =====================================================================
% SECTION 4 — FULL MATHEMATICAL DERIVATION
% Self-contained derivation suitable for independent verification.
% Every algebraic step, tensor identity, IBP, and boundary term
% is shown explicitly.
% Uses ONLY frozen notation from Sections 1-3 (Master Version 1.2).
% Compile with pdflatex + bibtex (or Overleaf).
% =====================================================================
\documentclass[11pt]{article}

\usepackage[margin=1in]{geometry}
\usepackage{amsmath,amssymb,bm}
\usepackage{booktabs}
\usepackage{cite}
\usepackage{array}

\DeclareMathOperator{\tr}{tr}
\DeclareMathOperator{\sym}{sym}

\newcommand{\gradz}{\nabla_{0}}
\newcommand{\bX}{\mathbf{X}}
\newcommand{\bx}{\mathbf{x}}
\newcommand{\bp}{\boldsymbol{\varphi}}
\newcommand{\bu}{\mathbf{u}}
\newcommand{\bF}{\mathbf{F}}
\newcommand{\bC}{\mathbf{C}}
\newcommand{\bb}{\mathbf{b}}
\newcommand{\bE}{\mathbf{E}}
\newcommand{\bep}{\boldsymbol{\varepsilon}}
\newcommand{\bH}{\mathbf{H}}
\newcommand{\bG}{\bm{\mathcal{G}}}
\newcommand{\bR}{\mathbf{R}}
\newcommand{\bS}{\mathbf{S}}
\newcommand{\bSig}{\boldsymbol{\Sigma}}
\newcommand{\bD}{\mathbb{D}}

\title{Section~4: Variational Formulation and Weak Form \\
\large Complete Mathematical Derivation (Self-Contained)}
\author{Internal working document -- full derivation for independent verification}
\date{}

===========================
\section{Introduction and scope}
% =====================================================================

This document provides the complete, self-contained mathematical
derivation of the variational formulation and weak form for a
geometrically nonlinear, displacement-based, $C^{1}$ Bogner--Fox--Schmit
finite element formulation of Mindlin Form-II strain-gradient
elasticity. Every algebraic manipulation, tensor contraction, index
operation, integration-by-parts step, product-rule expansion, and
boundary contribution is shown explicitly.

The derivation starts solely from the frozen mathematical framework
of Sections~1--3 of Master Version~1.2. No new constitutive assumptions
or kinematic definitions are introduced.

% =====================================================================
\section{Frozen definitions inherited from Sections~1--3}
% =====================================================================

\subsection{Kinematic definitions (Section~1)}

\textbf{Motion and displacement:}
\begin{align}
\bx &= \bp(\bX,t), \qquad \bu(\bX,t) = \bp(\bX,t) - \bX.
\end{align}

\textbf{Deformation gradient:}
\begin{align}
F_{iJ} &= \frac{\partial x_{i}}{\partial X_{J}}
= \delta_{iJ} + u_{i,J},
\qquad
J = \det\mathbf{F} > 0.
\end{align}

\textbf{Right Cauchy--Green tensor and Green--Lagrange strain:}
\begin{align}
C_{IJ} &= F_{kI}F_{kJ},
\\
E_{IJ} &= \tfrac12(C_{IJ}-\delta_{IJ})
= \tfrac12(F_{kI}F_{kJ}-\delta_{IJ}).
\end{align}

\textbf{Displacement form of strain:}

Substituting $F_{kI}=\delta_{kI}+u_{k,I}$:
\begin{align}
F_{kI}F_{kJ} &= (\delta_{kI}+u_{k,I})(\delta_{kJ}+u_{k,J})
\nonumber\\
&= \delta_{kI}\delta_{kJ} + \delta_{kI}u_{k,J}
   + u_{k,I}\delta_{kJ} + u_{k,I}u_{k,J}
\nonumber\\
&= \delta_{IJ} + u_{I,J} + u_{J,I} + u_{k,I}u_{k,J}.
\end{align}

Therefore:
\begin{equation}
E_{IJ} = \tfrac12\bigl(u_{I,J} + u_{J,I} + u_{K,I}u_{K,J}\bigr).
\end{equation}

\textbf{Strain gradient:}
\begin{equation}
\mathcal{G}_{IJK} = \frac{\partial E_{IJ}}{\partial X_{K}}
= E_{IJ,K}.
\end{equation}

Symmetry $\mathcal{G}_{IJK}=\mathcal{G}_{JIK}$ follows directly
from $E_{IJ}=E_{JI}$.

Component expansion (product rule, mixed partials commute):
\begin{align}
\mathcal{G}_{IJK}
&= \tfrac12\bigl(u_{I,J}+u_{J,I}+u_{K,I}u_{K,J}\bigr)_{,K}
\nonumber\\
&= \tfrac12\bigl(u_{I,JK}+u_{J,IK}+u_{K,IK}u_{K,J}
   +u_{K,I}u_{K,JK}\bigr)
\nonumber\\
&= \tfrac12\bigl(u_{J,IK}+u_{I,JK}
   +u_{k,IK}u_{k,J}+u_{k,I}u_{k,JK}\bigr).
\end{align}

\subsection{Constitutive definitions (Section~2)}

\textbf{Helmholtz free-energy density:}
\begin{equation}
\psi(\mathbf{E},\bm{\mathcal{G}})
= \frac{\lambda}{2}(\tr\mathbf{E})^{2}
  + \mu\,\mathbf{E}:\mathbf{E}
  + \frac{1}{2}\mathbb{D}_{ABCPQR}\,
    \mathcal{G}_{ABC}\mathcal{G}_{PQR}.
\end{equation}

\textbf{Work-conjugate stresses:}
\begin{align}
S_{IJ} &:= \frac{\partial\psi}{\partial E_{IJ}}
\quad\text{(second Piola--Kirchhoff stress)},
\\
\Sigma_{IJK} &:= \frac{\partial\psi}{\partial\mathcal{G}_{IJK}}
\quad\text{(double stress)}.
\end{align}

\textbf{Symmetry of $\boldsymbol{\Sigma}$:}

Since $\psi$ depends on $\mathcal{G}_{ABC}$ only through the
combination $\mathbb{D}_{ABCPQR}\mathcal{G}_{ABC}\mathcal{G}_{PQR}$,
and $\mathbb{D}_{ABCPQR}=\mathbb{D}_{BACPQR}$ (symmetry inherited
from $\mathcal{G}_{ABC}=\mathcal{G}_{BAC}$):

\begin{align}
\Sigma_{IJK}
= \frac{\partial\psi}{\partial\mathcal{G}_{IJK}}
= \mathbb{D}_{IJKPQR}\,\mathcal{G}_{PQR}.
\end{align}

Then:
\begin{align}
\Sigma_{JIK}
= \mathbb{D}_{JIKPQR}\,\mathcal{G}_{PQR}
= \mathbb{D}_{IJKPQR}\,\mathcal{G}_{PQR}
= \Sigma_{IJK},
\end{align}

using $\mathbb{D}_{JIKPQR}=\mathbb{D}_{IJKPQR}$. Therefore
$\Sigma_{IJK}=\Sigma_{JIK}$. $\checkmark$

\textbf{Major symmetry of $\mathbb{D}$:}

$\mathbb{D}_{ABCPQR}=\mathbb{D}_{PQRABC}$ follows from the scalar
quadratic form:
$\tfrac12\mathbb{D}_{ABCPQR}\mathcal{G}_{ABC}\mathcal{G}_{PQR}
= \tfrac12\mathbb{D}_{PQRABC}\mathcal{G}_{PQR}\mathcal{G}_{ABC}$
(relabeling $A\leftrightarrow P$, $B\leftrightarrow Q$,
$C\leftrightarrow R$).

\textbf{First Piola--Kirchhoff stress and hyperstress:}
\begin{align}
P_{iJ} &:= S_{IJ}\,F_{iI}
\quad\text{(first Piola--Kirchhoff stress)},
\\
\mathcal{H}_{iJK} &:= \Sigma_{IJK}\,F_{iI}
\quad\text{(hyperstress)}.
\end{align}

$P_{iJ}$ carries one spatial index $i$ and one material index $J$
(two-point tensor). $\mathcal{H}_{iJK}$ carries one spatial index $i$
and two material indices $J,K$.

\textbf{Symmetry of $S_{IJ}$:}

Since $\psi$ depends on $E_{IJ}$ through $\tr\mathbf{E}=E_{AA}$ and
$\mathbf{E}:\mathbf{E}=E_{AB}E_{AB}$:
\begin{align}
S_{IJ} &= \frac{\partial\psi}{\partial E_{IJ}}
= \lambda\,E_{AA}\,\frac{\partial E_{BB}}{\partial E_{IJ}}
  + 2\mu\,E_{AB}\,\frac{\partial E_{AB}}{\partial E_{IJ}}
\nonumber\\
&= \lambda\,(\tr\mathbf{E})\,\delta_{IJ} + 2\mu\,E_{IJ}.
\end{align}

Since $\delta_{IJ}=\delta_{JI}$ and $E_{IJ}=E_{JI}$:
$S_{IJ}=\lambda(\tr\mathbf{E})\delta_{IJ}+2\mu E_{IJ}
= \lambda(\tr\mathbf{E})\delta_{JI}+2\mu E_{JI} = S_{JI}$. $\checkmark$

\subsection{Variational objects (Section~3)}

\textbf{Hamilton's principle:}
\begin{equation}
\delta\int_{t_{1}}^{t_{2}}\bigl[\mathcal{T}-\Pi_{\mathrm{int}}\bigr]
\,\mathrm{d}t
+ \int_{t_{1}}^{t_{2}}\delta\mathcal{W}_{\mathrm{ext}}\,\mathrm{d}t
= 0,
\end{equation}

for admissible $\delta\bp$ with $\delta\bp(t_{1})=\delta\bp(t_{2})=\mathbf{0}$
and $\delta\bp|_{\partial\mathcal{B}_{0u}}=\mathbf{0}$.

\textbf{Kinetic energy:}
\begin{equation}
\mathcal{T} = \int_{\mathcal{B}_{0}}
\tfrac12\,\rho_{0}\,\dot{\bu}\cdot\dot{\bu}\,
\mathrm{d}V_{0}.
\end{equation}

\textbf{Internal energy:}
\begin{equation}
\Pi_{\mathrm{int}} = \int_{\mathcal{B}_{0}}
\psi\bigl(\mathbf{E}(\bu),\bm{\mathcal{G}}(\bu)\bigr)\,
\mathrm{d}V_{0}.
\end{equation}

\textbf{External virtual work:}
\begin{equation}
\delta\mathcal{W}_{\mathrm{ext}}
= \int_{\mathcal{B}_{0}} f_{0i}\,\delta u_{i}\,\mathrm{d}V_{0}
+ \int_{\partial\mathcal{B}_{0t}} \bar{T}_{0i}\,\delta u_{i}\,
  \mathrm{d}A_{0}.
\end{equation}

\textbf{Inertia virtual work:}
\begin{equation}
\delta\mathcal{W}_{\mathrm{iner}}
= \int_{\mathcal{B}_{0}} \rho_{0}\,\ddot{u}_{i}\,\delta u_{i}\,
  \mathrm{d}V_{0}.
\end{equation}

% =====================================================================
\section{Derivation of the instantaneous variational statement}
% =====================================================================

Starting from Hamilton's principle \cite{Hamilton1834,Goldstein2002} for the finite-strain Mindlin Form-II strain-gradient solid \cite{Gurtin1981,MarsdenHughes1983}.

\subsection{Time integration by parts on $\mathcal{T}$}

The variational framework follows the virtual power method of Germain \cite{Germain1973} and Maugin \cite{Maugin1980}, extended to finite strains within the Total Lagrangian framework \cite{Holzapfel2000,BonetWood2008,SimoHughes1998,Bathe2006}.

\textbf{Step 1:} Compute $\delta\mathcal{T}$ at fixed time $t$.

\begin{align}
\delta\mathcal{T}
&= \delta\int_{\mathcal{B}_{0}}
   \tfrac12\,\rho_{0}\,\dot{\varphi}_{i}\,\dot{\varphi}_{i}\,
   \mathrm{d}V_{0}
\nonumber\\
&= \int_{\mathcal{B}_{0}}
   \rho_{0}\,\dot{\varphi}_{i}\,\delta\dot{\varphi}_{i}\,
   \mathrm{d}V_{0},
\end{align}

since $\delta(\dot{\varphi}_{i}\dot{\varphi}_{i})
= 2\dot{\varphi}_{i}\delta\dot{\varphi}_{i}$ and $\rho_{0}$ is
independent of $\bu$.

\textbf{Step 2:} Show $\delta\dot{\varphi}_{i}
= \partial(\delta\varphi_{i})/\partial t$.

Since $\delta$ and $\partial/\partial t$ act on different variables
(variation vs.\ time, with $\bX$ held fixed):
\begin{equation}
\delta\dot{\varphi}_{i}
= \delta\!\left(\frac{\partial\varphi_{i}}{\partial t}\bigg|_{\bX}\right)
= \frac{\partial}{\partial t}\bigg|_{\bX}(\delta\varphi_{i})
= \frac{\partial(\delta\varphi_{i})}{\partial t}.
\end{equation}

\textbf{Step 3:} Integrate $\int\delta\mathcal{T}\,\mathrm{d}t$
by parts in time.

\begin{align}
\int_{t_{1}}^{t_{2}}\!\delta\mathcal{T}\,\mathrm{d}t
&= \int_{t_{1}}^{t_{2}}\!\int_{\mathcal{B}_{0}}
   \rho_{0}\,\dot{\varphi}_{i}\,
   \frac{\partial(\delta\varphi_{i})}{\partial t}\,
   \mathrm{d}V_{0}\,\mathrm{d}t.
\end{align}

Integration by parts in $t$:
\begin{align}
&= \int_{\mathcal{B}_{0}}\rho_{0}
   \left[\int_{t_{1}}^{t_{2}}
   \dot{\varphi}_{i}\,
   \frac{\partial(\delta\varphi_{i})}{\partial t}\,
   \mathrm{d}t\right]\mathrm{d}V_{0}
\nonumber\\
&= \int_{\mathcal{B}_{0}}\rho_{0}
   \Bigl[\dot{\varphi}_{i}\,\delta\varphi_{i}\Bigr]_{t_{1}}^{t_{2}}
   \mathrm{d}V_{0}
   - \int_{t_{1}}^{t_{2}}\!\int_{\mathcal{B}_{0}}
   \rho_{0}\,\ddot{\varphi}_{i}\,\delta\varphi_{i}\,
   \mathrm{d}V_{0}\,\mathrm{d}t.
\end{align}

\textbf{Step 4:} Apply endpoint conditions
$\delta\varphi_{i}(t_{1})=\delta\varphi_{i}(t_{2})=0$.

The boundary term vanishes:
$\bigl[\dot{\varphi}_{i}\,\delta\varphi_{i}\bigr]_{t_{1}}^{t_{2}}=0$.

Since $\ddot{\varphi}_{i}=\ddot{u}_{i}$ and $\delta\varphi_{i}=\delta u_{i}$
(because $\bX$ is time-independent):

\begin{equation}
\delta\int_{t_{1}}^{t_{2}}\mathcal{T}\,\mathrm{d}t
= -\int_{t_{1}}^{t_{2}}\!\int_{\mathcal{B}_{0}}
  \rho_{0}\,\ddot{u}_{i}\,\delta u_{i}\,
  \mathrm{d}V_{0}\,\mathrm{d}t
= -\int_{t_{1}}^{t_{2}}\delta\mathcal{W}_{\mathrm{iner}}\,\mathrm{d}t.
\end{equation}

$\checkmark$

\subsection{Instantaneous variational statement}

Substituting into Hamilton's principle:

\begin{equation}
-\int_{t_{1}}^{t_{2}}\delta\mathcal{W}_{\mathrm{iner}}\,\mathrm{d}t
- \int_{t_{1}}^{t_{2}}\delta\Pi_{\mathrm{int}}\,\mathrm{d}t
+ \int_{t_{1}}^{t_{2}}\delta\mathcal{W}_{\mathrm{ext}}\,\mathrm{d}t
= 0.
\end{equation}

Since this holds for arbitrary $[t_{1},t_{2}]$:

\begin{equation}
\boxed{
\delta\Pi_{\mathrm{int}}
= \delta\mathcal{W}_{\mathrm{ext}}
- \delta\mathcal{W}_{\mathrm{iner}}.
}
\end{equation}

% =====================================================================
\section{Variation of internal energy}
% =====================================================================

\subsection{Chain rule for $\delta\psi$}

Since $\psi=\psi(\mathbf{E},\bm{\mathcal{G}})$ and
$\mathcal{B}_{0}$ is fixed (TL description):

\begin{equation}
\delta\Pi_{\mathrm{int}}
= \int_{\mathcal{B}_{0}}\delta\psi\,\mathrm{d}V_{0}.
\end{equation}

The chain rule for a function of two tensor arguments:

\begin{align}
\delta\psi
&= \frac{\partial\psi}{\partial E_{AB}}\,\delta E_{AB}
 + \frac{\partial\psi}{\partial\mathcal{G}_{ABC}}\,\delta\mathcal{G}_{ABC}
\nonumber\\
&= S_{AB}\,\delta E_{AB}
 + \Sigma_{ABC}\,\delta\mathcal{G}_{ABC}.
\end{align}

Relabeling dummy indices $A\to I$, $B\to J$, $C\to K$:

\begin{equation}
\delta\psi = S_{IJ}\,\delta E_{IJ}
+ \Sigma_{IJK}\,\delta\mathcal{G}_{IJK}.
\end{equation}

\textbf{Verification that each contraction is a scalar:}

$S_{IJ}\delta E_{IJ}$: $S_{IJ}$ is 2nd-order, $\delta E_{IJ}$ is
2nd-order, summing over $I,J$ gives a scalar. $\checkmark$

$\Sigma_{IJK}\delta\mathcal{G}_{IJK}$: $\Sigma_{IJK}$ is 3rd-order,
$\delta\mathcal{G}_{IJK}$ is 3rd-order, summing over $I,J,K$ gives
a scalar. $\checkmark$

Therefore:

\begin{equation}
\boxed{
\delta\Pi_{\mathrm{int}}
= \int_{\mathcal{B}_{0}}
  \bigl[\,S_{IJ}\,\delta E_{IJ}
        + \Sigma_{IJK}\,\delta\mathcal{G}_{IJK}\bigr]\,
  \mathrm{d}V_{0}.
}
\end{equation}

\subsection{Variation chain: $\delta\bu \to \delta\bF \to \delta\bE \to \delta\bm{\mathcal{G}}$}

\textbf{Step 1:} $\delta\bF$.

\begin{equation}
\delta F_{iJ}
= \delta(\delta_{iJ}+u_{i,J})
= (\delta u_{i})_{,J}
= \delta u_{i,J}.
\end{equation}

($\delta_{iJ}$ is constant, $\delta$ commutes with $\partial/\partial X_J$
since $\bX$ is the independent variable.)

\textbf{Step 2:} $\delta\bE$.

\begin{align}
\delta E_{IJ}
&= \delta\!\left[\tfrac12(F_{kI}F_{kJ}-\delta_{IJ})\right]
\nonumber\\
&= \tfrac12\bigl(\delta F_{kI}\cdot F_{kJ}
              + F_{kI}\cdot\delta F_{kJ}\bigr)
\nonumber\\
&= \tfrac12\bigl(\delta u_{k,I}\,F_{kJ}
              + F_{kI}\,\delta u_{k,J}\bigr).
\end{align}

Relabeling $k\to i$:

\begin{equation}
\delta E_{IJ}
= \tfrac12\bigl(F_{iI}\,\delta u_{i,J}
              + F_{iJ}\,\delta u_{i,I}\bigr).
\end{equation}

\textbf{Symmetry check:} $\delta E_{IJ}=\delta E_{JI}$ because
swapping $I\leftrightarrow J$ exchanges the two terms. $\checkmark$

\textbf{Step 3:} $\delta\bm{\mathcal{G}}$.

\begin{equation}
\delta\mathcal{G}_{IJK}
= \delta(E_{IJ,K})
= (\delta E_{IJ})_{,K}
= \delta E_{IJ,K}.
\end{equation}

($\delta$ commutes with $\partial/\partial X_K$ since $\bX$ is
independent of the variation.)

Substituting the expression for $\delta E_{IJ}$:

\begin{align}
\delta\mathcal{G}_{IJK}
&= \tfrac12\bigl(F_{iI}\,\delta u_{i,J}
              + F_{iJ}\,\delta u_{i,I}\bigr)_{,K}
\nonumber\\
&= \tfrac12\bigl(F_{iI,K}\,\delta u_{i,J}
              + F_{iI}\,\delta u_{i,JK}
              + F_{iJ,K}\,\delta u_{i,I}
              + F_{iJ}\,\delta u_{i,IK}\bigr).
\end{align}

This shows that $\delta\mathcal{G}_{IJK}$ involves second derivatives
of $\delta\bu$, which is the origin of the $C^{1}$ requirement.

% =====================================================================
\section{Assembly of the weak form}
% =====================================================================

\subsection{Full substitution}

Substituting all components into $\delta\Pi_{\mathrm{int}}
= \delta\mathcal{W}_{\mathrm{ext}} - \delta\mathcal{W}_{\mathrm{iner}}$:

\begin{align}
\int_{\mathcal{B}_{0}}
&\bigl[\,S_{IJ}\,\delta E_{IJ}
      + \Sigma_{IJK}\,\delta\mathcal{G}_{IJK}\bigr]\,
  \mathrm{d}V_{0}
\nonumber\\
&= \int_{\mathcal{B}_{0}} f_{0i}\,\delta u_{i}\,\mathrm{d}V_{0}
 + \int_{\partial\mathcal{B}_{0t}} \bar{T}_{0i}\,\delta u_{i}\,
   \mathrm{d}A_{0}
 - \int_{\mathcal{B}_{0}} \rho_{0}\,\ddot{u}_{i}\,\delta u_{i}\,
   \mathrm{d}V_{0}.
\label{fd:weak_full}
\end{align}

\subsection{Simplification of the classical branch}
\label{fd:classical_simplify}

\textbf{Goal:} Show that $S_{IJ}\,\delta E_{IJ} = P_{iJ}\,\delta u_{i,J}$.

\begin{align}
S_{IJ}\,\delta E_{IJ}
&= S_{IJ}\cdot\tfrac12\bigl(F_{iI}\,\delta u_{i,J}
                        + F_{iJ}\,\delta u_{i,I}\bigr)
\nonumber\\
&= \tfrac12\,S_{IJ}\,F_{iI}\,\delta u_{i,J}
 + \tfrac12\,S_{IJ}\,F_{iJ}\,\delta u_{i,I}.
\end{align}

\textbf{Index swap in the second term:}
Relabel $I\leftrightarrow J$ (both are dummy indices summed from 1 to 3):

\begin{align}
\tfrac12\,S_{IJ}\,F_{iJ}\,\delta u_{i,I}
&\xrightarrow{I\to J',\;J\to I'}
\tfrac12\,S_{J'I'}\,F_{iI'}\,\delta u_{i,J'}
\nonumber\\
&= \tfrac12\,S_{I'J'}\,F_{iI'}\,\delta u_{i,J'}
\quad\text{(using }S_{J'I'}=S_{I'J'}\text{)}
\nonumber\\
&= \tfrac12\,S_{IJ}\,F_{iI}\,\delta u_{i,J}
\quad\text{(renaming }I'\to I,\;J'\to J\text{)}.
\end{align}

Therefore both terms are identical:

\begin{align}
S_{IJ}\,\delta E_{IJ}
&= \tfrac12\,S_{IJ}\,F_{iI}\,\delta u_{i,J}
 + \tfrac12\,S_{IJ}\,F_{iI}\,\delta u_{i,J}
\nonumber\\
&= S_{IJ}\,F_{iI}\,\delta u_{i,J}
\nonumber\\
&= P_{iJ}\,\delta u_{i,J}.
\end{align}

$\checkmark$

\subsection{Simplification of the gradient branch}

\textbf{Step 1:} Write out $\Sigma_{IJK}\,\delta\mathcal{G}_{IJK}$.

\begin{equation}
\Sigma_{IJK}\,\delta\mathcal{G}_{IJK}
= \Sigma_{IJK}\,\delta E_{IJ,K}.
\end{equation}

\textbf{Step 2:} Substitute $\delta E_{IJ}$:

\begin{align}
\Sigma_{IJK}\,\delta E_{IJ,K}
&= \Sigma_{IJK}\cdot
   \tfrac12\bigl(F_{iI}\,\delta u_{i,J}
              + F_{iJ}\,\delta u_{i,I}\bigr)_{,K}.
\end{align}

However, it is more useful to keep $\delta E_{IJ,K}$ in the form
$(\delta E_{IJ})_{,K}$ for the first IBP, and only substitute
after IBP.

The gradient branch in integral form is:

\begin{equation}
\int_{\mathcal{B}_{0}}\Sigma_{IJK}\,(\delta E_{IJ})_{,K}\,
  \mathrm{d}V_{0}.
\end{equation}

% =====================================================================
\section{Integration by parts -- complete derivation}
% =====================================================================

\subsection{Classical branch IBP}
\label{fd:ibp_P_summary}

\textbf{Term:} $\displaystyle\int_{\mathcal{B}_{0}}
P_{iJ}\,\delta u_{i,J}\,\mathrm{d}V_{0}$.

\textbf{Apply product rule:}

\begin{equation}
P_{iJ}\,\delta u_{i,J}
= (P_{iJ}\,\delta u_{i})_{,J}
- P_{iJ,J}\,\delta u_{i}.
\end{equation}

Proof: $(P_{iJ}\delta u_{i})_{,J}
= P_{iJ,J}\delta u_{i} + P_{iJ}\delta u_{i,J}$ (product rule). $\checkmark$

\textbf{Apply divergence theorem:}

\begin{equation}
\int_{\mathcal{B}_{0}}(P_{iJ}\,\delta u_{i})_{,J}\,\mathrm{d}V_{0}
= \int_{\partial\mathcal{B}_{0}}P_{iJ}\,\delta u_{i}\,N_{J}\,
  \mathrm{d}A_{0},
\end{equation}

where $N_{J}$ is the outward unit normal on $\partial\mathcal{B}_{0}$.

\textbf{Result:}

\begin{equation}
\int_{\mathcal{B}_{0}}P_{iJ}\,\delta u_{i,J}\,\mathrm{d}V_{0}
= \int_{\partial\mathcal{B}_{0}}P_{iJ}\,N_{J}\,\delta u_{i}\,
  \mathrm{d}A_{0}
- \int_{\mathcal{B}_{0}}P_{iJ,J}\,\delta u_{i}\,\mathrm{d}V_{0}.
\end{equation}

$\checkmark$

\subsection{Gradient branch -- First IBP (over $X_K$)}

\textbf{Term:} $\displaystyle\int_{\mathcal{B}_{0}}
\Sigma_{IJK}\,\delta E_{IJ,K}\,\mathrm{d}V_{0}$.

\textbf{Apply product rule:}

\begin{equation}
\Sigma_{IJK}\,\delta E_{IJ,K}
= (\Sigma_{IJK}\,\delta E_{IJ})_{,K}
- \Sigma_{IJK,K}\,\delta E_{IJ}.
\end{equation}

\textbf{Apply divergence theorem:}

\begin{equation}
\int_{\mathcal{B}_{0}}(\Sigma_{IJK}\,\delta E_{IJ})_{,K}\,
  \mathrm{d}V_{0}
= \int_{\partial\mathcal{B}_{0}}\Sigma_{IJK}\,N_{K}\,
  \delta E_{IJ}\,\mathrm{d}A_{0}.
\end{equation}

\textbf{Result:}

\begin{equation}
\int_{\mathcal{B}_{0}}\Sigma_{IJK}\,\delta E_{IJ,K}\,
  \mathrm{d}V_{0}
= \int_{\partial\mathcal{B}_{0}}\Sigma_{IJK}\,N_{K}\,
  \delta E_{IJ}\,\mathrm{d}A_{0}
- \int_{\mathcal{B}_{0}}\Sigma_{IJK,K}\,\delta E_{IJ}\,
  \mathrm{d}V_{0}.
\label{fd:ibp1}
\end{equation}

$\checkmark$

\subsection{Gradient branch -- Index manipulation of domain term}

\textbf{Substitute $\delta E_{IJ}$ into the remaining domain integral:}

\begin{align}
-\int_{\mathcal{B}_{0}}\Sigma_{IJK,K}\,\delta E_{IJ}\,
  \mathrm{d}V_{0}
&= -\int_{\mathcal{B}_{0}}\Sigma_{IJK,K}\cdot
   \tfrac12\bigl(F_{iI}\,\delta u_{i,J}
              + F_{iJ}\,\delta u_{i,I}\bigr)\,
  \mathrm{d}V_{0}
\nonumber\\
&= -\tfrac12\int_{\mathcal{B}_{0}}
   \Sigma_{IJK,K}\,F_{iI}\,\delta u_{i,J}\,
   \mathrm{d}V_{0}
 - \tfrac12\int_{\mathcal{B}_{0}}
   \Sigma_{IJK,K}\,F_{iJ}\,\delta u_{i,I}\,
   \mathrm{d}V_{0}.
\label{fd:two_terms}
\end{align}

\textbf{Index swap in the second integral:}

In the second integral, relabel $I\leftrightarrow J$:

\begin{align}
-\tfrac12\int_{\mathcal{B}_{0}}
   \Sigma_{IJK,K}\,F_{iJ}\,\delta u_{i,I}\,
   \mathrm{d}V_{0}
&\xrightarrow{I\to J',\;J\to I'}
-\tfrac12\int_{\mathcal{B}_{0}}
   \Sigma_{J'I'K,K}\,F_{iI'}\,\delta u_{i,J'}\,
   \mathrm{d}V_{0}.
\end{align}

Using $\Sigma_{J'I'K}=\Sigma_{I'J'K}$ (symmetry of $\boldsymbol{\Sigma}$):

\begin{equation}
= -\tfrac12\int_{\mathcal{B}_{0}}
   \Sigma_{I'J'K,K}\,F_{iI'}\,\delta u_{i,J'}\,
   \mathrm{d}V_{0}.
\end{equation}

Renaming dummy indices $I'\to I$, $J'\to J$:

\begin{equation}
= -\tfrac12\int_{\mathcal{B}_{0}}
   \Sigma_{IJK,K}\,F_{iI}\,\delta u_{i,J}\,
   \mathrm{d}V_{0}.
\end{equation}

This is \emph{identical} to the first integral in \eqref{fd:two_terms}.

\textbf{Combining:}

\begin{equation}
-\int_{\mathcal{B}_{0}}\Sigma_{IJK,K}\,\delta E_{IJ}\,
  \mathrm{d}V_{0}
= -\int_{\mathcal{B}_{0}}\Sigma_{IJK,K}\,F_{iI}\,
  \delta u_{i,J}\,\mathrm{d}V_{0}.
\label{fd:combined_domain}
\end{equation}

$\checkmark$

\subsection{Gradient branch -- Second IBP (over $X_J$)}

\textbf{Term:}
$\displaystyle-\int_{\mathcal{B}_{0}}\Sigma_{IJK,K}\,F_{iI}\,
  \delta u_{i,J}\,\mathrm{d}V_{0}$.

Define $A_{iJ}:=\Sigma_{IJK,K}\,F_{iI}$. Then the integral is

**Remark:** $A_{iJ}$ is introduced solely as an auxiliary tensor to simplify the integration by parts operation. It is not an independent constitutive quantity and does not introduce any new material parameters. The tensor $A_{iJ}$ is merely a notational convenience representing the contraction of the double-stress divergence $\Sigma_{IJK,K}$ with the deformation gradient $F_{iI}$. All material information remains contained in $\Sigma_{IJK}$ and $F_{iI}$.

\begin{equation}
-\int_{\mathcal{B}_{0}}A_{iJ}\,\delta u_{i,J}\,\mathrm{d}V_{0}.
\end{equation}

\textbf{Apply product rule:}

\begin{equation}
A_{iJ}\,\delta u_{i,J}
= (A_{iJ}\,\delta u_{i})_{,J}
- A_{iJ,J}\,\delta u_{i}.
\end{equation}

\textbf{Apply divergence theorem:}

\begin{equation}
\int_{\mathcal{B}_{0}}(A_{iJ}\,\delta u_{i})_{,J}\,\mathrm{d}V_{0}
= \int_{\partial\mathcal{B}_{0}}A_{iJ}\,N_{J}\,\delta u_{i}\,
  \mathrm{d}A_{0}.
\end{equation}

\textbf{Result:}

\begin{align}
-\int_{\mathcal{B}_{0}}A_{iJ}\,\delta u_{i,J}\,\mathrm{d}V_{0}
&= -\int_{\partial\mathcal{B}_{0}}A_{iJ}\,N_{J}\,\delta u_{i}\,
   \mathrm{d}A_{0}
 + \int_{\mathcal{B}_{0}}A_{iJ,J}\,\delta u_{i}\,\mathrm{d}V_{0}.
\end{align}

Substituting back $A_{iJ}=\Sigma_{IJK,K}F_{iI}$:

\begin{align}
-\int_{\mathcal{B}_{0}}\Sigma_{IJK,K}\,F_{iI}\,
  \delta u_{i,J}\,\mathrm{d}V_{0}
&= -\int_{\partial\mathcal{B}_{0}}\bigl(\Sigma_{IJK,K}\,F_{iI}\bigr)\,
   N_{J}\,\delta u_{i}\,\mathrm{d}A_{0}
\nonumber\\
&\quad + \int_{\mathcal{B}_{0}}\bigl(\Sigma_{IJK,K}\,F_{iI}\bigr)_{,J}\,
   \delta u_{i}\,\mathrm{d}V_{0}.
\label{fd:ibp2}
\end{align}

$\checkmark$

\textbf{Expansion of the domain term by product rule:}

\begin{align}
\bigl(\Sigma_{IJK,K}\,F_{iI}\bigr)_{,J}
&= \Sigma_{IJK,KJ}\,F_{iI}
 + \Sigma_{IJK,K}\,F_{iI,J}.
\end{align}

This is a fourth-order differential operator on $\delta\bu$
(since $\Sigma_{IJK,KJ}$ involves $\mathcal{G}_{IJK,JK}$, which
contains fourth derivatives of $\bu$).

\textbf{Note on the hyperstress divergence:}

\begin{equation}
\mathcal{H}_{iJK,K}
= (\Sigma_{IJK}F_{iI})_{,K}
= \Sigma_{IJK,K}F_{iI} + \Sigma_{IJK}F_{iI,K}.
\end{equation}

At finite strain, $\Sigma_{IJK}F_{iI,K}\neq 0$ in general, so
$(\Sigma_{IJK,K}F_{iI}) \neq \mathcal{H}_{iJK,K}$.
This distinction is important for the boundary terms.

% =====================================================================
\section{Assembly of the strong form and boundary conditions}
% =====================================================================

\subsection{Domain (strong form) assembly}

Collecting all volume contributions from the classical branch
(Section~\ref{fd:ibp_P_summary}), the gradient branch
(Eqs.~\eqref{fd:combined_domain} and \eqref{fd:ibp2}), and the
external/inertia terms:

\begin{align}
\int_{\mathcal{B}_{0}}\Bigl[
&-P_{iJ,J}
+ \bigl(\Sigma_{IJK,K}\,F_{iI}\bigr)_{,J}
+ f_{0i}
- \rho_{0}\,\ddot{u}_{i}
\Bigr]\delta u_{i}\,\mathrm{d}V_{0}
\nonumber\\
&\quad + \text{boundary terms} = 0.
\end{align}

Since $\delta\bu$ is arbitrary in $\mathcal{B}_{0}$:

\begin{equation}
\boxed{
-P_{iJ,J} + \bigl(\Sigma_{IJK,K}\,F_{iI}\bigr)_{,J}
- f_{0i} + \rho_{0}\,\ddot{u}_{i} = 0
\qquad\text{in }\mathcal{B}_{0}.
}
\end{equation}

Equivalently, moving source terms to the RHS:

\begin{equation}
P_{iJ,J} - \bigl(\Sigma_{IJK,K}\,F_{iI}\bigr)_{,J} + f_{0i}
= \rho_{0}\,\ddot{u}_{i}.
\end{equation}

Defining $\mathrm{Div}_{0}\mathbf{T}
:= P_{iJ,J}-(\Sigma_{IJK,K}F_{iI})_{,J}$:

\begin{equation}
\mathrm{Div}_{0}\mathbf{T} + \mathbf{f}_{0}
= \rho_{0}\,\ddot{\mathbf{u}}.
\end{equation}

$\checkmark$

\subsection{Boundary terms assembly}

The boundary contributions are:

\begin{align}
\text{From classical IBP:}&\quad
+\int_{\partial\mathcal{B}_{0}}P_{iJ}\,N_{J}\,\delta u_{i}\,
  \mathrm{d}A_{0},
\\
\text{From gradient 1st IBP:}&\quad
+\int_{\partial\mathcal{B}_{0}}\Sigma_{IJK}\,N_{K}\,
  \delta E_{IJ}\,\mathrm{d}A_{0},
\\
\text{From gradient 2nd IBP:}&\quad
-\int_{\partial\mathcal{B}_{0}}\bigl(\Sigma_{IJK,K}\,F_{iI}\bigr)\,
  N_{J}\,\delta u_{i}\,\mathrm{d}A_{0},
\\
\text{From external work:}&\quad
+\int_{\partial\mathcal{B}_{0t}}\bar{T}_{0i}\,\delta u_{i}\,
  \mathrm{d}A_{0}.
\end{align}

\textbf{Combining terms proportional to $\delta u_{i}$:}

\begin{equation}
\int_{\partial\mathcal{B}_{0}}
\bigl(P_{iJ} - \Sigma_{IJK,K}\,F_{iI}\bigr)\,N_{J}\,
\delta u_{i}\,\mathrm{d}A_{0}.
\end{equation}

\textbf{Higher-order boundary term:}

\begin{equation}
\int_{\partial\mathcal{B}_{0}}\Sigma_{IJK}\,N_{K}\,
  \delta E_{IJ}\,\mathrm{d}A_{0}.
\end{equation}

Substituting $\delta E_{IJ} = \tfrac12(F_{lI}\delta u_{l,J}+F_{lJ}\delta u_{l,I})$:

\begin{align}
\Sigma_{IJK}\,N_{K}\,\delta E_{IJ}
&= \Sigma_{IJK}\,N_{K}\cdot\tfrac12\bigl(F_{lI}\delta u_{l,J}
                                      + F_{lJ}\delta u_{l,I}\bigr)
\nonumber\\
&= \tfrac12\Sigma_{IJK}\,F_{lI}\,N_{K}\,\delta u_{l,J}
 + \tfrac12\Sigma_{IJK}\,F_{lJ}\,N_{K}\,\delta u_{l,I}.
\end{align}

In the second term, swap $I\leftrightarrow J$:

\begin{align}
\tfrac12\Sigma_{IJK}\,F_{lJ}\,N_{K}\,\delta u_{l,I}
&\xrightarrow{I\to J',\;J\to I'}
\tfrac12\Sigma_{J'I'K}\,F_{lI'}\,N_{K}\,\delta u_{l,J'}
\nonumber\\
&= \tfrac12\Sigma_{I'J'K}\,F_{lI'}\,N_{K}\,\delta u_{l,J'}
\quad(\text{since }\Sigma_{J'I'K}=\Sigma_{I'J'K})
\nonumber\\
&= \tfrac12\Sigma_{IJK}\,F_{lI}\,N_{K}\,\delta u_{l,J}.
\end{align}

Combining:

\begin{equation}
\Sigma_{IJK}\,N_{K}\,\delta E_{IJ}
= \Sigma_{IJK}\,F_{lI}\,N_{K}\,\delta u_{l,J}
= \mathcal{H}_{lJK}\,N_{K}\,\delta u_{l,J}.
\end{equation}

\textbf{Verification:} $\mathcal{H}_{lJK}N_K\delta u_{l,J}$
$= \Sigma_{IJK}F_{lI}N_K\delta u_{l,J}$.
All indices ($l$, $J$, $K$) are contracted. Result is a scalar. $\checkmark$

\subsection{Extraction of boundary conditions}

The boundary condition derivation follows Toupin \cite{Toupin1964} and Germain \cite{Germain1973}. The $C^1$ requirement for displacement-based strain-gradient elements is discussed by Amanatidou \& Aravas \cite{AmanatidouAravas2002} and Engel et al.\ \cite{EngelGarikipati2002}.

The total boundary virtual work is:

\begin{align}
\delta\mathcal{W}_{\mathrm{bnd}}
&= \int_{\partial\mathcal{B}_{0}}
   \bigl(P_{iJ}-\Sigma_{IJK,K}F_{iI}\bigr)N_J\,\delta u_i\,
   \mathrm{d}A_0
 + \int_{\partial\mathcal{B}_{0}}
   \mathcal{H}_{lJK}\,N_K\,\delta u_{l,J}\,
   \mathrm{d}A_0.
\end{align}

The boundary $\partial\mathcal{B}_{0}$ is partitioned into
$\partial\mathcal{B}_{0u}\cup\partial\mathcal{B}_{0t}$.

\textbf{On $\partial\mathcal{B}_{0u}$:}
$\delta\bu=\mathbf{0}$ and $\partial(\delta\bu)/\partial N=\mathbf{0}$.
All boundary terms vanish.

\textbf{On $\partial\mathcal{B}_{0t}$:}
$\delta\bu$ and $\partial(\delta\bu)/\partial N$ are free.

Decomposing $\delta u_{l,J}$ into normal and tangential parts:

\begin{equation}
\delta u_{l,J} = N_J\,(\delta u_{l,M}N_M)
+ (\delta_{JM}-N_JN_M)\,\delta u_{l,M}
= N_J\,\delta u_{l,N} + \delta u_{l,J}^{\mathrm{tan}}.
\end{equation}

The higher-order boundary term becomes:

\begin{equation}
\mathcal{H}_{lJK}N_K\delta u_{l,J}
= \mathcal{H}_{lJK}N_KN_J\,\delta u_{l,N}
+ \mathcal{H}_{lJK}N_K\,\delta u_{l,J}^{\mathrm{tan}}.
\end{equation}

**Rigorous justification of independent boundary variables** (following Toupin \cite{Toupin1964} and Germain \cite{Germain1973}):

For a displacement-based $C^1$ formulation, the virtual displacement field $\delta\mathbf{u}$ and its gradient $\nabla_0\delta\mathbf{u}$ must be specified on $\partial\mathcal{B}_{0u}$. The gradient decomposition on the boundary yields:
\begin{equation}
\delta u_{l,J} = N_J\,\delta u_{l,N} + \delta u_{l,J}^{\mathrm{tan}},
\end{equation}
where $\delta u_{l,N} = N_M\,\delta u_{l,M}$ is the normal derivative variation and $\delta u_{l,J}^{\mathrm{tan}} = (\delta_{JM}-N_JN_M)\,\delta u_{l,M}$ contains tangential derivatives.

**Key argument:** Since $\delta\mathbf{u} = \mathbf{0}$ on $\partial\mathcal{B}_{0u}$, we have $\delta u_{l,M} = 0$ for all $M$. Therefore:
\begin{equation}
\delta u_{l,J}^{\mathrm{tan}} = (\delta_{JM}-N_JN_M)\cdot 0 = 0.
\end{equation}

**Functional-analytic justification (H² Sobolev regularity and trace theory):**

The displacement-based strain-gradient formulation requires test and trial functions in $H^2(\mathcal{B}_0)$, since the weak form involves second derivatives of the displacement field. By the Sobolev trace theorem for $H^2$ functions, both the function values $\delta\mathbf{u}|_{\partial\mathcal{B}_{0}}$ and the full gradient $\nabla\delta\mathbf{u}|_{\partial\mathcal{B}_{0}}$ are well-defined trace operators into $H^{3/2}(\partial\mathcal{B}_{0})$ and $H^{1/2}(\partial\mathcal{B}_{0})$, respectively. However, the gradient trace decomposes into normal and tangential components: $\nabla\delta\mathbf{u} = \mathbf{N}\otimes\mathbf{N}\nabla\delta\mathbf{u} + (\mathbf{I}-\mathbf{N}\otimes\mathbf{N})\nabla\delta\mathbf{u}$.

For a $C^1$ finite element interpolation (e.g., Bogner--Fox--Schmit bicubic Hermite elements), the nodal degrees of freedom are precisely the displacement values and the full gradient at element vertices. On an essential boundary $\partial\mathcal{B}_{0u}$ where $\delta\mathbf{u} = \mathbf{0}$ is prescribed, the $C^1$ interpolant automatically enforces $\nabla\delta\mathbf{u} = \mathbf{0}$ at all boundary nodes. Since the tangential derivatives $\delta u_{l,J}^{\mathrm{tan}}$ are computed from differences of displacement values along the boundary, and these values are identically zero on $\partial\mathcal{B}_{0u}$, the tangential derivatives are not independent degrees of freedom—they are completely determined by the boundary condition $\delta\mathbf{u} = \mathbf{0}$.

In contrast, the normal derivative $\partial(\delta u_l)/\partial N$ involves differences in the direction transverse to the boundary, where the displacement field is not constrained by essential boundary conditions on $\partial\mathcal{B}_{0t}$. Therefore, the normal derivative remains an independent degree of freedom that can be freely varied on $\partial\mathcal{B}_{0t}$. This distinction follows Toupin's principle of virtual power \cite{Toupin1964} and Germain's method of virtual work \cite{Germain1973}, which identify the normal derivative as the independent kinematic quantity conjugate to the double traction in strain-gradient theories.

The tangential derivatives vanish identically on $\partial\mathcal{B}_{0u}$ and cannot contribute to boundary virtual work. Only the normal derivative $\delta u_{l,N}$ remains as an independent virtual quantity on $\partial\mathcal{B}_{0t}$.

Following Toupin's principle of virtual power \cite{Toupin1964} and Germain's method of virtual work \cite{Germain1973}, the higher-order boundary term $\mathcal{H}_{lJK}N_K\,\delta u_{l,J}$ on $\partial\mathcal{B}_{0t}$ decomposes as:
\begin{equation}
\mathcal{H}_{lJK}N_K\,\delta u_{l,J} = \underbrace{\mathcal{H}_{lJK}N_KN_J}_{\mathcal{D}_l}\,\delta u_{l,N} + \underbrace{\mathcal{H}_{lJK}N_K\,\delta u_{l,J}^{\mathrm{tan}}}_{=\,0\;\text{on}\;\partial\mathcal{B}_{0u}}.
\end{equation}

Since $\delta u_{l,N}$ is the sole independent variation in the normal direction, its conjugate $\mathcal{D}_l = \mathcal{H}_{lJK}N_KN_J$ must be prescribed as the double-traction boundary condition on $\partial\mathcal{B}_{0t}$:
\begin{equation}
\mathcal{D}_i = \mathcal{H}_{iJK}N_KN_J = \bar{d}_i \quad \text{on } \partial\mathcal{B}_{0t}.
\end{equation}

The standard treatment \cite{Toupin1964,Germain1973} identifies:

\textbf{Natural BCs on $\partial\mathcal{B}_{0t}$:}

\begin{align}
\bigl(P_{iJ}-\Sigma_{IJK,K}F_{iI}\bigr)N_J &= \bar{T}_{0i}
\quad\text{(effective single traction)},
\\
\mathcal{D}_{i} := \mathcal{H}_{iJK}N_KN_J &= \bar{d}_{i}
\quad\text{(double traction)}.
\end{align}

\textbf{Essential BCs on $\partial\mathcal{B}_{0u}$:}

\begin{equation}
u_{i} = \bar{u}_{i},
\qquad
\frac{\partial u_{i}}{\partial N} = \bar{g}_{i}.
\end{equation}

% =====================================================================
\section{The weak form for FEM discretization}
% =====================================================================

\subsection{Unreduced weak form (working equation)}

The weak form is \emph{not} the strong form. It is the integral
statement \emph{before} IBP, which is what the finite element method
actually discretizes:

\begin{equation}
\boxed{
\int_{\mathcal{B}_{0}}
\bigl[\,S_{IJ}\,\delta E_{IJ}
      + \Sigma_{IJK}\,\delta\mathcal{G}_{IJK}\bigr]\,
\mathrm{d}V_{0}
= \int_{\mathcal{B}_{0}} f_{0i}\,\delta u_{i}\,\mathrm{d}V_{0}
+ \int_{\partial\mathcal{B}_{0t}} \bar{T}_{0i}\,\delta u_{i}\,
  \mathrm{d}A_{0}
- \int_{\mathcal{B}_{0}} \rho_{0}\,\ddot{u}_{i}\,\delta u_{i}\,
  \mathrm{d}V_{0}.
}
\end{equation}

Using the identity $S_{IJ}\delta E_{IJ}=P_{iJ}\delta u_{i,J}$
(derived in Section~\ref{fd:classical_simplify}):

\begin{equation}
\int_{\mathcal{B}_{0}}
\bigl[\,P_{iJ}\,\delta u_{i,J}
      + \Sigma_{IJK}\,(\delta E_{IJ})_{,K}\bigr]\,
\mathrm{d}V_{0}
= \int_{\mathcal{B}_{0}} (f_{0i}-\rho_{0}\ddot{u}_{i})\,
  \delta u_{i}\,\mathrm{d}V_{0}
+ \int_{\partial\mathcal{B}_{0t}} \bar{T}_{0i}\,\delta u_{i}\,
  \mathrm{d}A_{0}.
\end{equation}

This requires $\bu,\delta\bu\in H^{2}(\mathcal{B}_{0})$ (i.e.,
$C^{1}$ continuity), which the BFS element provides.

\subsection{Static (quasi-static) weak form}

Setting $\ddot{\mathbf{u}}=\mathbf{0}$:

\begin{equation}
\int_{\mathcal{B}_{0}}
\bigl[\,P_{iJ}\,\delta u_{i,J}
      + \Sigma_{IJK}\,(\delta E_{IJ})_{,K}\bigr]\,
\mathrm{d}V_{0}
= \int_{\mathcal{B}_{0}} f_{0i}\,\delta u_{i}\,\mathrm{d}V_{0}
+ \int_{\partial\mathcal{B}_{0t}} \bar{T}_{0i}\,\delta u_{i}\,
  \mathrm{d}A_{0}.
\end{equation}

This is the form used by the quasi-static 2022 benchmark.

% =====================================================================
\section{Small-strain limit and 2022 benchmark recovery}
% =====================================================================

The recovery of classical linear elasticity \cite{Holzapfel2000} and Mindlin Form-II strain-gradient theory \cite{Mindlin1964,MindlinEshel1968} is verified below.

\subsection{Limit of each quantity}

As $\|\nabla_{0}\mathbf{u}\|\to 0$:

\textbf{Deformation gradient:}
$F_{iI} = \delta_{iI}+u_{i,I} \to \delta_{iI}$.

\textbf{Strain:}
$E_{IJ} = \tfrac12(u_{I,J}+u_{J,I}+u_{K,I}u_{K,J})
\to \tfrac12(u_{I,J}+u_{J,I}) = \varepsilon_{IJ}$.

\textbf{Strain gradient:}
$\mathcal{G}_{IJK} = E_{IJ,K} \to \varepsilon_{IJ,K}$.

\textbf{Stresses:}
$S_{IJ} = \lambda(\tr\mathbf{E})\delta_{IJ}+2\mu E_{IJ}
\to \lambda(\tr\boldsymbol{\varepsilon})\delta_{ij}+2\mu\varepsilon_{ij}
= \sigma_{ij}$.

$\Sigma_{IJK} = \mathbb{D}_{IJKPQR}\mathcal{G}_{PQR}
\to \mathbb{D}_{ijkpqr}\varepsilon_{pq,r} = h_{ijk}$.

\textbf{First Piola--Kirchhoff stress:}
$P_{iJ} = S_{IJ}F_{iI} \to \sigma_{ij}\delta_{jJ} = \sigma_{iJ}$.

\textbf{Density:}
$\rho = \rho_{0}/J \to \rho_{0}$. This follows from $J \to 1$ under the infinitesimal deformation limit.

\textbf{Body force:}
$\mathbf{f}_{0} \to \mathbf{b}$ (per unit current volume).

\textbf{Traction:}
$\bar{\mathbf{T}}_{0} \to \bar{\mathbf{t}}$ (Cauchy traction).

\subsection{Weak form recovery}

\textbf{Classical branch:}

\begin{equation}
P_{iJ}\,\delta u_{i,J} \to \sigma_{iJ}\,\delta u_{i,J}
= \sigma_{ij}\,\delta u_{i,j} = \sigma_{ij}\,\delta\varepsilon_{ij}.
\end{equation}

(Using $\delta\varepsilon_{ij}=\tfrac12(\delta u_{i,j}+\delta u_{j,i})$
and $\sigma_{ij}=\sigma_{ji}$:
$\sigma_{ij}\delta u_{i,j}
= \sigma_{ij}\cdot\tfrac12(\delta u_{i,j}+\delta u_{j,i})
= \sigma_{ij}\delta\varepsilon_{ij}$.)

\textbf{Gradient branch:}

\begin{equation}
\Sigma_{IJK}\,(\delta E_{IJ})_{,K}
\to h_{ijk}\,(\delta\varepsilon_{ij})_{,k}
= h_{ijk}\,\delta\varepsilon_{ij,k}.
\end{equation}

\textbf{Full weak form:}

\begin{equation}
\boxed{
\int_{\Omega}
\bigl[\,\sigma_{ij}\,\delta\varepsilon_{ij}
+ h_{ijk}\,\delta\varepsilon_{ij,k}\bigr]\,
\mathrm{d}v
= \int_{\Omega} b_{i}\,\delta u_{i}\,\mathrm{d}v
+ \int_{\partial\Omega_{t}} \bar{t}_{i}\,\delta u_{i}\,
  \mathrm{d}a
- \int_{\Omega} \rho\,\ddot{u}_{i}\,\delta u_{i}\,\mathrm{d}v.
}
\end{equation}

This is \emph{exactly} the validated 2022 Mindlin Form-II weak form.
$\checkmark$

\subsection{Strong form recovery}

The strong form reduces to:

\begin{equation}
-P_{iJ,J}+(\Sigma_{IJK,K}F_{iI})_{,J}-f_{0i}+\rho_{0}\ddot{u}_{i}=0
\;\;\to\;\;
-\sigma_{ij,j}+h_{ijk,jk}-b_{i}+\rho\ddot{u}_{i}=0,
\end{equation}

i.e.,

\begin{equation}
\sigma_{ij,j}-h_{ijk,jk}+b_{i} = \rho\ddot{u}_{i}.
\end{equation}

Setting $h_{ijk}=0$: $\sigma_{ij,j}+b_{i}=\rho\ddot{u}_{i}$
(classical Cauchy/Navier equation). $\checkmark$

\subsection{Boundary condition recovery}

\begin{align}
\text{Essential:}&\quad
u_{i}=\bar{u}_{i},\quad
\partial u_{i}/\partial n = \bar{g}_{i}.
\\
\text{Natural:}&\quad
(P_{iJ}-\Sigma_{IJK,K}F_{iI})N_{J}
\to (\sigma_{ij}-h_{ijk,k})n_{j} = \bar{t}_{i},
\\
&\quad
\mathcal{D}_{i}=\mathcal{H}_{iJK}N_{K}N_{J}
\to h_{ijk}n_{k}n_{j} = \bar{d}_{i}.
\end{align}

These are exactly the 2022 benchmark boundary conditions. $\checkmark$

\textbf{Mandatory recovery rule: SATISFIED at the weak-form level.}

% =====================================================================
\section{Summary of the complete derivation}
% =====================================================================

\begin{enumerate}
\item Hamilton's principle is stated with kinetic energy, internal
energy, external virtual work, and inertia virtual work.

\item Time integration by parts converts $\delta\int\mathcal{T}\,\mathrm{d}t$
to $-\int\delta\mathcal{W}_{\mathrm{iner}}\,\mathrm{d}t$.

\item The instantaneous variational statement
$\delta\Pi_{\mathrm{int}} = \delta\mathcal{W}_{\mathrm{ext}}
- \delta\mathcal{W}_{\mathrm{iner}}$ is obtained.

\item $\delta\Pi_{\mathrm{int}}$ is expressed as
$\int[S_{IJ}\delta E_{IJ}+\Sigma_{IJK}\delta\mathcal{G}_{IJK}]\mathrm{d}V_0$.

\item The classical branch simplifies: $S_{IJ}\delta E_{IJ}
= P_{iJ}\delta u_{i,J}$ (via index swap using $S_{IJ}=S_{JI}$).

\item IBP on classical branch gives domain term $-P_{iJ,J}\delta u_i$
and boundary term $P_{iJ}N_J\delta u_i$.

\item First IBP on gradient branch (over $X_K$) gives boundary term
$\Sigma_{IJK}N_K\delta E_{IJ}$ and domain term
$-\Sigma_{IJK,K}\delta E_{IJ}$.

\item Index manipulation (swap $I\leftrightarrow J$ using
$\Sigma_{IJK}=\Sigma_{JIK}$) converts the domain term to
$-\Sigma_{IJK,K}F_{iI}\delta u_{i,J}$.

\item Second IBP (over $X_J$) gives domain term
$(\Sigma_{IJK,K}F_{iI})_{,J}\delta u_i$ and boundary term
$-(\Sigma_{IJK,K}F_{iI})N_J\delta u_i$.

\item The higher-order boundary term is correctly rewritten as
$\mathcal{H}_{lJK}N_K\delta u_{l,J}$ (scalar, all indices contracted).

\item Setting the domain coefficient of $\delta u_i$ to zero gives
the strong form:
$-P_{iJ,J}+(\Sigma_{IJK,K}F_{iI})_{,J}-f_{0i}+\rho_0\ddot{u}_i=0$.

\item Boundary terms give the natural BCs:
effective traction $(P_{iJ}-\Sigma_{IJK,K}F_{iI})N_J=\bar{T}_{0i}$
and double traction $\mathcal{D}_i=\mathcal{H}_{iJK}N_KN_J=\bar{d}_i$.

\item Small-strain limit reduces every quantity exactly to the
2022 benchmark: $\sigma_{ij,j}-h_{ijk,jk}+b_i=\rho\ddot{u}_i$, with
matching boundary conditions and weak form.

\item The mandatory recovery rule is satisfied at every level.
\end{enumerate}



% =====================================================================
% SECTION 6 — FINITE ELEMENT DISCRETIZATION
% Master Version 1.2 — Frozen Sections 1-5 unchanged.
% Complete FE discretization using C1 Bogner-Fox-Schmit elements.
% Uses ONLY frozen notation from Sections 1-5.
% Compile with pdflatex + bibtex (or Overleaf).
% =====================================================================
\documentclass[11pt]{article}

\usepackage[margin=1in]{geometry}
\usepackage{amsmath,amssymb,bm}
\usepackage{booktabs}
\usepackage{longtable}
\usepackage{cite}
\usepackage{array}

\DeclareMathOperator{\tr}{tr}
\DeclareMathOperator{\sym}{sym}

\newcommand{\gradz}{\nabla_{0}}
\newcommand{\bX}{\mathbf{X}}
\newcommand{\bx}{\mathbf{x}}
\newcommand{\bp}{\boldsymbol{\varphi}}
\newcommand{\bu}{\mathbf{u}}
\newcommand{\bF}{\mathbf{F}}
\newcommand{\bC}{\mathbf{C}}
\newcommand{\bb}{\mathbf{b}}
\newcommand{\bE}{\mathbf{E}}
\newcommand{\bep}{\boldsymbol{\varepsilon}}
\newcommand{\bH}{\mathbf{H}}
\newcommand{\bG}{\bm{\mathcal{G}}}
\newcommand{\bR}{\mathbf{R}}
\newcommand{\bS}{\mathbf{S}}
\newcommand{\bSig}{\boldsymbol{\Sigma}}
\newcommand{\bD}{\mathbb{D}}
\newcommand{\bCC}{\mathbb{C}}

\title{Section~6: Finite Element Discretization using $C^{1}$ Bogner--Fox--Schmit Elements \\
\large Master Version~1.2 (Frozen Sections~1--5 + New Section~6)}
\author{Internal working document -- complete FE discretization for MATLAB implementation}
\date{}


\subsection{Introduction and purpose}
\label{ssec:fem_intro}

This section derives the complete finite element discretization of the
frozen weak formulation from Sections~1--5 using $C^{1}$ continuous
Bogner--Fox--Schmit (BFS) rectangular elements. The result is the
discrete nonlinear system of equations that forms the basis for the
MATLAB implementation in the next section.

The BFS element is a bicubic Hermite rectangular element that provides
$C^{1}$ continuity of both the displacement field $\bu$ and its gradient
$\gradz\bu$ across element boundaries. This $C^{1}$ continuity is
mandatory for strain-gradient elasticity because the strain gradient
$\bm{\mathcal{G}}=\gradz\bE$ involves second derivatives of the
displacement field, which must be square-integrable
(Section~1, \S\ref{sec:bfs}).

The derivation proceeds by:
\begin{enumerate}
\item Defining the BFS shape functions and nodal degrees of freedom.
\item Interpolating the displacement field and its derivatives.
\item Computing the kinematic quantities ($\bF$, $\bE$, $\bm{\mathcal{G}}$)
  from the interpolated displacement.
\item Substituting into the frozen weak form (Section~4).
\item Extracting the discrete residual vector and the final nonlinear
  system.
\end{enumerate}

The resulting discrete equations are exact (no approximation beyond the
FE interpolation) and reduce exactly to the validated 2022 benchmark in
the infinitesimal-strain limit.

Supporting references: the BFS element was introduced by
Bogner, Fox, and Schmitt \cite{BognerFoxSchmitt1965}. Its application
to strain-gradient elasticity follows Amanatidou \& Aravas
\cite{AmanatidouAravas2002} and Engel et al.\
\cite{EngelGarikipati2002}. The general FEM framework follows
Bathe \cite{Bathe2006}, Belytschko et al.\ \cite{Belytschko2000},
Zienkiewicz \& Taylor \cite{ZienkiewiczTaylor2000}, and Hughes
\cite{Hughes1987}. The $C^{1}$ continuity requirement for higher-order
continua is discussed by Toupin \cite{Toupin1964} and Germain
\cite{Germain1973}.

% =====================================================================
\subsection{Why $C^{1}$ continuity is mandatory}
\label{ssec:c1_requirement}
% =====================================================================

The weak form (Section~4) contains the term
$\int_{\mathcal{B}_{0}}\Sigma_{IJK}\,\delta\mathcal{G}_{IJK}\,\mathrm{d}V_{0}$,
where $\delta\mathcal{G}_{IJK}=(\delta E_{IJ})_{,K}$ involves first
derivatives of the virtual strain $\delta E_{IJ}$. Since
$\delta E_{IJ}$ itself contains first derivatives of the virtual
displacement $\delta\bu$, the term $\delta\mathcal{G}_{IJK}$ contains
\emph{second} derivatives of $\delta\bu$.

For the integral to be well-defined, the test and trial functions must
have square-integrable second derivatives, i.e.,
$\bu,\delta\bu\in H^{2}(\mathcal{B}_{0})$. This is the functional
setting for fourth-order PDEs and requires $C^{1}$ continuity across
element boundaries.

\textbf{Why Lagrange $C^{0}$ elements are insufficient:}

Standard Lagrange elements provide only $C^{0}$ continuity (continuous
displacement, discontinuous first derivatives). Their second derivatives
are distributions (Dirac delta functions at element boundaries) and are
not square-integrable. Therefore, the strain-gradient term
$\int\Sigma_{IJK}\,\delta\mathcal{G}_{IJK}\,\mathrm{d}V_{0}$ is
ill-defined for $C^{0}$ elements.

To use $C^{0}$ elements, one would need to reformulate the problem as a
mixed formulation, introducing the strain gradient or double stress as
an independent field variable with its own interpolation
\cite{AmanatidouAravas2002}. This increases the number of unknowns and
complicates the implementation.

The BFS $C^{1}$ element avoids this complication by directly providing
the required $H^{2}$ regularity. Each node carries four degrees of
freedom (displacement, two first derivatives, and one mixed second
derivative), which is the minimal set required for $C^{1}$ continuity
in 2D.

% =====================================================================
\subsection{Finite element approximation}
\label{ssec:fe_approx}
% =====================================================================

\subsubsection{Element geometry}

Consider a rectangular element $e$ with four nodes labeled $1,2,3,4$ in
counterclockwise order. The element occupies the region
$\mathcal{B}_{0}^{e}\subset\mathcal{B}_{0}$ with material coordinates
$\bX=(X_{1},X_{2})$.

For a BFS element, the element is typically defined in the parent
domain $(\xi,\eta)\in[-1,1]^{2}$, with the mapping to physical
coordinates given by:
\begin{equation}
X_{1} = \frac{h_{1}}{2}(\xi+1) + X_{1}^{\min},
\qquad
X_{2} = \frac{h_{2}}{2}(\eta+1) + X_{2}^{\min},
\end{equation}
where $h_{1}$ and $h_{2}$ are the element dimensions in the $X_{1}$ and
$X_{2}$ directions, respectively, and $(X_{1}^{\min},X_{2}^{\min})$ are
the coordinates of the lower-left corner.

\subsubsection{BFS shape functions}

The BFS element uses bicubic Hermite interpolation. In each direction
$(\xi,\eta)$, the Hermite cubic shape functions are:
\begin{align}
H_{0}^{(1)}(\xi) &= \tfrac14(1-\xi)^{2}(2+\xi),
&
H_{1}^{(1)}(\xi) &= \tfrac14(1-\xi)^{2}(1+\xi),
\label{s6:H1}\\
H_{0}^{(2)}(\xi) &= \tfrac14(1+\xi)^{2}(2-\xi),
&
H_{1}^{(2)}(\xi) &= -\tfrac14(1+\xi)^{2}(1-\xi).
\label{s6:H2}
\end{align}

These satisfy the interpolation conditions at $\xi=\pm 1$:
\begin{equation}
H_{0}^{(1)}(\pm 1)=\delta_{1,\pm 1},
\qquad
\frac{\mathrm{d}H_{0}^{(1)}}{\mathrm{d}\xi}(\pm 1)=0,
\qquad
H_{1}^{(1)}(\pm 1)=0,
\qquad
\frac{\mathrm{d}H_{1}^{(1)}}{\mathrm{d}\xi}(\pm 1)=\delta_{1,\pm 1}.
\end{equation}

The 2D BFS shape functions are tensor products of the 1D Hermite
functions. For node $a\in\{1,2,3,4\}$ located at $(\xi_{a},\eta_{a})$,
the four shape functions are:
\begin{align}
N_{a}^{(1)}(\xi,\eta) &= H_{0}^{(i)}(\xi)\,H_{0}^{(j)}(\eta),
\label{s6:N1}\\
N_{a}^{(2)}(\xi,\eta) &= H_{1}^{(i)}(\xi)\,H_{0}^{(j)}(\eta)\cdot\frac{h_{1}}{2},
\label{s6:N2}\\
N_{a}^{(3)}(\xi,\eta) &= H_{0}^{(i)}(\xi)\,H_{1}^{(j)}(\eta)\cdot\frac{h_{2}}{2},
\label{s6:N3}\\
N_{a}^{(4)}(\xi,\eta) &= H_{1}^{(i)}(\xi)\,H_{1}^{(j)}(\eta)\cdot\frac{h_{1}h_{2}}{4},
\label{s6:N4}
\end{align}
where $(i,j)\in\{(1,1),(2,1),(2,2),(1,2)\}$ correspond to nodes
$1,2,3,4$, respectively. The scaling factors ensure that $N_{a}^{(2)}$
and $N_{a}^{(3)}$ have units of length (to interpolate derivatives
$\partial u/\partial X_{1}$ and $\partial u/\partial X_{2}$), and
$N_{a}^{(4)}$ has units of length squared (to interpolate the mixed
derivative $\partial^{2}u/(\partial X_{1}\partial X_{2})$).

\textbf{Interpolation properties:}
\begin{align}
N_{a}^{(1)}(\xi_{b},\eta_{b}) &= \delta_{ab},
&
\frac{\partial N_{a}^{(1)}}{\partial X_{1}}(\xi_{b},\eta_{b}) &= 0,
\label{s6:interp1}\\
N_{a}^{(2)}(\xi_{b},\eta_{b}) &= 0,
&
\frac{\partial N_{a}^{(2)}}{\partial X_{1}}(\xi_{b},\eta_{b}) &= \delta_{ab},
\label{s6:interp2}\\
N_{a}^{(3)}(\xi_{b},\eta_{b}) &= 0,
&
\frac{\partial N_{a}^{(3)}}{\partial X_{2}}(\xi_{b},\eta_{b}) &= \delta_{ab},
\label{s6:interp3}\\
N_{a}^{(4)}(\xi_{b},\eta_{b}) &= 0,
&
\frac{\partial^{2}N_{a}^{(4)}}{\partial X_{1}\partial X_{2}}(\xi_{b},\eta_{b}) &= \delta_{ab}.
\label{s6:interp4}
\end{align}

\subsubsection{Nodal degrees of freedom}

For each spatial component $i\in\{1,2\}$ (2D problem), each node $a$
carries four degrees of freedom representing the displacement field and
its derivatives at the node:
\begin{align}
u_{i}^{(a)} &:= u_{i}(\xi_{a},\eta_{a}),
\label{s6:DOF1}\\
\left(\frac{\partial u_{i}}{\partial X_{1}}\right)^{(a)} &:= \frac{\partial u_{i}}{\partial X_{1}}(\xi_{a},\eta_{a}),
\label{s6:DOF2}\\
\left(\frac{\partial u_{i}}{\partial X_{2}}\right)^{(a)} &:= \frac{\partial u_{i}}{\partial X_{2}}(\xi_{a},\eta_{a}),
\label{s6:DOF3}\\
\left(\frac{\partial^{2}u_{i}}{\partial X_{1}\partial X_{2}}\right)^{(a)} &:= \frac{\partial^{2}u_{i}}{\partial X_{1}\partial X_{2}}(\xi_{a},\eta_{a}).
\label{s6:DOF4}
\end{align}

These nodal variables are the minimal set required for $C^{1}$ continuity
in 2D. They interpolate the field values $u_{i}(\mathbf{X})$ and the
gradient components $u_{i,1}(\mathbf{X})$, $u_{i,2}(\mathbf{X})$,
$u_{i,12}(\mathbf{X})$ defined in Sections~1--5.

For a 2D problem with 4 nodes per element, the total number of DOFs per
element is $4\text{ nodes}\times 2\text{ components}\times 4\text{ DOFs}=32$.

The element DOF vector is:
\begin{equation}
\mathbf{d}^{e}
= \bigl[\,
   u_{1}^{(1)},\,u_{1,1}^{(1)},\,u_{1,2}^{(1)},\,u_{1,12}^{(1)},
   \,u_{2}^{(1)},\,u_{2,1}^{(1)},\,u_{2,2}^{(1)},\,u_{2,12}^{(1)},
\nonumber\\
   u_{1}^{(2)},\,u_{1,1}^{(2)},\,u_{1,2}^{(2)},\,u_{1,12}^{(2)},
   \,u_{2}^{(2)},\,u_{2,1}^{(2)},\,u_{2,2}^{(2)},\,u_{2,12}^{(2)},
\nonumber\\
   u_{1}^{(3)},\,\ldots,\,u_{2}^{(4)},\,u_{2,12}^{(4)}
\,\bigr]^{T}
\in\mathbb{R}^{32}.
\label{s6:d_elem}
\end{equation}

\subsection{Displacement interpolation}
\label{ssec:disp_interp}
% =====================================================================

\subsubsection{Field interpolation}

The displacement field within element $e$ is approximated as:
\begin{equation}
u_{i}^{h}(\bX)
= \sum_{a=1}^{4}\bigl[\,
   N_{a}^{(1)}(\xi,\eta)\,u_{i}^{(a)}
 + N_{a}^{(2)}(\xi,\eta)\,u_{i,1}^{(a)}
 + N_{a}^{(3)}(\xi,\eta)\,u_{i,2}^{(a)}
 + N_{a}^{(4)}(\xi,\eta)\,u_{i,12}^{(a)}
\,\bigr].
\label{s6:u_interp}
\end{equation}

This can be written in matrix form as:
\begin{equation}
\boxed{
u_{i}^{h}(\bX) = \mathbf{N}_{i}^{T}(\bX)\,\mathbf{d}_{i}^{e},
}
\label{s6:u_matrix}
\end{equation}
where $\mathbf{d}_{i}^{e}\in\mathbb{R}^{16}$ is the DOF vector for
component $i$:
\begin{equation}
\mathbf{d}_{i}^{e}
= \bigl[\,
   u_{i}^{(1)},\,u_{i,1}^{(1)},\,u_{i,2}^{(1)},\,u_{i,12}^{(1)},
   \,u_{i}^{(2)},\,\ldots,\,u_{i}^{(4)},\,u_{i,12}^{(4)}
\,\bigr]^{T},
\end{equation}
and $\mathbf{N}_{i}(\bX)\in\mathbb{R}^{16}$ is the shape function vector:
\begin{equation}
\mathbf{N}_{i}(\bX)
= \bigl[\,
   N_{1}^{(1)},\,N_{1}^{(2)},\,N_{1}^{(3)},\,N_{1}^{(4)},
   \,N_{2}^{(1)},\,\ldots,\,N_{4}^{(4)}
\,\bigr]^{T}.
\end{equation}

\textbf{Matrix dimensions:}
\begin{itemize}
\item $u_{i}^{h}$: scalar (displacement component $i$)
\item $\mathbf{N}_{i}$: $16\times 1$ vector
\item $\mathbf{d}_{i}^{e}$: $16\times 1$ vector
\item $\mathbf{N}_{i}^{T}\,\mathbf{d}_{i}^{e}$: scalar $\checkmark$
\end{itemize}

\subsubsection{Virtual displacement interpolation}

The virtual displacement field uses the same interpolation:
\begin{equation}
\delta u_{i}^{h}(\bX) = \mathbf{N}_{i}^{T}(\bX)\,\delta\mathbf{d}_{i}^{e},
\label{s6:delta_u_matrix}
\end{equation}
where $\delta\mathbf{d}_{i}^{e}\in\mathbb{R}^{16}$ is the virtual DOF
vector:
\begin{equation}
\delta\mathbf{d}_{i}^{e}
= \bigl[\,
   \delta u_{i}^{(1)},\,\delta u_{i,1}^{(1)},\,\delta u_{i,2}^{(1)},\,\delta u_{i,12}^{(1)},
   \,\ldots,\,
   \delta u_{i}^{(4)},\,\delta u_{i,12}^{(4)}
\,\bigr]^{T}.
\end{equation}

% =====================================================================
\subsection{Interpolation of displacement gradient}
\label{ssec:grad_interp}
% =====================================================================

\subsubsection{First derivatives}

The displacement gradient is obtained by differentiating
\eqref{s6:u_interp}:
\begin{equation}
u_{i,J}^{h}(\bX)
= \sum_{a=1}^{4}\bigl[\,
   N_{a,J}^{(1)}\,u_{i}^{(a)}
 + N_{a,J}^{(2)}\,u_{i,1}^{(a)}
 + N_{a,J}^{(3)}\,u_{i,2}^{(a)}
 + N_{a,J}^{(4)}\,u_{i,12}^{(a)}
\,\bigr].
\label{s6:u_grad_interp}
\end{equation}

This can be written as:
\begin{equation}
\boxed{
u_{i,J}^{h}(\bX) = \mathbf{B}_{iJ}^{T}(\bX)\,\mathbf{d}_{i}^{e},
}
\label{s6:u_grad_matrix}
\end{equation}
where $\mathbf{B}_{iJ}(\bX)\in\mathbb{R}^{16}$ is the gradient shape
function vector:
\begin{equation}
\mathbf{B}_{iJ}(\bX)
= \bigl[\,
   N_{1,J}^{(1)},\,N_{1,J}^{(2)},\,N_{1,J}^{(3)},\,N_{1,J}^{(4)},
   \,N_{2,J}^{(1)},\,\ldots,\,N_{4,J}^{(4)}
\,\bigr]^{T}.
\end{equation}

\textbf{Matrix dimensions:}
\begin{itemize}
\item $u_{i,J}^{h}$: scalar (component $i$, derivative index $J$)
\item $\mathbf{B}_{iJ}$: $16\times 1$ vector
\item $\mathbf{B}_{iJ}^{T}\,\mathbf{d}_{i}^{e}$: scalar $\checkmark$
\end{itemize}

\subsubsection{Second derivatives (required for strain gradient)}

The second derivatives are:
\begin{equation}
u_{i,JK}^{h}(\bX)
= \sum_{a=1}^{4}\bigl[\,
   N_{a,JK}^{(1)}\,u_{i}^{(a)}
 + N_{a,JK}^{(2)}\,u_{i,1}^{(a)}
 + N_{a,JK}^{(3)}\,u_{i,2}^{(a)}
 + N_{a,JK}^{(4)}\,u_{i,12}^{(a)}
\,\bigr].
\label{s6:u_2nd_interp}
\end{equation}

This can be written as:
\begin{equation}
\boxed{
u_{i,JK}^{h}(\bX) = \mathbf{G}_{iJK}^{T}(\bX)\,\mathbf{d}_{i}^{e},
}
\label{s6:u_2nd_matrix}
\end{equation}
where $\mathbf{G}_{iJK}(\bX)\in\mathbb{R}^{16}$ is the second-gradient
shape function vector:
\begin{equation}
\mathbf{G}_{iJK}(\bX)
= \bigl[\,
   N_{1,JK}^{(1)},\,N_{1,JK}^{(2)},\,N_{1,JK}^{(3)},\,N_{1,JK}^{(4)},
   \,N_{2,JK}^{(1)},\,\ldots,\,N_{4,JK}^{(4)}
\,\bigr]^{T}.
\end{equation}

\textbf{Matrix dimensions:}
\begin{itemize}
\item $u_{i,JK}^{h}$: scalar (component $i$, derivative indices $J,K$)
\item $\mathbf{G}_{iJK}$: $16\times 1$ vector
\item $\mathbf{G}_{iJK}^{T}\,\mathbf{d}_{i}^{e}$: scalar $\checkmark$
\end{itemize}

\textbf{Note:} The BFS element provides all second derivatives
$u_{i,11}^{h}$, $u_{i,22}^{h}$, and $u_{i,12}^{h}=u_{i,21}^{h}$ as
continuous fields within each element. This is the key advantage over
$C^{0}$ elements.

% =====================================================================
\subsection{Interpolation of Green--Lagrange strain}
\label{ssec:strain_interp}
% =====================================================================

\subsubsection{Strain interpolation}

The Green--Lagrange strain is:
\begin{equation}
E_{IJ}^{h}(\bX)
= \tfrac12\bigl(F_{kI}^{h}\,F_{kJ}^{h} - \delta_{IJ}\bigr),
\end{equation}
where the deformation gradient is:
\begin{equation}
F_{iJ}^{h} = \delta_{iJ} + u_{i,J}^{h}
= \delta_{iJ} + \mathbf{B}_{iJ}^{T}\,\mathbf{d}_{i}^{e}.
\label{s6:F_h}
\end{equation}

Substituting \eqref{s6:F_h}:
\begin{equation}
\boxed{
E_{IJ}^{h}(\bX)
= \tfrac12\bigl[\,
   (\delta_{kI}+\mathbf{B}_{kI}^{T}\mathbf{d}_{k}^{e})\,
   (\delta_{kJ}+\mathbf{B}_{kJ}^{T}\mathbf{d}_{k}^{e})
   - \delta_{IJ}
\,\bigr].
}
\label{s6:E_h}
\end{equation}

This is a \emph{nonlinear} function of the DOF vector
$\mathbf{d}_{i}^{e}$ due to the product of two deformation gradients.

\textbf{Matrix dimensions:}
\begin{itemize}
\item $E_{IJ}^{h}$: scalar (strain component $IJ$)
\item $F_{kI}^{h}$: scalar
\item $F_{kJ}^{h}$: scalar
\item Product: scalar $\checkmark$
\end{itemize}

\subsubsection{Strain variation}

The variation of the strain is (from Section~4, \eqref{s5:dE_def}):
\begin{equation}
\delta E_{IJ}^{h}
= \tfrac12\bigl(F_{kI}^{h}\,\delta u_{k,J}^{h}
              + \delta u_{k,I}^{h}\,F_{kJ}^{h}\bigr).
\end{equation}

Substituting \eqref{s6:delta_u_matrix} differentiated:
\begin{equation}
\delta u_{k,J}^{h} = \mathbf{B}_{kJ}^{T}\,\delta\mathbf{d}_{k}^{e},
\end{equation}
we obtain:
\begin{equation}
\boxed{
\delta E_{IJ}^{h}
= \tfrac12\bigl(\mathbf{B}_{kJ}^{T}\,\delta\mathbf{d}_{k}^{e}\,F_{kI}^{h}
              + F_{kJ}^{h}\,\mathbf{B}_{kI}^{T}\,\delta\mathbf{d}_{k}^{e}\bigr)
= \bigl(\mathbf{B}_{\delta E,IJ}^{h}\bigr)^{T}\,\delta\mathbf{d}_{k}^{e},
}
\label{s6:deltaE_h}
\end{equation}
where the strain-displacement matrix is:
\begin{equation}
\mathbf{B}_{\delta E,IJ}^{h}
= \tfrac12\bigl(F_{kI}^{h}\,\mathbf{B}_{kJ} + F_{kJ}^{h}\,\mathbf{B}_{kI}\bigr).
\label{s6:B_deltaE}
\end{equation}

\textbf{Matrix dimensions:}
\begin{itemize}
\item $\delta E_{IJ}^{h}$: scalar
\item $\mathbf{B}_{\delta E,IJ}^{h}$: $16\times 1$ vector (for each $k$)
\item $\delta\mathbf{d}_{k}^{e}$: $16\times 1$ vector
\item Product: scalar $\checkmark$
\end{itemize}

\textbf{Note:} The strain-displacement matrix $\mathbf{B}_{\delta E,IJ}^{h}$
depends on the current deformation through $F_{kI}^{h}$ and $F_{kJ}^{h}$,
making it configuration-dependent (geometric nonlinearity).

% =====================================================================
\subsection{Interpolation of strain gradient}
\label{ssec:grad_strain_interp}
% =====================================================================

\subsubsection{Strain gradient interpolation}

The strain gradient is:
\begin{equation}
\mathcal{G}_{IJK}^{h}(\bX)
= E_{IJ,K}^{h}(\bX)
= \tfrac12\bigl(F_{mI,K}^{h}\,F_{mJ}^{h}
              + F_{mI}^{h}\,F_{mJ,K}^{h}\bigr).
\end{equation}

The derivative of the deformation gradient is:
\begin{equation}
F_{mI,K}^{h} = u_{m,I,K}^{h} = \mathbf{G}_{mIK}^{T}\,\mathbf{d}_{m}^{e}.
\end{equation}

Substituting:
\begin{equation}
\boxed{
\mathcal{G}_{IJK}^{h}
= \tfrac12\bigl[\,
   (\mathbf{G}_{mIK}^{T}\mathbf{d}_{m}^{e})\,F_{mJ}^{h}
 + F_{mI}^{h}\,(\mathbf{G}_{mJK}^{T}\mathbf{d}_{m}^{e})
\,\bigr].
}
\label{s6:G_h}
\end{equation}

This is also a \emph{nonlinear} function of the DOF vector, but now
involves second derivatives of the shape functions.

\textbf{Matrix dimensions:}
\begin{itemize}
\item $\mathcal{G}_{IJK}^{h}$: scalar (gradient component $IJK$)
\item $\mathbf{G}_{mIK}^{T}\mathbf{d}_{m}^{e}$: scalar
\item $F_{mJ}^{h}$: scalar
\item Product: scalar $\checkmark$
\end{itemize}

\subsubsection{Strain gradient variation}

The variation of the strain gradient is (from Section~5, \eqref{s5:Delta_dG}):
\begin{equation}
\delta\mathcal{G}_{IJK}^{h}
= (\delta E_{IJ}^{h})_{,K}
= \tfrac12\bigl(F_{mI,K}^{h}\,\delta u_{m,J}^{h}
              + F_{mI}^{h}\,\delta u_{m,J,K}^{h}
              + \delta u_{m,I,K}^{h}\,F_{mJ}^{h}
              + \delta u_{m,I}^{h}\,F_{mJ,K}^{h}\bigr).
\end{equation}

Substituting the interpolations:
\begin{align}
\delta u_{m,J}^{h} &= \mathbf{B}_{mJ}^{T}\,\delta\mathbf{d}_{m}^{e},
\label{s6:delta_u_J}\\
\delta u_{m,J,K}^{h} &= \mathbf{G}_{mJK}^{T}\,\delta\mathbf{d}_{m}^{e},
\label{s6:delta_u_JK}\\
F_{mI,K}^{h} &= \mathbf{G}_{mIK}^{T}\,\mathbf{d}_{m}^{e},
\label{s6:F_mIK}
\end{align}
we obtain:
\begin{equation}
\delta\mathcal{G}_{IJK}^{h}
= \tfrac12\bigl[\,
   (\mathbf{G}_{mIK}^{T}\mathbf{d}_{m}^{e})\,(\mathbf{B}_{mJ}^{T}\delta\mathbf{d}_{m}^{e})
 + F_{mI}^{h}\,(\mathbf{G}_{mJK}^{T}\delta\mathbf{d}_{m}^{e})
 + (\mathbf{G}_{mIK}^{T}\delta\mathbf{d}_{m}^{e})\,F_{mJ}^{h}
 + (\mathbf{B}_{mI}^{T}\delta\mathbf{d}_{m}^{e})\,(\mathbf{G}_{mJK}^{T}\mathbf{d}_{m}^{e})
\,\bigr].
\label{s6:deltaG_h_raw}
\end{equation}

This can be written as:
\begin{equation}
\boxed{
\delta\mathcal{G}_{IJK}^{h}
= \bigl(\mathbf{B}_{\delta\mathcal{G},IJK}^{h}\bigr)^{T}\,\delta\mathbf{d}_{m}^{e},
}
\label{s6:deltaG_h}
\end{equation}
where the higher-order strain-displacement matrix is:
\begin{align}
\mathbf{B}_{\delta\mathcal{G},IJK}^{h}
&= \tfrac12\bigl[\,
   (\mathbf{G}_{mIK}^{T}\mathbf{d}_{m}^{e})\,\mathbf{B}_{mJ}
 + F_{mI}^{h}\,\mathbf{G}_{mJK}
 + F_{mJ}^{h}\,\mathbf{G}_{mIK}
 + (\mathbf{G}_{mJK}^{T}\mathbf{d}_{m}^{e})\,\mathbf{B}_{mI}
\,\bigr].
\label{s6:B_deltaG}
\end{align}

\textbf{Matrix dimensions:}
\begin{itemize}
\item $\delta\mathcal{G}_{IJK}^{h}$: scalar
\item $\mathbf{B}_{\delta\mathcal{G},IJK}^{h}$: $16\times 1$ vector (for each $m$)
\item $\delta\mathbf{d}_{m}^{e}$: $16\times 1$ vector
\item Product: scalar $\checkmark$
\end{itemize}

\textbf{Key observation:} The higher-order strain-displacement matrix
$\mathbf{B}_{\delta\mathcal{G},IJK}^{h}$ depends on:
\begin{enumerate}
\item The current deformation gradient $F_{mI}^{h}$, $F_{mJ}^{h}$
  (geometric nonlinearity).
\item The current second derivatives of displacement
  $\mathbf{G}_{mIK}^{T}\mathbf{d}_{m}^{e}$, $\mathbf{G}_{mJK}^{T}\mathbf{d}_{m}^{e}$
  (initial curvature effect).
\item The first-derivative shape functions $\mathbf{B}_{mI}$, $\mathbf{B}_{mJ}$.
\item The second-derivative shape functions $\mathbf{G}_{mIK}$, $\mathbf{G}_{mJK}$.
\end{enumerate}

This confirms that the consistent linearization (Section~5) will produce
both material and geometric tangent contributions.

% =====================================================================
\subsection{Discrete weak form}
\label{ssec:discrete_weak}
% =====================================================================

\subsubsection{Assembly of element contributions}

The continuous weak form (Section~4, \eqref{s5:residual}) is:
\begin{equation}
G(\mathbf{u};\,\delta\mathbf{u})
= \int_{\mathcal{B}_{0}}
  \bigl[\,P_{iJ}\,\delta u_{i,J}
  + \Sigma_{IJK}\,\delta\mathcal{G}_{IJK}\bigr]\,
  \mathrm{d}V_{0}
- \delta\mathcal{W}_{\mathrm{ext}}
+ \delta\mathcal{W}_{\mathrm{iner}}
= 0.
\end{equation}

Substituting the FE interpolations \eqref{s6:deltaE_h} and
\eqref{s6:deltaG_h}:
\begin{equation}
G^{h}(\mathbf{d};\,\delta\mathbf{d})
= \sum_{e}\int_{\mathcal{B}_{0}^{e}}
  \bigl[\,P_{iJ}^{h}\,\mathbf{B}_{iJ}^{T}\,\delta\mathbf{d}_{i}^{e}
  + \Sigma_{IJK}^{h}\,\bigl(\mathbf{B}_{\delta\mathcal{G},IJK}^{h}\bigr)^{T}\,\delta\mathbf{d}_{k}^{e}
  \bigr]\,\mathrm{d}V_{0}
- \delta\mathcal{W}_{\mathrm{ext}}^{h}
+ \delta\mathcal{W}_{\mathrm{iner}}^{h}
= 0.
\end{equation}

Since $\delta\mathbf{d}^{e}$ is arbitrary, we extract the element
residual vector:
\begin{equation}
\boxed{
\mathbf{r}^{e}(\mathbf{d}^{e})
= \int_{\mathcal{B}_{0}^{e}}
  \bigl[\,\mathbf{B}_{iJ}\,P_{iJ}^{h}
  + \mathbf{B}_{\delta\mathcal{G},IJK}^{h}\,\Sigma_{IJK}^{h}
  \bigr]\,\mathrm{d}V_{0}
- \mathbf{f}_{\mathrm{ext}}^{e}
+ \mathbf{f}_{\mathrm{iner}}^{e}
= \mathbf{0},
}
\label{s6:residual_elem}
\end{equation}
where:
\begin{align}
\mathbf{f}_{\mathrm{ext}}^{e}
&= \int_{\mathcal{B}_{0}^{e}} \mathbf{N}_{i}\,f_{0i}\,\mathrm{d}V_{0}
 + \int_{\partial\mathcal{B}_{0t}^{e}} \mathbf{N}_{i}\,\bar{T}_{0i}\,\mathrm{d}A_{0},
\label{s6:f_ext}\\
\mathbf{f}_{\mathrm{iner}}^{e}
&= \int_{\mathcal{B}_{0}^{e}} \mathbf{N}_{i}\,\rho_{0}\,\ddot{u}_{i}^{h}\,\mathrm{d}V_{0}.
\label{s6:f_iner}
\end{align}

\textbf{Matrix dimensions:}
\begin{itemize}
\item $\mathbf{r}^{e}$: $32\times 1$ element residual vector
\item $\mathbf{B}_{iJ}$: $16\times 1$ vector (for each $i,J$)
\item $P_{iJ}^{h}$: scalar
\item $\mathbf{B}_{iJ}\,P_{iJ}^{h}$: $16\times 1$ vector
\item Sum over $i,J$: $32\times 1$ vector $\checkmark$
\end{itemize}

\subsubsection{Global assembly}

The global residual vector is assembled from element contributions:
\begin{equation}
\boxed{
\mathbf{R}(\mathbf{D})
= \mathbb{A}_{e=1}^{N_{\mathrm{elem}}} \mathbf{r}^{e}(\mathbf{d}^{e})
= \mathbf{0},
}
\label{s6:residual_global}
\end{equation}
where $\mathbb{A}$ denotes the standard FE assembly operator, and
$\mathbf{D}\in\mathbb{R}^{N_{\mathrm{dof}}}$ is the global DOF vector.

The consistent linearization of this discrete residual, which produces the tangent matrix required for Newton--Raphson iteration, is developed in Section~7.

% =====================================================================
\subsection{Summary of matrices and vectors}
\label{ssec:summary}
% =====================================================================

\begin{longtable}[c]{>{\raggedright\arraybackslash}p{0.18\linewidth}
>{\raggedright\arraybackslash}p{0.15\linewidth}
>{\raggedright\arraybackslash}p{0.67\linewidth}}
\caption{Summary of FE matrices and vectors.}
\label{tab:fe_matrices}\\
\toprule
Symbol & Dimensions & Definition \\
\midrule
\endfirsthead
\toprule
Symbol & Dimensions & Definition \\
\midrule
\endhead
$\mathbf{N}_{i}$ & $16\times 1$ & Displacement shape function vector \\
$\mathbf{B}_{iJ}$ & $16\times 1$ & First-derivative shape function vector \\
$\mathbf{G}_{iJK}$ & $16\times 1$ & Second-derivative shape function vector \\
$\mathbf{B}_{\delta E,IJ}^{h}$ & $16\times 1$ & Strain-displacement matrix (nonlinear) \\
$\mathbf{B}_{\delta\mathcal{G},IJK}^{h}$ & $16\times 1$ & Higher-order strain-displacement matrix (nonlinear) \\
$\mathbf{d}_{i}^{e}$ & $16\times 1$ & Element DOF vector (component $i$) \\
$\mathbf{d}^{e}$ & $32\times 1$ & Element DOF vector (both components) \\
$\mathbf{r}^{e}$ & $32\times 1$ & Element residual vector \\
$\mathbf{R}$ & $N_{\mathrm{dof}}\times 1$ & Global residual vector \\
\bottomrule
\end{longtable}

% =====================================================================
\subsection{Small-strain limit and 2022 benchmark recovery}
\label{ssec:recovery}
% =====================================================================

In the infinitesimal-strain limit ($\|\nabla_{0}\mathbf{u}\|\to 0$):
\begin{align}
F_{iJ}^{h} &\to \delta_{iJ},
\qquad
E_{IJ}^{h} \to \varepsilon_{IJ}^{h},
\qquad
\mathcal{G}_{IJK}^{h} \to \varepsilon_{IJ,K}^{h},
\label{s6:lim_kin}\\
P_{iJ}^{h} &\to \sigma_{ij},
\qquad
\Sigma_{IJK}^{h} \to h_{ijk},
\label{s6:lim_stress}\\
\mathbf{B}_{\delta E,IJ}^{h} &\to \mathbf{B}_{\varepsilon,ij},
\qquad
\mathbf{B}_{\delta\mathcal{G},IJK}^{h} \to \mathbf{B}_{\nabla\varepsilon,ijk}.
\label{s6:lim_B}
\end{align}

The residual reduces to:
\begin{equation}
\mathbf{R} \to \mathbb{A}_{e}
  \int_{\mathcal{B}_{0}^{e}}
  \bigl[\,\mathbf{B}_{\varepsilon,ij}^{T}\,\sigma_{ij}
  + \mathbf{B}_{\nabla\varepsilon,ijk}^{T}\,h_{ijk}
  \bigr]\,\mathrm{d}V_{0}
- \mathbf{F}_{\mathrm{ext}}.
\end{equation}

This is \emph{exactly} the discrete form of the validated 2022
Mindlin Form-II weak formulation. $\checkmark$

\textbf{Mandatory recovery rule: SATISFIED at the discrete FE level.}

% =====================================================================
\subsection{Summary of Section~6}
\label{ssec:section_summary}
% =====================================================================

\begin{enumerate}
\item The BFS $C^{1}$ element provides the required $H^{2}$ regularity
  for strain-gradient elasticity.

\item Each node carries 4 DOFs: $u_{i}$, $\partial u_{i}/\partial X_{1}$,
  $\partial u_{i}/\partial X_{2}$, $\partial^{2}u_{i}/(\partial X_{1}\partial X_{2})$.

\item The displacement, first derivatives, and second derivatives are
  interpolated using the shape function vectors $\mathbf{N}_{i}$,
  $\mathbf{B}_{iJ}$, and $\mathbf{G}_{iJK}$.

\item The Green--Lagrange strain and strain gradient are computed from
  the interpolated deformation gradient and its derivatives.

\item The strain-displacement matrices $\mathbf{B}_{\delta E,IJ}^{h}$ and
  $\mathbf{B}_{\delta\mathcal{G},IJK}^{h}$ are configuration-dependent
  (geometric nonlinearity).

\item The discrete residual vector \eqref{s6:residual_elem} is assembled
  from element contributions.

\item In the small-strain limit, the discrete equations reduce exactly
  to the validated 2022 Mindlin Form-II formulation.
\end{enumerate}

\textbf{Section~6 is now complete. The consistent linearization and
tangent operator are developed in Section~7.}

% =====================================================================
% SECTION 7 — ELEMENT MATRICES, NUMERICAL INTEGRATION, AND
%             NEWTON-RAPHSON ALGORITHM
% Master Version 1.2 — Frozen Sections 1-6 unchanged.
% Complete element-level FE equations for MATLAB implementation.
% Uses ONLY frozen notation from Sections 1-6.
% Compile with pdflatex + bibtex (or Overleaf).
% =====================================================================
\documentclass[11pt]{article}

\usepackage[margin=1in]{geometry}
\usepackage{amsmath,amssymb,bm}
\usepackage{booktabs}
\usepackage{longtable}
\usepackage{cite}
\usepackage{array}
\usepackage{algorithmic}
\usepackage{algorithm}

\DeclareMathOperator{\tr}{tr}
\DeclareMathOperator{\sym}{sym}

\newcommand{\gradz}{\nabla_{0}}
\newcommand{\bX}{\mathbf{X}}
\newcommand{\bx}{\mathbf{x}}
\newcommand{\bp}{\boldsymbol{\varphi}}
\newcommand{\bu}{\mathbf{u}}
\newcommand{\bF}{\mathbf{F}}
\newcommand{\bC}{\mathbf{C}}
\newcommand{\bb}{\mathbf{b}}
\newcommand{\bE}{\mathbf{E}}
\newcommand{\bep}{\boldsymbol{\varepsilon}}
\newcommand{\bH}{\mathbf{H}}
\newcommand{\bG}{\bm{\mathcal{G}}}
\newcommand{\bR}{\mathbf{R}}
\newcommand{\bS}{\mathbf{S}}
\newcommand{\bSig}{\boldsymbol{\Sigma}}
\newcommand{\bD}{\mathbb{D}}
\newcommand{\bCC}{\mathbb{C}}

\title{Section~7: Element Matrices, Numerical Integration and Newton--Raphson Algorithm \\
\large Master Version~1.2 (Frozen Sections~1--6 + New Section~7)}
\author{Internal working document -- complete solver specification for MATLAB implementation}
\date{}


\section{Element Matrices, Numerical Integration and Newton--Raphson Algorithm}
\label{sec:solver}
% =====================================================================

\subsection{Introduction and scope}
\label{ssec:s7_intro}

This section derives the complete element-level finite element equations
and the Newton--Raphson solution algorithm for the geometrically
nonlinear, displacement-based, $C^{1}$ BFS Mindlin Form-II
strain-gradient formulation developed in Sections~1--6.

Starting from the discrete residual
\eqref{s6:residual_elem} and the consistent linearization framework
of Section~5, we derive:
\begin{enumerate}
\item The four contributions to the element tangent stiffness matrix
  (classical material, classical geometric, gradient material, gradient
  geometric).
\item The element internal, external, and inertia force vectors.
\item The numerical integration procedure for BFS elements.
\item The assembly operator mapping element quantities to the global
  system.
\item The complete Newton--Raphson iterative algorithm with convergence
  criteria and load stepping strategy.
\end{enumerate}

All derivations follow the frozen notation and constitutive
definitions of Sections~1--6. No re-derivation of constitutive
equations, weak forms, or BFS interpolation is performed.

Supporting references: the computational framework for consistent
linearization and tangent operators follows
Simo \& Hughes \cite{SimoHughes1998} and Holzapfel
\cite{Holzapfel2000}. The nonlinear solution algorithm follows
Bathe \cite{Bathe2006}, Belytschko et al.\
\cite{Belytschko2000}, and Zienkiewicz \& Taylor
\cite{ZienkiewiczTaylor2000}. The BFS element integration follows
Bogner, Fox, and Schmitt \cite{BognerFoxSchmitt1965} and
Hughes \cite{Hughes1987}.

% =====================================================================
\subsection{Element tangent matrix}
\label{ssec:tangent}
% =====================================================================

\subsubsection{Block structure of the tangent}

The element tangent stiffness matrix is the G\^ateaux derivative of the
element residual \eqref{s6:residual_elem} with respect to the element
DOF vector $\mathbf{d}^e \in \mathbb{R}^{32}$:

\begin{equation}
\boxed{
\mathbf{K}_{T}^{e}
:= \frac{\partial \mathbf{r}^{e}}{\partial \mathbf{d}^{e}}
= \mathbf{K}_{\mathrm{mat,cl}}^{e}
+ \mathbf{K}_{\mathrm{geo,cl}}^{e}
+ \mathbf{K}_{\mathrm{mat,gr}}^{e}
+ \mathbf{K}_{\mathrm{geo,gr}}^{e}
\in \mathbb{R}^{32 \times 32}.
}
\label{s7:tangent_total}
\end{equation}

The tangent is symmetric ($(\mathbf{K}_{T}^{e})^{T}=\mathbf{K}_{T}^{e}$)
because the formulation is hyperelastic
(Section~2, \eqref{eq:sec2_complete_energy_tensor}).

\textbf{Dimensional check:}
\begin{itemize}
\item $\mathbf{r}^{e} \in \mathbb{R}^{32}$ (Section~6)
\item $\mathbf{d}^{e} \in \mathbb{R}^{32}$ (Section~6)
\item $\mathbf{K}_{T}^{e} \in \mathbb{R}^{32 \times 32}$ $\checkmark$
\end{itemize}

\subsubsection{Notation for assembled shape function vectors}

For clarity, we define the \emph{assembled} shape function vectors
that operate on the full $32\times 1$ DOF vector
$\mathbf{d}^{e} = [\mathbf{d}_{1}^{e};\,\mathbf{d}_{2}^{e}]$
(where $\mathbf{d}_{k}^{e}\in\mathbb{R}^{16}$ is the DOF vector for
component $k$):

\begin{align}
\mathbf{N}_{k}^{*} &:= \begin{bmatrix} \mathbf{0}_{16\times 1} \\ \vdots \\ \mathbf{N}_{k} \\ \vdots \\ \mathbf{0}_{16\times 1} \end{bmatrix}
\in \mathbb{R}^{32},
&
\mathbf{B}_{kJ}^{*} &:= \begin{bmatrix} \mathbf{0} \\ \vdots \\ \mathbf{B}_{kJ} \\ \vdots \\ \mathbf{0} \end{bmatrix}
\in \mathbb{R}^{32},
&
\mathbf{G}_{kJK}^{*} &:= \begin{bmatrix} \mathbf{0} \\ \vdots \\ \mathbf{G}_{kJK} \\ \vdots \\ \mathbf{0} \end{bmatrix}
\in \mathbb{R}^{32}.
\label{s7:assembled_SF}
\end{align}

The vector $\mathbf{N}_{k}^{*}$ (resp.\ $\mathbf{B}_{kJ}^{*}$,
$\mathbf{G}_{kJK}^{*}$) has the $16\times 1$ shape function vector
$\mathbf{N}_{k}$ (resp.\ $\mathbf{B}_{kJ}$, $\mathbf{G}_{kJK}$) placed
in the block corresponding to spatial component $k$, with zeros in the
other component block. These allow us to write:

\begin{align}
u_{k}^{h} &= (\mathbf{N}_{k}^{*})^{T}\mathbf{d}^{e},
&
u_{k,J}^{h} &= (\mathbf{B}_{kJ}^{*})^{T}\mathbf{d}^{e},
&
u_{k,JK}^{h} &= (\mathbf{G}_{kJK}^{*})^{T}\mathbf{d}^{e}.
\label{s7:assembled_interp}
\end{align}

Similarly, we define the assembled strain-displacement vectors
$\mathbf{B}_{\delta E,IJ}^{*}\in\mathbb{R}^{32}$ and the assembled
gradient-displacement vectors
$\mathbf{B}_{\delta\mathcal{G},IJK}^{*}\in\mathbb{R}^{32}$:

\begin{align}
\mathbf{B}_{\delta E,IJ}^{*}
&= \tfrac12\sum_{k=1}^{2}\bigl(F_{kI}^{h}\,\mathbf{B}_{kJ}^{*}
                               + F_{kJ}^{h}\,\mathbf{B}_{kI}^{*}\bigr)
\in \mathbb{R}^{32},
\label{s7:B_deltaE_assembled}\\
\mathbf{B}_{\delta\mathcal{G},IJK}^{*}
&= \tfrac12\sum_{m=1}^{2}\bigl[\,(\mathbf{G}_{mIK}^{*})^{T}\mathbf{d}^{e}\,\mathbf{B}_{mJ}^{*}
 + F_{mI}^{h}\,\mathbf{G}_{mJK}^{*}
 + F_{mJ}^{h}\,\mathbf{G}_{mIK}^{*}
 + (\mathbf{G}_{mJK}^{*})^{T}\mathbf{d}^{e}\,\mathbf{B}_{mI}^{*}\,\bigr]
\in \mathbb{R}^{32}.
\label{s7:B_deltaG_assembled}
\end{align}

These satisfy:
\begin{equation}
\delta E_{IJ}^{h} = (\mathbf{B}_{\delta E,IJ}^{*})^{T}\,\delta\mathbf{d}^{e},
\qquad
\delta\mathcal{G}_{IJK}^{h} = (\mathbf{B}_{\delta\mathcal{G},IJK}^{*})^{T}\,\delta\mathbf{d}^{e}.
\end{equation}

\subsubsection{Classical material stiffness matrix}
\label{ssec:K_mat_cl}

From the consistent linearization of Section~5 (Eq.~\eqref{s5:K_mat_cl_elem}),
the classical material stiffness arises from the constitutive
linearization $\Delta S_{IJ} = \mathbb{C}_{IJKL}\,\Delta E_{KL}$:

\begin{equation}
\boxed{
\mathbf{K}_{\mathrm{mat,cl}}^{e}
= \int_{\mathcal{B}_{0}^{e}}
  \sum_{I,J,K,L}
  \mathbb{C}_{IJKL}\,
  \mathbf{B}_{\delta E,KL}^{*}\,
  (\mathbf{B}_{\delta E,IJ}^{*})^{T}\,
  \mathrm{d}V_{0}
\in \mathbb{R}^{32\times 32}.
}
\label{s7:K_mat_cl}
\end{equation}

\textbf{Proof of symmetry:}
The major symmetry $\mathbb{C}_{IJKL}=\mathbb{C}_{KLIJ}$
(Section~5, \eqref{s5:C_tensor}) ensures:
\begin{equation}
(\mathbf{K}_{\mathrm{mat,cl}}^{e})^{T}
= \int \sum_{I,J,K,L}
  \mathbb{C}_{IJKL}\,
  \mathbf{B}_{\delta E,IJ}^{*}\,
  (\mathbf{B}_{\delta E,KL}^{*})^{T}\,
  \mathrm{d}V_{0}
= \int \sum_{K,L,I,J}
  \mathbb{C}_{KLIJ}\,
  \mathbf{B}_{\delta E,KL}^{*}\,
  (\mathbf{B}_{\delta E,IJ}^{*})^{T}\,
  \mathrm{d}V_{0}
= \mathbf{K}_{\mathrm{mat,cl}}^{e}. \quad\checkmark
\end{equation}

\textbf{Derivation:}
The contribution to the virtual work is
$\int \mathbb{C}_{IJKL}\,\Delta E_{KL}\,\delta E_{IJ}\,\mathrm{d}V_{0}$.
Substituting $\Delta E_{KL} = (\mathbf{B}_{\delta E,KL}^{*})^{T}\Delta\mathbf{d}^{e}$
and $\delta E_{IJ} = (\mathbf{B}_{\delta E,IJ}^{*})^{T}\delta\mathbf{d}^{e}$:
\begin{equation}
= \delta\mathbf{d}^{T}
  \left[\int \sum_{I,J,K,L} \mathbb{C}_{IJKL}\,
  \mathbf{B}_{\delta E,KL}^{*}\,
  (\mathbf{B}_{\delta E,IJ}^{*})^{T}\,
  \mathrm{d}V_{0}\right]
  \Delta\mathbf{d}
= \delta\mathbf{d}^{T}\,\mathbf{K}_{\mathrm{mat,cl}}^{e}\,\Delta\mathbf{d}. \quad\checkmark
\end{equation}

\textbf{Matrix dimensions:}
\begin{itemize}
\item $\mathbb{C}_{IJKL}$: $[\mathrm{Pa}]$
\item $\mathbf{B}_{\delta E,KL}^{*}$: $32\times 1$
\item $(\mathbf{B}_{\delta E,IJ}^{*})^{T}$: $1\times 32$
\item Outer product: $32\times 32$
\item $\mathrm{d}V_{0}$: $[\mathrm{m}^{3}]$
\item $\mathbf{K}_{\mathrm{mat,cl}}^{e}$: $[\mathrm{Pa}\cdot\mathrm{m}^{3}] = [\mathrm{N\,m}]$ $\checkmark$
\end{itemize}

\subsubsection{Classical geometric stiffness matrix}
\label{ssec:K_geo_cl}

The classical geometric stiffness arises from the push-forward
linearization $S_{IJ}\,\Delta F_{iI}$
(Section~5, Eq.~\eqref{s5:K_geo_cl_elem}):

\begin{equation}
\boxed{
\mathbf{K}_{\mathrm{geo,cl}}^{e}
= \int_{\mathcal{B}_{0}^{e}}
  \sum_{i,I,J}
  S_{IJ}^{h}\,
  \mathbf{B}_{iJ}^{*}\,
  (\mathbf{B}_{iI}^{*})^{T}\,
  \mathrm{d}V_{0}
\in \mathbb{R}^{32\times 32}.
}
\label{s7:K_geo_cl}
\end{equation}

\textbf{Proof of symmetry:}
Swapping $I\leftrightarrow J$ and using $S_{IJ}=S_{JI}$:
\begin{equation}
(\mathbf{K}_{\mathrm{geo,cl}}^{e})^{T}
= \int \sum_{i,I,J} S_{IJ}\,
  \mathbf{B}_{iI}^{*}\,(\mathbf{B}_{iJ}^{*})^{T}\,\mathrm{d}V_{0}
= \int \sum_{i,I,J} S_{JI}\,
  \mathbf{B}_{iJ}^{*}\,(\mathbf{B}_{iI}^{*})^{T}\,\mathrm{d}V_{0}
= \mathbf{K}_{\mathrm{geo,cl}}^{e}. \quad\checkmark
\end{equation}

\textbf{Physical meaning:}
This is the initial stress (geometric stiffness) contribution. It
vanishes when $S_{IJ}=\mathbf{0}$ (reference configuration).

\subsubsection{Gradient material stiffness matrix}
\label{ssec:K_mat_gr}

The gradient material stiffness arises from the gradient constitutive
law $\Delta\Sigma_{IJK}=\mathbb{D}_{IJKPQR}\,\Delta\mathcal{G}_{PQR}$
(Section~5, Eq.~\eqref{s5:K_mat_gr_elem}):

\begin{equation}
\boxed{
\mathbf{K}_{\mathrm{mat,gr}}^{e}
= \int_{\mathcal{B}_{0}^{e}}
  \sum_{I,J,K,P,Q,R}
  \mathbb{D}_{IJKPQR}\,
  \mathbf{B}_{\delta\mathcal{G},PQR}^{*}\,
  (\mathbf{B}_{\delta\mathcal{G},IJK}^{*})^{T}\,
  \mathrm{d}V_{0}
\in \mathbb{R}^{32\times 32}.
}
\label{s7:K_mat_gr}
\end{equation}

\textbf{Proof of symmetry:}
The major symmetry $\mathbb{D}_{IJKPQR}=\mathbb{D}_{PQRIJK}$
(Section~2, \eqref{eq:sec2_D_symmetries}) ensures:
\begin{equation}
(\mathbf{K}_{\mathrm{mat,gr}}^{e})^{T}
= \mathbf{K}_{\mathrm{mat,gr}}^{e}. \quad\checkmark
\end{equation}

\textbf{Matrix dimensions:}
\begin{itemize}
\item $\mathbb{D}_{IJKPQR}$: $[\mathrm{Pa\,m^{2}}]$
\item $\mathbf{B}_{\delta\mathcal{G},PQR}^{*}$: $32\times 1$
\item $\mathrm{d}V_{0}$: $[\mathrm{m}^{3}]$
\item $\mathbf{K}_{\mathrm{mat,gr}}^{e}$: $[\mathrm{Pa\,m^{2}}\cdot\mathrm{m}^{-2}\cdot\mathrm{m}^{3}] = [\mathrm{N\,m}]$ $\checkmark$
\end{itemize}

\textbf{Computational note:}
This matrix requires second derivatives of the BFS shape functions
(through $\mathbf{B}_{\delta\mathcal{G}}^{*}$), confirming the $C^{1}$
continuity requirement established in Sections~1 and~6.

\subsubsection{Gradient geometric stiffness matrix}
\label{ssec:K_geo_gr}

The gradient geometric stiffness arises from the product-rule
contributions involving $\Sigma_{IJK}$ multiplied by displacement
gradient variations (Section~5, Eq.~\eqref{s6:K_geo_gr_elem}):

\begin{equation}
\boxed{
\mathbf{K}_{\mathrm{geo,gr}}^{e}
= \int_{\mathcal{B}_{0}^{e}}
  \sum_{i,I,J,K}
  \Sigma_{IJK}^{h}\,
  \bigl[\,\mathbf{G}_{iIK}^{*}\,(\mathbf{B}_{iJ}^{*})^{T}
       + \mathbf{B}_{iI}^{*}\,(\mathbf{G}_{iJK}^{*})^{T}\,\bigr]\,
  \mathrm{d}V_{0}
\in \mathbb{R}^{32\times 32}.
}
\label{s7:K_geo_gr}
\end{equation}

\textbf{Proof of symmetry:}
The transpose of the bracket is:
\begin{equation}
\bigl[\mathbf{G}_{iIK}^{*}(\mathbf{B}_{iJ}^{*})^{T}
 + \mathbf{B}_{iI}^{*}(\mathbf{G}_{iJK}^{*})^{T}\bigr]^{T}
= \mathbf{B}_{iJ}^{*}(\mathbf{G}_{iIK}^{*})^{T}
 + \mathbf{G}_{iJK}^{*}(\mathbf{B}_{iI}^{*})^{T}.
\end{equation}

Relabeling dummy indices $I\leftrightarrow J$ and using
$\Sigma_{IJK}=\Sigma_{JIK}$ (Section~2, \eqref{eq:sec2_D_symmetries}):
\begin{equation}
(\mathbf{K}_{\mathrm{geo,gr}}^{e})^{T}
= \int \sum_{i,J,I,K}
  \Sigma_{JIK}\,
  \bigl[\mathbf{B}_{iI}^{*}(\mathbf{G}_{iJK}^{*})^{T}
       + \mathbf{G}_{iIK}^{*}(\mathbf{B}_{iJ}^{*})^{T}\bigr]\,
  \mathrm{d}V_{0}
= \mathbf{K}_{\mathrm{geo,gr}}^{e}. \quad\checkmark
\end{equation}

\textbf{Physical meaning:}
This is the initial double-stress (gradient geometric stiffness)
contribution. It depends on the current double stress
$\Sigma_{IJK}^{h}$ and involves products of the second-derivative
shape functions $\mathbf{G}^{*}$ with the first-derivative shape
functions $\mathbf{B}^{*}$. It vanishes when
$\Sigma_{IJK}=\mathbf{0}$.

\subsubsection{Total element tangent}
\label{ssec:total_tangent}

The total element tangent stiffness matrix is:

\begin{equation}
\boxed{
\mathbf{K}_{T}^{e}
= \mathbf{K}_{\mathrm{mat,cl}}^{e}
+ \mathbf{K}_{\mathrm{geo,cl}}^{e}
+ \mathbf{K}_{\mathrm{mat,gr}}^{e}
+ \mathbf{K}_{\mathrm{geo,gr}}^{e}
\in \mathbb{R}^{32\times 32}.
}
\label{s7:K_total}
\end{equation}

The total tangent is symmetric (each contribution is individually
symmetric) and configuration-dependent (evaluated at the current
displacement $\mathbf{d}^{e}$).

\textbf{Summary table:}

\begin{longtable}[c]{>{\raggedright\arraybackslash}p{0.18\linewidth}
>{\raggedright\arraybackslash}p{0.42\linewidth}
>{\raggedright\arraybackslash}p{0.40\linewidth}}
\caption{Element tangent matrix contributions.}
\label{tab:tangent_contributions}\\
\toprule
Matrix & Expression & Physical origin \\
\midrule
\endfirsthead
\toprule
Matrix & Expression & Physical origin \\
\midrule
\endhead
$\mathbf{K}_{\mathrm{mat,cl}}^{e}$
& $\int \mathbb{C}_{IJKL}\,\mathbf{B}_{\delta E,KL}^{*}\,(\mathbf{B}_{\delta E,IJ}^{*})^{T}\,\mathrm{d}V_{0}$
& Material stiffness (classical): constitutive resistance to strain. \\
\midrule
$\mathbf{K}_{\mathrm{geo,cl}}^{e}$
& $\int S_{IJ}^{h}\,\mathbf{B}_{iJ}^{*}\,(\mathbf{B}_{iI}^{*})^{T}\,\mathrm{d}V_{0}$
& Geometric stiffness (classical): initial stress effect. \\
\midrule
$\mathbf{K}_{\mathrm{mat,gr}}^{e}$
& $\int \mathbb{D}_{IJKPQR}\,\mathbf{B}_{\delta\mathcal{G},PQR}^{*}\,(\mathbf{B}_{\delta\mathcal{G},IJK}^{*})^{T}\,\mathrm{d}V_{0}$
& Gradient material stiffness: constitutive resistance to strain gradient. \\
\midrule
$\mathbf{K}_{\mathrm{geo,gr}}^{e}$
& $\int \Sigma_{IJK}^{h}\,[\mathbf{G}_{iIK}^{*}(\mathbf{B}_{iJ}^{*})^{T}+\mathbf{B}_{iI}^{*}(\mathbf{G}_{iJK}^{*})^{T}]\,\mathrm{d}V_{0}$
& Gradient geometric stiffness: initial double-stress effect. \\
\bottomrule
\end{longtable}

% =====================================================================
\subsection{Element force vectors}
\label{ssec:force_vectors}
% =====================================================================

\subsubsection{Element internal force vector}

The element internal force vector follows from the discrete weak form
(Section~6, Eq.~\eqref{s6:residual_elem}):

\begin{equation}
\boxed{
\mathbf{r}_{\mathrm{int}}^{e}(\mathbf{d}^{e})
= \int_{\mathcal{B}_{0}^{e}}
  \left[\,\sum_{i,J}\mathbf{B}_{iJ}^{*}\,P_{iJ}^{h}
  + \sum_{I,J,K}\mathbf{B}_{\delta\mathcal{G},IJK}^{*}\,\Sigma_{IJK}^{h}\,\right]\,
  \mathrm{d}V_{0}
\in \mathbb{R}^{32}.
}
\label{s7:r_int}
\end{equation}

where:
\begin{align}
P_{iJ}^{h} &= S_{IJ}^{h}\,F_{iI}^{h}
= \bigl[\mathbb{C}_{IJKL}\,E_{KL}^{h}\bigr]\,F_{iI}^{h}
\quad\text{(first Piola--Kirchhoff stress)},
\label{s7:P_h}\\
\Sigma_{IJK}^{h} &= \mathbb{D}_{IJKPQR}\,\mathcal{G}_{PQR}^{h}
\quad\text{(double stress)}.
\label{s7:Sig_h}
\end{align}

\textbf{Matrix dimensions:}
\begin{itemize}
\item $\mathbf{B}_{iJ}^{*}$: $32\times 1$
\item $P_{iJ}^{h}$: $[\mathrm{Pa}]$
\item $\mathbf{B}_{iJ}^{*}\,P_{iJ}^{h}$: $32\times 1$
\item Sum over $i\in\{1,2\}$, $J\in\{1,2\}$: 4 terms
\item $\mathbf{B}_{\delta\mathcal{G},IJK}^{*}$: $32\times 1$
\item $\Sigma_{IJK}^{h}$: $[\mathrm{Pa\,m}]$
\item $\mathbf{B}_{\delta\mathcal{G},IJK}^{*}\,\Sigma_{IJK}^{h}$: $32\times 1$
\item Sum over $I,J,K\in\{1,2\}$: 8 terms
\item $\mathbf{r}_{\mathrm{int}}^{e}$: $32\times 1$
\item Units: $[\mathrm{Pa}\cdot\mathrm{m}^{3}] = [\mathrm{N\,m}]$ $\checkmark$
\end{itemize}

\textbf{Nonlinearity:}
$\mathbf{r}_{\mathrm{int}}^{e}$ is nonlinear in $\mathbf{d}^{e}$
because:
\begin{enumerate}
\item $P_{iJ}^{h}$ depends on $E_{KL}^{h}$ and $F_{iI}^{h}$, both
  functions of $\mathbf{d}^{e}$.
\item $\Sigma_{IJK}^{h}$ depends on $\mathcal{G}_{PQR}^{h}$, a
  function of $\mathbf{d}^{e}$.
\item $\mathbf{B}_{\delta\mathcal{G},IJK}^{*}$ depends on $F_{mI}^{h}$
  and $(\mathbf{G}_{mIK}^{*})^{T}\mathbf{d}^{e}$, both functions of
  $\mathbf{d}^{e}$.
\end{enumerate}

\subsubsection{Element external force vector}

\begin{equation}
\boxed{
\mathbf{f}_{\mathrm{ext}}^{e}
= \int_{\mathcal{B}_{0}^{e}}
  \sum_{i}\mathbf{N}_{i}^{*}\,f_{0i}\,
  \mathrm{d}V_{0}
+ \int_{\partial\mathcal{B}_{0t}^{e}}
  \sum_{i}\mathbf{N}_{i}^{*}\,\bar{T}_{0i}\,
  \mathrm{d}A_{0}
\in \mathbb{R}^{32}.
}
\label{s7:f_ext}
\end{equation}

For prescribed dead loads ($f_{0i}$ and $\bar{T}_{0i}$ independent of
deformation), $\mathbf{f}_{\mathrm{ext}}^{e}$ is \emph{constant} with
respect to $\mathbf{d}^{e}$ and does not contribute to the tangent.

\textbf{Matrix dimensions:}
\begin{itemize}
\item $\mathbf{N}_{i}^{*}$: $32\times 1$
\item $f_{0i}$: $[\mathrm{N\,m}^{-3}]$
\item $\mathbf{N}_{i}^{*}\,f_{0i}$: $32\times 1$
\item Volume integral: $[\mathrm{N\,m}^{-3}\cdot\mathrm{m}^{3}] = [\mathrm{N}]$ $\checkmark$
\item $\bar{T}_{0i}$: $[\mathrm{N\,m}^{-2}]$
\item Surface integral: $[\mathrm{N\,m}^{-2}\cdot\mathrm{m}^{2}] = [\mathrm{N}]$ $\checkmark$
\end{itemize}

\textbf{Note:} The units of $\mathbf{f}_{\mathrm{ext}}^{e}$ are $[\mathrm{N}]$ (force),
while $\mathbf{r}_{\mathrm{int}}^{e}$ has units $[\mathrm{N\,m}]$ (energy).
This apparent inconsistency is resolved by noting that $\mathbf{r}_{\mathrm{int}}^{e}$
is the \emph{energy-conjugate} residual (derivative of the internal
virtual work with respect to the DOF vector), while
$\mathbf{f}_{\mathrm{ext}}^{e}$ is the force vector. In the
discrete weak form $\mathbf{r}^{e}=\mathbf{r}_{\mathrm{int}}^{e}
- \mathbf{f}_{\mathrm{ext}}^{e} + \mathbf{f}_{\mathrm{iner}}^{e}$, all terms
must have the same units. This requires that $\mathbf{r}_{\mathrm{int}}^{e}$
and $\mathbf{f}_{\mathrm{ext}}^{e}$ are both expressed in force units,
with the BFS shape functions carrying the appropriate dimensional
scaling through their derivatives (Section~6).

\subsubsection{Element inertia vector}

\begin{equation}
\boxed{
\mathbf{f}_{\mathrm{iner}}^{e}
= \int_{\mathcal{B}_{0}^{e}}
  \sum_{i}\mathbf{N}_{i}^{*}\,\rho_{0}\,\ddot{u}_{i}^{h}\,
  \mathrm{d}V_{0}
= \mathbf{M}^{e}\,\ddot{\mathbf{d}}^{e}
\in \mathbb{R}^{32}.
}
\label{s7:f_iner}
\end{equation}

where the \emph{consistent mass matrix} is:

\begin{equation}
\boxed{
\mathbf{M}^{e}
= \int_{\mathcal{B}_{0}^{e}}
  \sum_{i}\rho_{0}\,\mathbf{N}_{i}^{*}\,(\mathbf{N}_{i}^{*})^{T}\,
  \mathrm{d}V_{0}
\in \mathbb{R}^{32\times 32}.
}
\label{s7:mass_matrix}
\end{equation}

\textbf{Proof of symmetry:}
$(\mathbf{M}^{e})^{T} = \int \rho_{0}\,\mathbf{N}_{i}^{*}\,(\mathbf{N}_{i}^{*})^{T}\,\mathrm{d}V_{0}
= \mathbf{M}^{e}$. $\checkmark$

\textbf{Proof of positive definiteness:}
For any $\mathbf{v}\neq\mathbf{0}$:
$\mathbf{v}^{T}\mathbf{M}^{e}\mathbf{v}
= \int \rho_{0}\sum_{i}((\mathbf{N}_{i}^{*})^{T}\mathbf{v})^{2}\,\mathrm{d}V_{0} > 0$. $\checkmark$

\textbf{Matrix dimensions:}
\begin{itemize}
\item $\rho_{0}$: $[\mathrm{kg\,m}^{-3}]$
\item $\mathbf{N}_{i}^{*}$: $32\times 1$
\item $\mathbf{N}_{i}^{*}(\mathbf{N}_{i}^{*})^{T}$: $32\times 32$
\item $\mathrm{d}V_{0}$: $[\mathrm{m}^{3}]$
\item $\mathbf{M}^{e}$: $[\mathrm{kg}]$ $\checkmark$
\end{itemize}

% =====================================================================
\subsection{Numerical integration}
\label{ssec:integration}
% =====================================================================

\subsubsection{Gauss quadrature for BFS elements}

The BFS element requires numerical integration of the element matrices.
For polynomial integrands, Gauss--Legendre quadrature provides high
accuracy with a small number of points.

\textbf{Volume integrals (stiffness matrices, internal force, mass):}
For integrands involving products of BFS shape functions and their
derivatives, a $3\times 3$ Gauss quadrature in the parent domain
$(\xi,\eta)\in[-1,1]^{2}$ is adopted in the present implementation
and verified numerically for the benchmark problem. The polynomial
degree of the integrands is at most $5$ in each direction for the
linear terms, but the nonlinear formulation involves configuration-dependent
quantities that are evaluated at the Gauss points, so exactness in the
classical polynomial sense is not claimed for the full nonlinear system.

\begin{equation}
\int_{\mathcal{B}_{0}^{e}} g(\mathbf{X})\,\mathrm{d}V_{0}
\approx \frac{h_{1}h_{2}}{4}\sum_{p=1}^{3}\sum_{q=1}^{3}
w_{p}\,w_{q}\,g(\xi_{p},\eta_{q})\,J_{\mathrm{map}},
\label{s7:gauss_volume}
\end{equation}

where:
\begin{itemize}
\item $(\xi_{p},\eta_{q})$ are the Gauss points:
  $\xi_{1,2,3} = -\sqrt{3/5},\;0,\;\sqrt{3/5}$
  (similarly for $\eta$).
\item $w_{1,2,3} = 5/9,\;8/9,\;5/9$ are the Gauss weights.
\item $J_{\mathrm{map}}=1$ for the standard BFS mapping from parent
  to physical domain.
\item $h_{1},h_{2}$ are the element dimensions.
\end{itemize}

\textbf{Surface integrals (external force):}
For boundary traction integrals, 1D Gauss quadrature along each
element edge is used:

\begin{equation}
\int_{\partial\mathcal{B}_{0t}^{e}} g(\mathbf{X})\,\mathrm{d}A_{0}
\approx \frac{h_{e}}{2}\sum_{p=1}^{n_{\mathrm{surf}}}
w_{p}\,g(\xi_{p})\,
\end{equation}

where $h_{e}$ is the edge length and $n_{\mathrm{surf}}=3$ is typically
sufficient.

\subsubsection{Summary of integration requirements}

\begin{longtable}[c]{>{\raggedright\arraybackslash}p{0.30\linewidth}
>{\raggedright\arraybackslash}p{0.25\linewidth}
>{\raggedright\arraybackslash}p{0.45\linewidth}}
\caption{Numerical integration requirements.}
\label{tab:integration}\\
\toprule
Integral & Quadrature rule & Maximum polynomial degree \\
\midrule
\endfirsthead
\toprule
Integral & Quadrature rule & Maximum polynomial degree \\
\midrule
\endhead
$\mathbf{K}_{\mathrm{mat,cl}}^{e}$ & $3\times 3$ Gauss & adopted for $C^{1}$ BFS \\
$\mathbf{K}_{\mathrm{geo,cl}}^{e}$ & $3\times 3$ Gauss & adopted \\
$\mathbf{K}_{\mathrm{mat,gr}}^{e}$ & $3\times 3$ Gauss & adopted \\
$\mathbf{K}_{\mathrm{geo,gr}}^{e}$ & $3\times 3$ Gauss & adopted \\
$\mathbf{r}_{\mathrm{int}}^{e}$ & $3\times 3$ Gauss & adopted \\
$\mathbf{f}_{\mathrm{ext}}^{e}$ (volume) & $3\times 3$ Gauss & adopted \\
$\mathbf{f}_{\mathrm{ext}}^{e}$ (surface) & 3-point 1D Gauss & adopted \\
$\mathbf{M}^{e}$ & $3\times 3$ Gauss & adopted \\
\bottomrule
\end{longtable}

% =====================================================================
\subsection{Assembly operator}
\label{ssec:assembly}
% =====================================================================

\subsubsection{Element-to-global assembly}

The assembly operator $\mathbb{A}$ maps element quantities to the
global system. Let $\mathcal{I}^{e}$ be the connectivity array for
element $e$, mapping the element's 32 local DOFs to the global DOF
indices.

\begin{equation}
\boxed{
\mathbf{R}_{I_{\mathcal{I}^{e}(a)}} = \sum_{e=1}^{N_{\mathrm{elem}}}
  \mathbf{r}^{e}_{a},
\qquad
\mathbf{K}_{I_{\mathcal{I}^{e}(a)},\,I_{\mathcal{I}^{e}(b)}}
= \sum_{e=1}^{N_{\mathrm{elem}}}
  \mathbf{K}_{T,a b}^{e},
}
\label{s7:assembly}
\end{equation}

for $a,b\in\{1,\ldots,32\}$, where $I_{\mathcal{I}^{e}(a)}$ is the
global DOF index corresponding to element $e$'s local DOF $a$.

\textbf{Assembly algorithm:}
\begin{enumerate}
\item Initialize $\mathbf{R}=\mathbf{0}\in\mathbb{R}^{N_{\mathrm{dof}}}$
  and $\mathbf{K}_{T}=\mathbf{0}\in\mathbb{R}^{N_{\mathrm{dof}}\times N_{\mathrm{dof}}}$.
\item For each element $e=1,\ldots,N_{\mathrm{elem}}$:
  \begin{enumerate}
  \item Compute $\mathbf{r}^{e}\in\mathbb{R}^{32}$ and
    $\mathbf{K}_{T}^{e}\in\mathbb{R}^{32\times 32}$.
  \item Scatter-add: $\mathbf{R}(\mathcal{I}^{e}) \mathrel{+}= \mathbf{r}^{e}$.
  \item Scatter-add: $\mathbf{K}_{T}(\mathcal{I}^{e},\mathcal{I}^{e})
    \mathrel{+}= \mathbf{K}_{T}^{e}$.
  \end{enumerate}
\end{enumerate}

\subsubsection{Global system}

After assembly:

\begin{equation}
\boxed{
\mathbf{R}(\mathbf{D})
= \mathbb{A}_{e=1}^{N_{\mathrm{elem}}}
  \mathbf{r}^{e}(\mathbf{d}^{e})
= \mathbf{R}_{\mathrm{int}}(\mathbf{D})
- \mathbf{F}_{\mathrm{ext}}
+ \mathbf{F}_{\mathrm{iner}}
= \mathbf{0},
}
\label{s7:global_residual}
\end{equation}

\begin{equation}
\boxed{
\mathbf{K}_{T}(\mathbf{D})
= \mathbb{A}_{e=1}^{N_{\mathrm{elem}}}
  \mathbf{K}_{T}^{e}(\mathbf{d}^{e})
\in \mathbb{R}^{N_{\mathrm{dof}}\times N_{\mathrm{dof}}}.
}
\label{s7:global_tangent}
\end{equation}

The global system is symmetric and sparse (banded structure for
structured BFS meshes).

% =====================================================================
\subsection{Newton--Raphson iterative algorithm}
\label{ssec:newton}
% =====================================================================

\subsubsection{Algorithm statement}

The Newton--Raphson method solves the nonlinear system
$\mathbf{R}(\mathbf{D})=\mathbf{0}$ iteratively. At each load/time
step with prescribed external loads $\mathbf{F}_{\mathrm{ext}}^{n+1}$:

\begin{enumerate}
\item \textbf{Initialize:}
  Set $\mathbf{D}^{(0)}=\mathbf{D}^{n}$ (solution at previous step).

\item \textbf{Iteration loop:} For $k=0,1,2,\ldots,k_{\max}$:
  \begin{enumerate}
  \item \textbf{Compute residual:}
    Evaluate $\mathbf{R}^{(k)}=\mathbf{R}_{\mathrm{int}}(\mathbf{D}^{(k)})
    - \mathbf{F}_{\mathrm{ext}}^{n+1}
    + \mathbf{M}\,\ddot{\mathbf{D}}^{(k)}$.

  \item \textbf{Check convergence:}
    If convergence criteria are satisfied (see below), set
    $\mathbf{D}^{n+1}=\mathbf{D}^{(k)}$ and proceed to next step.

  \item \textbf{Compute tangent:}
    Evaluate $\mathbf{K}_{T}^{(k)}=\mathbf{K}_{T}(\mathbf{D}^{(k)})$.

  \item \textbf{Solve linear system:}
    Solve $\mathbf{K}_{T}^{(k)}\,\Delta\mathbf{D}^{(k)}
    = -\mathbf{R}^{(k)}$ for $\Delta\mathbf{D}^{(k)}$.

  \item \textbf{Update:}
    Set $\mathbf{D}^{(k+1)}=\mathbf{D}^{(k)}+\Delta\mathbf{D}^{(k)}$.
  \end{enumerate}

\item \textbf{Failure:}
  If $k>k_{\max}$ without convergence, reduce load increment and
  retry (see \S\ref{ssec:stepping}).
\end{enumerate}

\subsubsection{Convergence criteria}
\label{ssec:convergence}

Convergence is assessed using three complementary norms
\cite{Bathe2006,Belytschko2000,ZienkiewiczTaylor2000}:

\textbf{1. Residual norm (force convergence):}
\begin{equation}
\eta_{R}^{(k)}
:= \frac{\|\mathbf{R}^{(k)}\|_{2}}{\|\mathbf{R}^{(0)}\|_{2}}
< \mathrm{tol}_{R}.
\label{s7:conv_residual}
\end{equation}

\textbf{2. Increment norm (displacement convergence):}
\begin{equation}
\eta_{D}^{(k)}
:= \frac{\|\Delta\mathbf{D}^{(k)}\|_{2}}{\|\mathbf{D}^{(k+1)}\|_{2}}
< \mathrm{tol}_{D}.
\label{s7:conv_increment}
\end{equation}

\textbf{3. Energy norm (work convergence):}
\begin{equation}
\eta_{E}^{(k)}
:= \frac{\bigl|\Delta\mathbf{D}^{(k)T}\,\mathbf{R}^{(k)}\bigr|}
        {\bigl|\Delta\mathbf{D}^{(0)T}\,\mathbf{R}^{(0)}\bigr|}
< \mathrm{tol}_{E}.
\label{s7:conv_energy}
\end{equation}

\textbf{Convergence test:} The iteration is declared converged when
\emph{all three} criteria are simultaneously satisfied:

\begin{equation}
\eta_{R}^{(k)}<\mathrm{tol}_{R}
\quad\wedge\quad
\eta_{D}^{(k)}<\mathrm{tol}_{D}
\quad\wedge\quad
\eta_{E}^{(k)}<\mathrm{tol}_{E}.
\label{s7:conv_test}
\end{equation}

\textbf{Recommended tolerances:}
\begin{equation}
\mathrm{tol}_{R}=\mathrm{tol}_{D}=\mathrm{tol}_{E}=10^{-8}.
\label{s7:tol_recommended}
\end{equation}

\textbf{Maximum iterations:}
\begin{equation}
k_{\max}=20.
\label{s7:kmax}
\end{equation}

Due to the consistent linearization, quadratic convergence is expected
near the solution, typically requiring $k=4$--$8$ iterations for
moderate load increments.

\subsubsection{Load/time stepping strategy}
\label{ssec:stepping}

The external loads and/or time are advanced incrementally. Let the
total analysis be divided into $N_{\mathrm{steps}}$ load/time steps,
with step $n$ going from $t^{n}$ to $t^{n+1}$.

\textbf{Standard (fixed) stepping:}
\begin{equation}
\mathbf{F}_{\mathrm{ext}}^{n}
= \frac{n}{N_{\mathrm{steps}}}\,\mathbf{F}_{\mathrm{ext}}^{\mathrm{total}},
\qquad
n=0,1,\ldots,N_{\mathrm{steps}}.
\label{s7:fixed_step}
\end{equation}

\textbf{Adaptive stepping:}
If convergence fails at step $n$ after $k_{\max}$ iterations,
the load increment is reduced by a factor $\alpha\in(0,1)$:
\begin{equation}
\Delta\mathbf{F}_{\mathrm{ext}}^{n}
\leftarrow \alpha\,\Delta\mathbf{F}_{\mathrm{ext}}^{n}.
\label{s7:adaptive_step}
\end{equation}
The value $\alpha=0.5$ is a typical value adopted in the present
implementation, consistent with standard nonlinear FEM practice
\cite{Bathe2006,Belytschko2000}.

After successful convergence, the increment may be increased:
\begin{equation}
\Delta\mathbf{F}_{\mathrm{ext}}^{n+1}
\leftarrow \beta\,\Delta\mathbf{F}_{\mathrm{ext}}^{n},
\label{s7:increase_step}
\end{equation}

provided the number of iterations $k$ satisfies $k<k_{\mathrm{target}}$
(recommended $k_{\mathrm{target}}=6$). The value $\beta=1.2$ is a
typical value adopted in the present implementation, consistent with
standard nonlinear FEM practice \cite{Bathe2006,Belytschko2000}.

% =====================================================================
\subsection{Complete solution algorithm}
\label{ssec:algorithm}
% =====================================================================

\subsubsection{Overall solution procedure}

\begin{algorithm}[H]
\caption{Complete nonlinear FEM solver}
\label{alg:complete}
\begin{algorithmic}[1]
\REQUIRE Mesh $\mathcal{M}$, material parameters $\lambda,\mu,a_{1},\ldots,a_{5}$
\REQUIRE Load schedule $\mathbf{F}_{\mathrm{ext}}^{0},\ldots,\mathbf{F}_{\mathrm{ext}}^{N_{\mathrm{steps}}}$
\REQUIRE Convergence tolerances $\mathrm{tol}_{R},\mathrm{tol}_{D},\mathrm{tol}_{E}$
\REQUIRE Maximum iterations $k_{\max}$
\ENSURE Solution history $\{\mathbf{D}^{n}\}_{n=0}^{N_{\mathrm{steps}}}$

\STATE Initialize: $\mathbf{D}^{0}=\mathbf{0}$, apply essential BCs
\FOR{$n=0,1,\ldots,N_{\mathrm{steps}}-1$}
  \STATE Set target load: $\mathbf{F}_{\mathrm{target}}=\mathbf{F}_{\mathrm{ext}}^{n+1}$
  \STATE $\mathbf{D}^{(0)}\leftarrow\mathbf{D}^{n}$
  \STATE $k\leftarrow 0$

  \WHILE{$k<k_{\max}$}
    \STATE \textit{// Element loop}
    \STATE $\mathbf{R}\leftarrow\mathbf{0}$, $\mathbf{K}_{T}\leftarrow\mathbf{0}$
    \FOR{each element $e$}
      \STATE Compute $\mathbf{d}^{e}$ from $\mathbf{D}^{(k)}$
      \STATE Evaluate shape functions at Gauss points
      \STATE Compute $F_{iJ}^{h}$, $E_{IJ}^{h}$, $\mathcal{G}_{IJK}^{h}$
      \STATE Evaluate stresses: $S_{IJ}^{h}$, $P_{iJ}^{h}$, $\Sigma_{IJK}^{h}$
      \STATE Assemble $\mathbf{r}_{\mathrm{int}}^{e}$ (Eq.~\eqref{s7:r_int})
      \STATE Assemble $\mathbf{K}_{\mathrm{mat,cl}}^{e}$ (Eq.~\eqref{s7:K_mat_cl})
      \STATE Assemble $\mathbf{K}_{\mathrm{geo,cl}}^{e}$ (Eq.~\eqref{s7:K_geo_cl})
      \STATE Assemble $\mathbf{K}_{\mathrm{mat,gr}}^{e}$ (Eq.~\eqref{s7:K_mat_gr})
      \STATE Assemble $\mathbf{K}_{\mathrm{geo,gr}}^{e}$ (Eq.~\eqref{s7:K_geo_gr})
      \STATE $\mathbf{K}_{T}^{e}\leftarrow\mathbf{K}_{\mathrm{mat,cl}}^{e}+\mathbf{K}_{\mathrm{geo,cl}}^{e}+\mathbf{K}_{\mathrm{mat,gr}}^{e}+\mathbf{K}_{\mathrm{geo,gr}}^{e}$
      \STATE Scatter-add to global: $\mathbf{R}\mathrel{+}=\mathbf{r}_{\mathrm{int}}^{e}$, $\mathbf{K}_{T}\mathrel{+}=\mathbf{K}_{T}^{e}$
    \ENDFOR
    \STATE \textit{// Apply loads and BCs}
    \STATE $\mathbf{R}\leftarrow\mathbf{R}-\mathbf{F}_{\mathrm{target}}$
    \STATE Apply essential BCs: zero rows/columns in $\mathbf{K}_{T}$, zero entries in $\mathbf{R}$

    \STATE \textit{// Check convergence}
    \IF{$\|\mathbf{R}\|/\|\mathbf{R}^{(0)}\|<\mathrm{tol}_{R}$ AND $\|\Delta\mathbf{D}\|/\|\mathbf{D}\|<\mathrm{tol}_{D}$ AND energy norm $<\mathrm{tol}_{E}$}
      \STATE $\mathbf{D}^{n+1}\leftarrow\mathbf{D}^{(k)}$
      \STATE \textbf{break} (converged)
    \ENDIF

    \STATE \textit{// Solve and update}
    \STATE Solve $\mathbf{K}_{T}\,\Delta\mathbf{D}=-\mathbf{R}$
    \STATE $\mathbf{D}^{(k+1)}\leftarrow\mathbf{D}^{(k)}+\Delta\mathbf{D}$
    \STATE $k\leftarrow k+1$
  \ENDWHILE

  \IF{$k\geq k_{\max}$}
    \STATE Reduce load increment (adaptive stepping)
    \STATE Retry current step
  \ENDIF
\ENDFOR
\end{algorithmic}
\end{algorithm}

% =====================================================================
\subsection{Pseudocode of the complete nonlinear FEM solver}
\label{ssec:pseudocode}
% =====================================================================

The following pseudocode specifies the complete solver at the module
level, suitable for direct translation to MATLAB.

\subsubsection{Main driver}
\begin{verbatim}
FUNCTION MainSolver(mesh, material, loads, BCs, tolerances)
    D = InitializeDOFs(mesh)
    D_history = {D}

    FOR n = 1 TO N_steps
        F_target = loads(n)
        D = NewtonRaphson(D, mesh, material, F_target, BCs, tolerances)
        APPEND D_history WITH D
    END FOR

    RETURN D_history
END FUNCTION
\end{verbatim}

\subsubsection{Newton--Raphson module}
\begin{verbatim}
FUNCTION NewtonRaphson(D_old, mesh, material, F_target, BCs, tol)
    D = D_old
    R0 = ComputeResidual(D, mesh, material, F_target, BCs)
    k = 0

    WHILE k < k_max
        [R, K_T] = AssembleSystem(D, mesh, material, F_target, BCs)

        IF CheckConvergence(R, R0, D, tol)
            RETURN D
        END IF

        dD = SolveLinearSystem(K_T, -R)
        D = D + dD
        k = k + 1
    END WHILE

    ERROR "Newton-Raphson did not converge"
END FUNCTION
\end{verbatim}

\subsubsection{Element computation module}
\begin{verbatim}
FUNCTION [r_e, K_e] = ComputeElement(d_e, element, material)
    [gp, gw] = GaussQuadrature(3, 3)

    r_int = zeros(32, 1)
    K_mat_cl = zeros(32, 32)
    K_geo_cl = zeros(32, 32)
    K_mat_gr = zeros(32, 32)
    K_geo_gr = zeros(32, 32)

    FOR each Gauss point (xi, eta, w)
        [N, B, G] = EvaluateBFS(xi, eta, element)
        [F_h, E_h, G_h] = ComputeKinematics(N, B, G, d_e)
        [S, P, Sigma] = ComputeStresses(E_h, G_h, material)
        [B_deltaE, B_deltaG] = ComputeBmatrices(F_h, G_h, B, G, d_e)

        r_int = r_int + w * (B' * P + B_deltaG' * Sigma)
        K_mat_cl = K_mat_cl + w * B_deltaE' * C_4th * B_deltaE
        K_geo_cl = K_geo_cl + w * B' * S_2nd * B
        K_mat_gr = K_mat_gr + w * B_deltaG' * D_6th * B_deltaG
        K_geo_gr = K_geo_gr + w * G_tensor * Sigma * B_tensor
    END FOR

    r_e = r_int - f_ext_e
    K_e = K_mat_cl + K_geo_cl + K_mat_gr + K_geo_gr
END FUNCTION
\end{verbatim}

% =====================================================================
\subsection{MATLAB module dependency diagram}
\label{ssec:modules}
% =====================================================================

The MATLAB implementation is organized into the following modules:

\begin{verbatim}
MainSolver.m
|
+-- Mesh.m                      % Mesh data structure, connectivity
|
+-- Material.m                  % Material parameters (lambda, mu, D_6th)
|
+-- BoundaryConditions.m        % Essential BCs, load schedule
|
+-- NewtonRaphson.m             % Iterative solver
|   |
|   +-- ComputeResidual.m       % Residual evaluation
|   +-- AssembleSystem.m        % Global assembly
|   |   |
|   |   +-- ComputeElement.m    % Element matrices and vectors
|   |   |   |
|   |   |   +-- EvaluateBFS.m   % Shape functions, derivatives
|   |   |   +-- GaussQuadrature.m % Gauss points and weights
|   |   |   +-- ComputeKinematics.m % F, E, G from DOFs
|   |   |   +-- ComputeStresses.m % S, P, Sigma from kinematics
|   |   |   +-- ComputeBmatrices.m % B_deltaE, B_deltaG
|   |   |
|   |   +-- ApplyBCs.m          % Apply essential BCs to K, R
|   |
|   +-- CheckConvergence.m      % Convergence criteria
|   +-- SolveLinearSystem.m     % Linear solver (sparse direct)
|
+-- PostProcess.m               % Results extraction, visualization
|   +-- ComputeDerivedFields.m  % Stress, strain from converged DOFs
|   +-- BenchmarkComparison.m   % Comparison with 2022 benchmark
|
+-- Utilities/
    +-- IndexMapping.m          % DOF numbering, assembly indices
    +-- TensorOperations.m      % 4th/6th-order tensor operations
    +-- IsotropicTensors.m      % C_4th, D_6th construction
\end{verbatim}

\textbf{Module descriptions:}

\begin{longtable}[c]{>{\raggedright\arraybackslash}p{0.25\linewidth}
>{\raggedright\arraybackslash}p{0.75\linewidth}}
\caption{MATLAB module descriptions.}
\label{tab:modules}\\
\toprule
Module & Description \\
\midrule
\endfirsthead
\toprule
Module & Description \\
\midrule
\endhead
\texttt{MainSolver} & Top-level driver. Reads input, calls Newton--Raphson, stores results. \\
\texttt{NewtonRaphson} & Iterative solver. Calls \texttt{AssembleSystem}, checks convergence, solves linear system. \\
\texttt{ComputeResidual} & Evaluates the global residual vector (called by \texttt{NewtonRaphson} for convergence normalization). \\
\texttt{AssembleSystem} & Global assembly. Loops over elements, calls \texttt{ComputeElement}, scatters to global arrays. \\
\texttt{ComputeElement} & Element-level computation. Calls sub-modules for shape functions, kinematics, stresses. \\
\texttt{EvaluateBFS} & Evaluates BFS bicubic Hermite shape functions and derivatives at parent coordinates. \\
\texttt{GaussQuadrature} & Returns Gauss points and weights for $3\times 3$ quadrature. \\
\texttt{ComputeKinematics} & Computes $F_{iJ}^{h}$, $E_{IJ}^{h}$, $\mathcal{G}_{IJK}^{h}$ from element DOFs using \eqref{s7:assembled_interp}. \\
\texttt{ComputeStresses} & Evaluates $S_{IJ}^{h}$, $P_{iJ}^{h}$, $\Sigma_{IJK}^{h}$ from kinematics using \eqref{s7:P_h}--\eqref{s7:Sig_h}. \\
\texttt{ComputeBmatrices} & Assembles $\mathbf{B}_{\delta E,IJ}^{*}$ and $\mathbf{B}_{\delta\mathcal{G},IJK}^{*}$ using \eqref{s7:B_deltaE_assembled}--\eqref{s7:B_deltaG_assembled}. \\
\texttt{IsotropicTensors} & Constructs $\mathbb{C}_{IJKL}$ and $\mathbb{D}_{IJKPQR}$ from Lam\'e parameters and gradient constants. \\
\bottomrule
\end{longtable}

% =====================================================================
\subsection{Small-strain limit and 2022 benchmark recovery}
\label{ssec:recovery}
% =====================================================================

In the infinitesimal-strain limit ($\|\nabla_{0}\mathbf{u}\|\to 0$):

\begin{align}
F_{iJ}^{h} &\to \delta_{iJ},
&
E_{IJ}^{h} &\to \varepsilon_{IJ}^{h},
&
\mathcal{G}_{IJK}^{h} &\to \varepsilon_{IJ,K}^{h},
\label{s7:lim_kin}\\
S_{IJ}^{h} &\to \sigma_{ij},
&
P_{iJ}^{h} &\to \sigma_{ij},
&
\Sigma_{IJK}^{h} &\to h_{ijk},
\label{s7:lim_stress}
\end{align}

and the strain-displacement matrices reduce to their linear forms:
\begin{equation}
\mathbf{B}_{\delta E,IJ}^{*} \to \mathbf{B}_{\varepsilon,ij}^{*},
\qquad
\mathbf{B}_{\delta\mathcal{G},IJK}^{*} \to \mathbf{B}_{\nabla\varepsilon,ijk}^{*}.
\end{equation}

\subsubsection{Recovery of element tangent}

\textbf{Classical material tangent:}
\begin{equation}
\mathbf{K}_{\mathrm{mat,cl}}^{e}
\to \int_{\mathcal{B}_{0}^{e}}
  \sum_{i,j,k,l}
  \mathbb{C}_{ijkl}\,
  \mathbf{B}_{\varepsilon,kl}^{*}\,
  (\mathbf{B}_{\varepsilon,ij}^{*})^{T}\,
  \mathrm{d}V_{0}.
\end{equation}

This is exactly the classical linear elastic stiffness matrix. $\checkmark$

\textbf{Classical geometric tangent:}
\begin{equation}
\mathbf{K}_{\mathrm{geo,cl}}^{e}
\to \int_{\mathcal{B}_{0}^{e}}
  \sum_{i,I,J}
  \underbrace{S_{IJ}^{h}}_{\to\,0}\,
  \mathbf{B}_{iJ}^{*}\,
  (\mathbf{B}_{iI}^{*})^{T}\,
  \mathrm{d}V_{0}
= \mathbf{0}.
\end{equation}

Vanishes because $S_{IJ}=\mathbf{0}$ in the linear theory. $\checkmark$

\textbf{Gradient material tangent:}
\begin{equation}
\mathbf{K}_{\mathrm{mat,gr}}^{e}
\to \int_{\mathcal{B}_{0}^{e}}
  \sum_{i,j,k,p,q,r}
  \mathbb{D}_{ijkpqr}\,
  \mathbf{B}_{\nabla\varepsilon,pqr}^{*}\,
  (\mathbf{B}_{\nabla\varepsilon,ijk}^{*})^{T}\,
  \mathrm{d}V_{0}.
\end{equation}

This is exactly the Mindlin Form-II gradient stiffness matrix of the
2022 benchmark. $\checkmark$

\textbf{Gradient geometric tangent:}
\begin{equation}
\mathbf{K}_{\mathrm{geo,gr}}^{e}
\to \int_{\mathcal{B}_{0}^{e}}
  \sum_{i,I,J,K}
  \underbrace{\Sigma_{IJK}^{h}}_{\to\,0}\,
  \bigl[\cdots\bigr]\,
  \mathrm{d}V_{0}
= \mathbf{0}.
\end{equation}

Vanishes because $\Sigma_{IJK}=\mathbf{0}$ in the linear theory. $\checkmark$

\subsubsection{Recovery of total tangent}

The total element tangent reduces to:

\begin{equation}
\boxed{
\mathbf{K}_{T}^{e}
\to \mathbf{K}_{\mathrm{mat,cl}}^{e}
+ \mathbf{K}_{\mathrm{mat,gr}}^{e}
\in \mathbb{R}^{32\times 32},
}
\label{s7:K_2022}
\end{equation}

which is \emph{exactly} the element stiffness matrix of the validated
2022 Mindlin Form-II formulation. The geometric contributions vanish
identically. $\checkmark$

\textbf{Mandatory recovery rule: SATISFIED at the element matrix level.}

% =====================================================================
\subsection{Summary of Section~7}
\label{ssec:summary}
% =====================================================================

\begin{enumerate}
\item The element tangent matrix has four symmetric contributions:
  classical material, classical geometric, gradient material, and
  gradient geometric stiffness.

\item The element force vectors (internal, external, inertia) follow
  directly from the discrete weak form of Section~6.

\item The $3\times 3$ Gauss quadrature is adopted for numerical
  integration of the BFS element matrices and verified for the
  benchmark problem.

\item The assembly operator maps $32\times 32$ element quantities to
  the global $N_{\mathrm{dof}}\times N_{\mathrm{dof}}$ system.

\item The Newton--Raphson algorithm with three convergence criteria
  (residual, increment, energy) ensures robust solution of the
  nonlinear system.

\item Adaptive load stepping handles difficult load paths.

\item The complete solution algorithm is specified in both algorithmic
  and pseudocode form, ready for MATLAB implementation.

\item The MATLAB module hierarchy is clearly defined.

\item In the small-strain limit, all element matrices reduce exactly
  to the validated 2022 Mindlin Form-II benchmark formulation.
\end{enumerate}

\textbf{Section~7 is complete. This section provides the complete
mathematical specification required for the MATLAB implementation.}

% =====================================================================
% SECTION 8 — VERIFICATION FRAMEWORK
% Master Version 1.2 — Frozen Sections 1-7 unchanged.
% Verification framework for 2022 Mindlin Form-II benchmark recovery.
% Uses ONLY frozen notation from Sections 1-7.
% Compile with pdflatex + bibtex (or Overleaf).
% =====================================================================

% =====================================================================
\section{Verification Framework}
\label{sec:verification}
% =====================================================================

\subsection{Purpose and scope}
\label{ssec:s8_purpose}

This section establishes the verification framework for the finite-strain
$C^{1}$ Bogner--Fox--Schmit (BFS) Mindlin Form-II strain-gradient
formulation developed in Sections~1--7. The objective is to demonstrate
that the proposed finite element formulation is mathematically and
computationally correct by recovering the validated 2022 Mindlin
Form-II benchmark in the infinitesimal-strain limit.

\textbf{Note:} This section provides the verification framework. Numerical values, convergence rates, and verification figures will be inserted after MATLAB computations are completed.

\textbf{Verification vs.\ validation distinction:}

\begin{itemize}
\item \textbf{Verification} asks: ``Are we solving the equations
  correctly?'' It confirms that the numerical implementation correctly
  solves the mathematical model derived in Sections~1--7. Verification
  compares numerical results against analytical solutions, benchmark
  problems with known solutions, or code-to-code comparisons
  \cite{OberkampfTrucano2002,ASME_Verification_2012}.

\item \textbf{Validation} asks: ``Are we solving the right equations?''
  It confirms that the mathematical model correctly represents the
  physical system. Validation requires comparison against experimental
  data \cite{OberkampfTrucano2002,ASME_Verification_2012}.
\end{itemize}

This section performs \textbf{verification}, not validation. We compare
the finite-strain formulation against the validated 2022 benchmark
\cite{TorabiNiiranen2022} in the small-strain limit. If the numerical
results converge to the benchmark solution, we verify that the
implementation correctly solves the derived equations. Validation
against experimental data is outside the scope of this work and will be
addressed in future publications.

\textbf{Scope:}
\begin{enumerate}
\item Describe the 2022 benchmark problem
\item Define material properties and geometry
\item Specify boundary and loading conditions
\item Establish error norms and convergence criteria
\item Provide mesh convergence study framework
\item Present verification tables and figures
\item Confirm exact benchmark recovery
\end{enumerate}

\subsection{Description of the 2022 benchmark problem}
\label{ssec:benchmark}

The validated 2022 benchmark \cite{TorabiNiiranen2022} considers a
two-dimensional strain-gradient elastic body in the small-strain regime.
The benchmark provides analytical or semi-analytical solutions for
specific loading and boundary conditions, allowing rigorous verification
of finite element formulations.

\textbf{Key characteristics of the 2022 benchmark:}
\begin{itemize}
\item \textbf{Theory:} Mindlin Form-II first strain-gradient elasticity
  \cite{Mindlin1964,MindlinEshel1968}
\item \textbf{Strain measure:} Infinitesimal strain tensor
  $\varepsilon_{ij} = \tfrac12(u_{i,j}+u_{j,i})$
\item \textbf{Higher-order kinematic variable:} Strain gradient tensor
  $\varepsilon_{ij,k} = \partial\varepsilon_{ij}/\partial x_{k}$
\item \textbf{Constitutive law:} Isotropic strain-gradient elasticity
  with five gradient constants $a_{1},\ldots,a_{5}$
\item \textbf{Governing equations:} Fourth-order partial differential
  equations in displacements
\item \textbf{Boundary conditions:} Essential (displacement + normal
  derivative) and natural (traction + double traction)
\end{itemize}

\textbf{Recovery rule:} The finite-strain formulation developed in
Sections~1--7 must reduce \emph{exactly} to the 2022 benchmark when
$\|\nabla_{0}\mathbf{u}\|\to 0$. Specifically:
\begin{align}
F_{iJ} &\to \delta_{iJ},
\qquad
E_{IJ} \to \varepsilon_{IJ},
\qquad
\mathcal{G}_{IJK} \to \varepsilon_{IJ,K},
\label{s8:limits}
\\
P_{iJ} &\to \sigma_{ij},
\qquad
\Sigma_{IJK} \to h_{ijk},
\label{s8:stress_limits}
\\
\mathbf{K}_{T}^{e} &\to \mathbf{K}_{\mathrm{mat,cl}}^{e}
+ \mathbf{K}_{\mathrm{mat,gr}}^{e}.
\label{s8:tangent_limit}
\end{align}

\subsection{Material properties}
\label{ssec:material}

The material is isotropic with the following constitutive parameters:

\begin{longtable}[c]{>{\raggedright\arraybackslash}p{0.15\linewidth}
>{\raggedright\arraybackslash}p{0.10\linewidth}
>{\raggedright\arraybackslash}p{0.15\linewidth}
>{\raggedright\arraybackslash}p{0.60\linewidth}}
\caption{Material properties for the 2022 benchmark.}
\label{tab:material}\\
\toprule
Parameter & Symbol & Value & Physical meaning \\
\midrule
\endfirsthead
\toprule
Parameter & Symbol & Value & Physical meaning \\
\midrule
\endhead
Young's modulus & $E$ & taken directly from the validated 2022 benchmark without modification & Classical elastic modulus \\
Poisson's ratio & $\nu$ & taken directly from the validated 2022 benchmark without modification & Classical elastic ratio \\
Lam\'e parameter & $\lambda$ & computed from the classical elastic constants & $\lambda = E\nu/[(1+\nu)(1-2\nu)]$ \\
Shear modulus & $\mu$ & computed from the classical elastic constants & $\mu = E/[2(1+\nu)]$ \\
Gradient constant & $a_{1}$ & taken directly from the validated 2022 benchmark without modification & Strain-gradient coefficient \\
Gradient constant & $a_{2}$ & taken directly from the validated 2022 benchmark without modification & Strain-gradient coefficient \\
Gradient constant & $a_{3}$ & taken directly from the validated 2022 benchmark without modification & Strain-gradient coefficient \\
Gradient constant & $a_{4}$ & taken directly from the validated 2022 benchmark without modification & Strain-gradient coefficient \\
Gradient constant & $a_{5}$ & taken directly from the validated 2022 benchmark without modification & Strain-gradient coefficient \\
Internal length scale & $\ell$ & computed from the gradient constants and shear modulus & Characteristic length \\
Mass density & $\rho_{0}$ & taken directly from the validated 2022 benchmark without modification & Reference density \\
\bottomrule
\end{longtable}

\textbf{Internal length scale:} The internal length scale $\ell$ is
defined as
\begin{equation}
\ell = \sqrt{\frac{\bar{a}}{\mu}},
\qquad
\bar{a} = \frac{1}{5}\sum_{i=1}^{5}a_{i},
\label{s8:length_scale}
\end{equation}
where $\bar{a}$ is the average gradient constant. The internal length
scale controls the size-dependent effects: as $\ell\to 0$, the
formulation reduces to classical elasticity.

\textbf{Gradient elasticity tensor:} The sixth-order gradient
elasticity tensor $\mathbb{D}_{IJKPQR}$ is constructed from the five
constants $a_{1},\ldots,a_{5}$ using the isotropic representation
(Section~2, Eq.~\eqref{eq:sec2_gradient_energy_invariants}):
\begin{equation}
\Psi_{\mathrm{g}} = a_{1}\,\varepsilon_{ii,k}\,\varepsilon_{jj,k}
+ a_{2}\,\varepsilon_{ij,j}\,\varepsilon_{kk,i}
+ a_{3}\,\varepsilon_{ij,i}\,\varepsilon_{kk,j}
+ a_{4}\,\varepsilon_{ij,k}\,\varepsilon_{ij,k}
+ a_{5}\,\varepsilon_{ij,k}\,\varepsilon_{kj,i}.
\label{s8:gradient_energy}
\end{equation}

\subsection{Geometry and boundary conditions}
\label{ssec:geometry}

The benchmark considers a two-dimensional rectangular domain
$\Omega = [0,L_{1}]\times[0,L_{2}]$ in the reference configuration.

\textbf{Geometry:}
\begin{itemize}
\item Length in $X_{1}$-direction: $L_{1}$ = taken directly from the validated 2022 benchmark without modification
\item Length in $X_{2}$-direction: $L_{2}$ = taken directly from the validated 2022 benchmark without modification
\item Thickness (plane strain): $t$ = taken directly from the validated 2022 benchmark without modification
\end{itemize}

\textbf{Boundary conditions:} The boundary $\partial\Omega$ is
partitioned into essential (Dirichlet) and natural (Neumann) parts:
\begin{equation}
\partial\Omega = \partial\Omega_{u} \cup \partial\Omega_{t},
\qquad
\partial\Omega_{u} \cap \partial\Omega_{t} = \emptyset.
\end{equation}

\textbf{Essential boundary conditions on $\partial\Omega_{u}$:}
\begin{equation}
u_{i} = \bar{u}_{i},
\qquad
\frac{\partial u_{i}}{\partial n} = \bar{g}_{i}
\qquad\text{on }\partial\Omega_{u},
\label{s8:essBC}
\end{equation}
where $n$ is the outward unit normal.

\textbf{Natural boundary conditions on $\partial\Omega_{t}$:}
\begin{equation}
(\sigma_{ij}-h_{ijk,k})n_{j} = \bar{t}_{i},
\qquad
h_{ijk}n_{j}n_{k} = \bar{d}_{i}
\qquad\text{on }\partial\Omega_{t},
\label{s8:natBC}
\end{equation}
where $\bar{t}_{i}$ is the prescribed traction and $\bar{d}_{i}$ is the
prescribed double traction.

\textbf{Specific boundary conditions for the benchmark:}
The detailed boundary conditions are specified by the 2022 benchmark and will be implemented accordingly.

\subsection{Loading conditions}
\label{ssec:loading}

The body is subjected to:
\begin{itemize}
\item \textbf{Body force:} $b_{i}(\bx)$ = specified by the 2022 benchmark
\item \textbf{Traction:} $\bar{t}_{i}(\bx)$ = specified by the 2022 benchmark
\item \textbf{Double traction:} $\bar{d}_{i}(\bx)$ = specified by the 2022 benchmark
\end{itemize}

For static problems, the inertia term vanishes ($\ddot{\mathbf{u}}=\mathbf{0}$),
and the equilibrium equation reduces to:
\begin{equation}
\sigma_{ij,j} - h_{ijk,jk} + b_{i} = 0
\qquad\text{in }\Omega.
\label{s8:equilibrium}
\end{equation}

\subsection{Mesh description}
\label{ssec:mesh}

The domain $\Omega$ is discretized using $C^{1}$ BFS rectangular
elements. The mesh is structured and uniform for the benchmark problem.

\textbf{Mesh parameters:}
\begin{itemize}
\item Number of elements in $X_{1}$-direction: $N_{e1}$ = varied for the mesh convergence study
\item Number of elements in $X_{2}$-direction: $N_{e2}$ = varied for the mesh convergence study
\item Total number of elements: $N_{e} = N_{e1}\times N_{e2}$
\item Element size in $X_{1}$-direction: $h_{1} = L_{1}/N_{e1}$
\item Element size in $X_{2}$-direction: $h_{2} = L_{2}/N_{e2}$
\item Characteristic element size: $h = \sqrt{h_{1}h_{2}}$
\end{itemize}

\textbf{Degrees of freedom:}
\begin{itemize}
\item DOFs per node: 4 (displacement $u_{i}$, first derivatives
  $u_{i,1}$, $u_{i,2}$, and mixed derivative $u_{i,12}$)
\item Nodes per element: 4
\item DOFs per element: $4\times 4\times 2 = 32$
\item Total number of nodes: $(N_{e1}+1)\times(N_{e2}+1)$
\item Total number of DOFs: $4\times 2\times (N_{e1}+1)\times(N_{e2}+1)$
\end{itemize}

\textbf{Mesh convergence study:} To assess convergence, we perform
analyses on a sequence of refined meshes:
\begin{equation}
N_{e}^{(k)} = 4^{k}\,N_{e}^{(0)},
\qquad
k=0,1,2,\ldots,K,
\label{s8:mesh_sequence}
\end{equation}
where $N_{e}^{(0)}$ is the coarsest mesh and $K$ is the number of
refinement levels. The element size decreases as $h^{(k)} = h^{(0)}/2^{k}$.

\subsection{Exact benchmark solution}
\label{ssec:exact_solution}

The 2022 benchmark provides the exact solution
$\mathbf{u}^{\mathrm{exact}}(\bx)$, strain field
$\boldsymbol{\varepsilon}^{\mathrm{exact}}(\bx)$, and stress field
$\boldsymbol{\sigma}^{\mathrm{exact}}(\bx)$ for the specified boundary
and loading conditions.

\textbf{Exact solution components:}
\begin{itemize}
\item Displacement field: $\mathbf{u}^{\mathrm{exact}}(\bx)$ = obtained from the 2022 benchmark
\item Strain field: $\varepsilon_{ij}^{\mathrm{exact}}(\bx) = \tfrac12(u_{i,j}^{\mathrm{exact}}+u_{j,i}^{\mathrm{exact}})$
\item Strain gradient field: $\varepsilon_{ij,k}^{\mathrm{exact}}(\bx) = \partial\varepsilon_{ij}^{\mathrm{exact}}/\partial x_{k}$
\item Stress field: $\sigma_{ij}^{\mathrm{exact}}(\bx)$ = computed from the constitutive law
\item Double stress field: $h_{ijk}^{\mathrm{exact}}(\bx)$ = computed from the gradient constitutive law
\end{itemize}

\textbf{Quantities used for comparison:}
\begin{enumerate}
\item Displacement field $\mathbf{u}^{h}(\bx)$ at Gauss points
\item Strain field $\boldsymbol{\varepsilon}^{h}(\bx)$ at Gauss points
\item Strain gradient field $\nabla\boldsymbol{\varepsilon}^{h}(\bx)$ at Gauss points
\item Stress field $\boldsymbol{\sigma}^{h}(\bx)$ at Gauss points
\item Double stress field $\mathbf{h}^{h}(\bx)$ at Gauss points
\item Reaction forces at constrained nodes
\item Total strain energy $U^{h}$
\end{enumerate}

\subsection{Error norms}
\label{ssec:error_norms}

We define three error norms to quantify the accuracy of the finite
element solution relative to the exact benchmark solution
\cite{Hughes1987,ZienkiewiczTaylor2000}.

\subsubsection{$L^{2}$ error norm}

The $L^{2}$ norm measures the pointwise error in the displacement or
strain field:
\begin{equation}
\|\mathbf{e}\|_{L^{2}}
:= \left(\int_{\Omega}
   \|\mathbf{u}^{h}(\bx)-\mathbf{u}^{\mathrm{exact}}(\bx)\|^{2}\,
   \mathrm{d}\Omega\right)^{1/2},
\label{s8:L2_norm}
\end{equation}
where $\|\cdot\|$ denotes the Euclidean norm of the vector field. For
strain fields:
\begin{equation}
\|\mathbf{e}_{\varepsilon}\|_{L^{2}}
:= \left(\int_{\Omega}
   \|\boldsymbol{\varepsilon}^{h}(\bx)-\boldsymbol{\varepsilon}^{\mathrm{exact}}(\bx)\|^{2}\,
   \mathrm{d}\Omega\right)^{1/2}.
\label{s8:L2_strain}
\end{equation}

\subsubsection{$H^{1}$ error norm}

The $H^{1}$ norm measures the error in both the field and its gradient:
\begin{equation}
\|\mathbf{e}\|_{H^{1}}
:= \left(\int_{\Omega}
   \bigl[\,\|\mathbf{u}^{h}-\mathbf{u}^{\mathrm{exact}}\|^{2}
   + \|\nabla\mathbf{u}^{h}-\nabla\mathbf{u}^{\mathrm{exact}}\|^{2}\,\bigr]\,
   \mathrm{d}\Omega\right)^{1/2}.
\label{s8:H1_norm}
\end{equation}

\subsubsection{$L^{\infty}$ error norm}

The $L^{\infty}$ norm measures the maximum pointwise error:
\begin{equation}
\|\mathbf{e}\|_{L^{\infty}}
:= \max_{\bx\in\Omega}
   \|\mathbf{u}^{h}(\bx)-\mathbf{u}^{\mathrm{exact}}(\bx)\|.
\label{s8:Linfinity_norm}
\end{equation}

\subsubsection{Energy norm}

The energy norm follows the bilinear form introduced in the weak formulation (Section~4, Eq.~\eqref{s4:Wint}).

\subsubsection{Relative error}

The relative error normalizes the absolute error by the exact solution:
\begin{equation}
e_{\mathrm{rel}}
:= \frac{\|\mathbf{e}\|}{\|\mathbf{u}^{\mathrm{exact}}\|}.
\label{s8:relative_error}
\end{equation}

\subsection{Mesh convergence study}
\label{ssec:convergence_study}

The mesh convergence study will demonstrate that the finite element solution
converges to the exact benchmark solution as the mesh is refined.

\textbf{Convergence rate:} The convergence rate will be evaluated numerically from the mesh convergence study. The numerical convergence rate is computed from successive mesh refinements:
\begin{equation}
r = \frac{\log(e_{k}/e_{k+1})}{\log(h_{k}/h_{k+1})}
= \frac{\log(e_{k}/e_{k+1})}{\log(2)},
\label{s8:numerical_rate}
\end{equation}
where $e_{k}$ is the error on mesh $k$ and $h_{k}/h_{k+1}=2$ for uniform refinement.

\subsection{Verification tables}
\label{ssec:tables}

The following tables present the verification results. Numerical values
are placeholders to be filled from MATLAB computations.

\subsubsection{Displacement error norms}

\begin{longtable}[c]{>{\raggedright\arraybackslash}p{0.12\linewidth}
>{\raggedright\arraybackslash}p{0.15\linewidth}
>{\raggedright\arraybackslash}p{0.15\linewidth}
>{\raggedright\arraybackslash}p{0.15\linewidth}
>{\raggedright\arraybackslash}p{0.15\linewidth}
>{\raggedright\arraybackslash}p{0.15\linewidth}
>{\raggedright\arraybackslash}p{0.15\linewidth}}
\caption{Displacement error norms for mesh convergence study.}
\label{tab:disp_errors}\\
\toprule
Mesh & $N_{e}$ & $h$ & $\|\mathbf{e}\|_{L^{2}}$ & Rate & $\|\mathbf{e}\|_{H^{1}}$ & Rate \\
\midrule
\endfirsthead
\toprule
Mesh & $N_{e}$ & $h$ & $\|\mathbf{e}\|_{L^{2}}$ & Rate & $\|\mathbf{e}\|_{H^{1}}$ & Rate \\
\midrule
\endhead
Coarse & to be computed & to be computed & to be computed & -- & to be computed & -- \\
Medium & to be computed & to be computed & to be computed & to be computed & to be computed & to be computed \\
Fine & to be computed & to be computed & to be computed & to be computed & to be computed & to be computed \\
Very fine & to be computed & to be computed & to be computed & to be computed & to be computed & to be computed \\
\bottomrule
\end{longtable}

\textbf{Acceptance criteria:} The convergence rates will be evaluated numerically and compared against the exact benchmark solution.

\subsubsection{Strain energy comparison}

\begin{longtable}[c]{>{\raggedright\arraybackslash}p{0.15\linewidth}
>{\raggedright\arraybackslash}p{0.25\linewidth}
>{\raggedright\arraybackslash}p{0.25\linewidth}
>{\raggedright\arraybackslash}p{0.35\linewidth}}
\caption{Strain energy comparison with exact benchmark.}
\label{tab:energy}\\
\toprule
Mesh & $U^{h}$ (FE) & $U^{\mathrm{exact}}$ (benchmark) & Relative error \\
\midrule
\endfirsthead
\toprule
Mesh & $U^{h}$ (FE) & $U^{\mathrm{exact}}$ (benchmark) & Relative error \\
\midrule
\endhead
Coarse & to be computed & to be computed & to be computed \\
Medium & to be computed & to be computed & to be computed \\
Fine & to be computed & to be computed & to be computed \\
Very fine & to be computed & to be computed & to be computed \\
\bottomrule
\end{longtable}

\textbf{Acceptance criteria:} The strain energy will be compared against the
exact benchmark value, and convergence will be verified as $h\to 0$.
\subsection{Verification figures}
\label{ssec:figures}

The following figures should be generated from MATLAB results to
visually demonstrate benchmark recovery.

\subsubsection{Figure 1: Mesh convergence plot}

Plot the error norms ($L^{2}$, $H^{1}$, $L^{\infty}$) versus element
size $h$ on a log-log scale to assess convergence rates.

\subsubsection{Figure 2: Displacement field comparison}

Plot the finite element displacement field $\mathbf{u}^{h}(\bx)$ and
the exact solution $\mathbf{u}^{\mathrm{exact}}(\bx)$ side by side to
visually assess solution accuracy.

\subsubsection{Figure 3: Error distribution}

Plot the pointwise error $\|\mathbf{u}^{h}(\bx)-\mathbf{u}^{\mathrm{exact}}(\bx)\|$
over the domain to visualize spatial error distribution.

\subsubsection{Figure 4: Strain field comparison}

Plot the finite element strain field $\boldsymbol{\varepsilon}^{h}(\bx)$
and the exact strain field $\boldsymbol{\varepsilon}^{\mathrm{exact}}(\bx)$.

\subsubsection{Figure 5: Stress field comparison}

Plot the finite element stress field $\boldsymbol{\sigma}^{h}(\bx)$
and the exact stress field $\boldsymbol{\sigma}^{\mathrm{exact}}(\bx)$.

\subsection{Discussion and verification statement}
\label{ssec:discussion}

\textbf{Verification results:}
To be completed after numerical computations are performed.

\textbf{Mesh convergence:} The finite element solution will converge to the
exact benchmark solution as the mesh is refined. The observed convergence rates will be reported after the numerical computations.

\textbf{Strain energy:} The strain energy will converge to the exact
benchmark value. The relative error will be computed.

\textbf{Field comparisons:} The displacement, strain, and stress fields
will be compared with the exact benchmark solution. The maximum
pointwise error will be computed.

\subsection{Final verification statement}
\label{ssec:final_statement}

\vspace{0.5cm}

\noindent\fbox{\parbox{0.95\textwidth}{
\centering
\textbf{VERIFICATION FRAMEWORK:}

\vspace{0.3cm}

This verification framework will demonstrate that the finite-strain
$C^{1}$ BFS Mindlin Form-II strain-gradient formulation developed in
Sections~1--7 \textbf{reduces exactly} to the validated 2022 Mindlin
Form-II benchmark formulation in the infinitesimal-strain limit.

\vspace{0.3cm}

The numerical implementation will be shown to correctly solve the
derived equations through:
\begin{itemize}
\item Convergence to the exact benchmark solution
\item Mesh convergence analysis
\item Field comparisons
\end{itemize}

\vspace{0.3cm}

Successful agreement with the validated 2022 benchmark will constitute verification of the numerical implementation.
}}

\vspace{0.5cm}

% =====================================================================
% SECTION 9 — NUMERICAL EXPERIMENT DESIGN
% Master Version 1.2 — Sections 1--8 frozen and unchanged.
% This section defines ONLY the experimental plan.  No numerical
% results, convergence rates, error values, CPU timings, or plots
% are reported; every statement is written in future tense.
% Compile as part of the master document (inherits all notation
% and macro definitions from Sections 1--8).
% =====================================================================

\section{Numerical Experiment Design}
\label{sec:numexp}

% =====================================================================
\subsection{Purpose, scope, and design principles}
\label{ssec:s9_purpose}
% =====================================================================

This section defines the complete numerical-experimental programme
that will be executed using the MATLAB implementation of the
geometrically nonlinear, displacement-based, $C^1$
Bogner--Fox--Schmit (BFS) Mindlin Form-II strain-gradient
finite-element formulation derived in Sections~1--7, using the
verification framework established in Section~8.

\textbf{Status.} No numerical computation has been performed at the
time of writing.  All statements in this section are therefore in
future tense.  Tables contain placeholder captions and column
headings only; numerical entries will be populated once MATLAB
results are available.  Figures are described by their intended
content; no convergence rate, error norm, load--displacement value,
stiffness ratio, or iteration count is stated or implied.

\textbf{Design principles.}
\begin{enumerate}
  \item \emph{Mandatory recovery first.}  The validated 2022
    small-strain Mindlin Form-II benchmark will be recovered exactly
    (to within discretisation error) before any finite-strain
    simulation is accepted.  This is the non-negotiable acceptance
    gate for the entire code base.
  \item \emph{Single variable at a time.}  Each parametric study will
    vary one well-defined control parameter while holding all others
    fixed, so that cause and effect remain unambiguously
    attributable.
  \item \emph{Isolation of nonlinearities.}  Every finite-strain
    result will be compared against both (i)~the small-strain
    strain-gradient solution (isolates geometric nonlinearity) and
    (ii)~the finite-strain classical ($\ell=0$) Saint-Venant--Kirchhoff
    solution (isolates gradient effects at finite deformation).
  \item \emph{Scope discipline.}  The programme is sized for a single
    Q1 journal paper: six numerical studies, 8--12 figures, and
    6--8 tables.  Ideas that exceed this scope are listed without
    elaboration in Subsection~\ref{ssec:future}.
  \item \emph{Notation locked.}  All kinematic, constitutive, and
    discretisation notation follows the frozen conventions of
    Sections~1--8 without modification.
\end{enumerate}

% =====================================================================
\subsection{Common computational setup}
\label{ssec:s9_setup}
% =====================================================================

\paragraph{Implementation platform.}
All simulations will be executed in MATLAB using the modular solver
architecture specified in Section~7 (main driver, Newton--Raphson
iterator, element-level computation of the four tangent
contributions $\mathbf{K}_{\mathrm{mat,cl}}^{e}$,
$\mathbf{K}_{\mathrm{geo,cl}}^{e}$, $\mathbf{K}_{\mathrm{mat,gr}}^{e}$,
$\mathbf{K}_{\mathrm{geo,gr}}^{e}$, $3\times3$ Gauss--Legendre
quadrature, and scatter-assembly into sparse global arrays).

\paragraph{Kinematic and constitutive model.}
All studies will use the plane-strain reduction (assumptions A7--A8 of
Section~2) with the Saint-Venant--Kirchhoff--Mindlin Form-II
free-energy density
\begin{equation}
\psi(\mathbf{E},\bm{\mathcal{G}})
  = \tfrac12\lambda(\tr\mathbf{E})^{2}
  + \mu\,\mathbf{E}{:}\mathbf{E}
  + \tfrac12\mathbb{D}_{ABCPQR}\,\mathcal{G}_{ABC}\mathcal{G}_{PQR},
\qquad \bm{\mathcal{G}}=\gradz\mathbf{E},
\end{equation}
of Section~2, eq.~\eqref{eq:sec2_complete_energy_tensor}.  The
classical Lam\'e constants $\lambda,\mu$ will be derived from a
Young's modulus $E$ and Poisson ratio $\nu$ reported per study;
the fifth-order isotropic gradient constants $a_{1},\ldots,a_{5}$
will be set per study and will define the internal length scale
\begin{equation}
\ell = \sqrt{\frac{\bar{a}}{\mu}},\qquad
\bar{a} = \tfrac{1}{5}\sum_{i=1}^{5}a_{i},
\end{equation}
consistent with Section~8, eq.~\eqref{s8:length_scale}.  A
non-dimensional gradient parameter will be defined in each study as
the ratio of $\ell$ to a characteristic geometric dimension~$L$
(plate width, beam height, etc.), so that size effects can be
reported in a scale-invariant form.

\paragraph{Reference domain and meshing.}
All problems will be defined on rectangular reference domains
$\Omega_{0}=[0,L_{1}]\times[0,L_{2}]$ compatible with the structured
BFS rectangle.  Uniform $N_{e1}\times N_{e2}$ meshes will be used
unless stated otherwise.  Nodal degrees of freedom will be the BFS
set $(u_{i},\,u_{i,1},\,u_{i,2},\,u_{i,12})$ per node per in-plane
component (Section~6, eqs.~\eqref{s6:DOF1}--\eqref{s6:DOF4}).

\paragraph{Loading and solver control.}
Loads will be applied incrementally via the dead-load schedule of
Section~7 (eqs.~\eqref{s7:fixed_step}--\eqref{s7:increase_step}).
Newton--Raphson iteration will use the three simultaneous
convergence criteria of eqs.~\eqref{s7:conv_residual}--\eqref{s7:conv_test}
(residual, increment, energy) with tolerances
$\mathrm{tol}_{R}=\mathrm{tol}_{D}=\mathrm{tol}_{E}=10^{-8}$ and
a ceiling of $k_{\max}=20$ iterations per step; adaptive
load-step cutback ($\alpha=0.5$) and growth ($\beta=1.2$,
$k_{\mathrm{target}}=6$) will be enabled for all finite-strain
studies.  Essential boundary conditions on $u_i$ and
$\partial u_i/\partial N$ will be imposed by the row-column
zeroing procedure described in Algorithm~\ref{alg:complete};
natural boundary conditions (single traction $\bar{\mathbf{T}}_{0}$
and double traction $\bar{\mathbf{d}}$) will be applied through
the external force vector $\mathbf{f}_{\mathrm{ext}}^{e}$
(eq.~\eqref{s7:f_ext}) via 3-point 1D Gauss edge integration.

\paragraph{Non-dimensionalisation.}
Where reported, displacements will be normalised by a geometric
length~$L$, strains are already dimensionless, stresses and
double stresses will be normalised by an elastic modulus
(typically $\mu$ or $E$), and total strain energy $\Pi_{\mathrm{int}}$
will be normalised by a reference energy scale $\mu L^{2}$
(per unit thickness in plane strain).  This will allow the
finite-strain, small-strain gradient, and classical
($\ell=0$) solutions to be compared on a common axis.

\paragraph{Overkill reference.}
For studies in which no closed-form analytical solution exists
(essentially every finite-strain case), a finely resolved
``overkill'' mesh will be computed and used as the numerical
reference for error and convergence assessment, in line with
standard nonlinear FEM verification practice
\cite{Bathe2006,ZienkiewiczTaylor2000}.

% =====================================================================
\subsection{Study~1 -- Benchmark verification (mandatory 2022 recovery)}
\label{ssec:s9_study1}
% =====================================================================

\paragraph{Scientific purpose.}
Satisfy the mandatory recovery rule: demonstrate that the finite-strain
BFS formulation, when exercised in the small-displacement-gradient
regime with the same material constants, geometry, boundary
conditions and loading as the validated 2022 Mindlin Form-II
benchmark \cite{TorabiNiiranen2022}, reproduces the benchmark
solution to within discretisation error.  This study is the
entry gate: no subsequent study will be accepted until this
verification passes.

\paragraph{What will be performed.}
The 2022 benchmark problem will be reproduced verbatim (same
domain dimensions, same $E,\nu,a_{1},\ldots,a_{5}$, same essential
boundary conditions $u_i=\bar{u}_i$, $\partial u_i/\partial n=\bar{g}_i$,
and same natural data $\bar{t}_i,\bar{d}_i$) on a sequence of
uniformly refined BFS meshes.  Deformation will be restricted to
the infinitesimal regime by scaling the load so that the maximum
Green--Lagrange strain satisfies $\|\mathbf{E}\|_{\infty}\le 10^{-6}$;
under this condition the geometric nonlinearity is $O(\|\mathbf{H}\|^{2})$
and is numerically negligible at the reported tolerance, so that the
finite-strain code exercises the same linear limit documented in
Sections~4 and~8.  The classical ($\ell=0$) case will also be run
to confirm that the code collapses to linear isotropic elasticity
when gradient moduli are zeroed.

\paragraph{Quantities to be extracted.}
Nodal displacements, displacement gradients, Green--Lagrange strain
$\mathbf{E}$, strain gradient $\bm{\mathcal{G}}$, second
Piola--Kirchhoff stress $\mathbf{S}$, double stress
$\boldsymbol{\Sigma}$, total strain energy $\Pi_{\mathrm{int}}$,
and reaction forces at constrained boundaries, on every mesh of
the refinement sequence.  These will be compared pointwise against
the analytical/semi-analytical benchmark fields.

\paragraph{Acceptance criteria.}
(i)~Monotonic mesh convergence in all error norms to the benchmark
solution as $h\to0$; (ii)~observed convergence rates consistent with
the $C^1$ BFS approximation order (rates will be reported as
computed, not pre-stated); (iii)~pointwise agreement of displacement,
strain, strain gradient, stress, and double-stress fields with the
benchmark at the finest mesh to within discretisation error;
(iv)~strain energy converging to the benchmark value;
(v)~reactions consistent with global equilibrium.

\paragraph{Planned tables.}
\begin{itemize}
  \item Table~\ref{tab:s9_disp_bench} -- displacement, gradient,
    strain, and energy error norms on each mesh, with observed
    convergence rates.
  \item Table~\ref{tab:s9_energy_bench} -- total strain energy on
    each mesh vs.\ the benchmark value, with relative error.
  \item Table~\ref{tab:s9_field_bench} -- pointwise field
    comparison at selected probe points vs.\ the benchmark.
\end{itemize}

\paragraph{Planned figures.}
\begin{itemize}
  \item Figure~\ref{fig:s9_conv_bench} -- log--log plot of $L^2$,
    $H^1$, and $L^\infty$ displacement error norms vs.\ element
    size $h$; benchmark convergence reference lines will be added
    once rates are measured.
  \item Figure~\ref{fig:s9_disp_field_bench} -- contour plots of
    $u_1$, $u_2$ from the finest BFS mesh side-by-side with the
    benchmark displacement field.
  \item Figure~\ref{fig:s9_err_field_bench} -- spatial distribution
    of the pointwise displacement error
    $\|\mathbf{u}^{h}-\mathbf{u}^{\mathrm{exact}}\|$ on the finest
    mesh.
\end{itemize}

% =====================================================================
\subsection{Study~2 -- Mesh convergence of the finite-strain formulation}
\label{ssec:s9_study2}
% =====================================================================

\paragraph{Scientific purpose.}
Establish the numerical consistency and convergence behaviour of the
fully nonlinear BFS formulation.  Study~1 confirmed convergence in
the linear limit; Study~2 extends this assessment to the
finite-deformation regime where no analytical solution is
available and where both geometric stiffness contributions
($\mathbf{K}_{\mathrm{geo,cl}}$, $\mathbf{K}_{\mathrm{geo,gr}}$)
are active.

\paragraph{What will be performed.}
A rectangular-domain boundary-value problem with a known nonlinear
character (e.g.\ a cantilever strip in end shear producing large
rotations, or a block under shear/extension) will be solved on a
nested sequence of uniformly refined meshes.  At least two load
levels will be considered: a small load (linear regime, cross-check
with Study~1) and a finite-deformation load producing appreciable
rotation and nonlinear strain.  For each load level the problem
will be solved on $K$ mesh densities (e.g.\ $N_e = N_{0}, 4N_{0},
16N_{0}, 64N_{0}$ elements, subject to computational cost).  The
finest-mesh solution will serve as the overkill reference against
which errors on coarser meshes will be measured.  The $\ell=0$
classical SVK case will also be run on the same mesh sequence to
isolate the influence of the gradient terms on convergence
behaviour.

\paragraph{Quantities to be extracted.}
Displacement field at Gauss points and nodes, strain energy,
$L^2$ and $H^1$ error norms relative to overkill, Newton--Raphson
iteration count per load step, number of adaptive step cutbacks,
and total (cumulative) iteration count as a proxy for solver
robustness.

\paragraph{Acceptance criteria.}
(i)~Monotonic convergence of displacement and energy to the
overkill solution with mesh refinement at both load levels;
(ii)~convergence rates on the linear problem consistent with
Study~1; (iii)~convergence of the finite-strain problem will
be reported as observed -- no rate is pre-supposed because the
Newton--Raphson path and linearisation around a nonlinear state
may produce pre-asymptotic behaviour on coarse meshes.

\paragraph{Planned tables.}
\begin{itemize}
  \item Table~\ref{tab:s9_conv_nl} -- error norms, rates, DOF
    count, and total Newton iterations per mesh, at small and
    finite load levels.
  \item Table~\ref{tab:s9_conv_noclg} -- same as above but for the
    $\ell=0$ classical SVK reference, to permit direct
    comparison of the effect of the gradient branch on
    convergence.
\end{itemize}

\paragraph{Planned figures.}
\begin{itemize}
  \item Figure~\ref{fig:s9_conv_nl_loglog} -- log--log plot of
    displacement and energy errors vs.\ $h$ at small and finite
    load levels.
  \item Figure~\ref{fig:s9_newton_hist} -- representative
    Newton--Raphson residual-norm histories ($\|\mathbf{R}^{(k)}\|_2$
    versus iteration $k$) at a mid-range load step on the finest
    mesh, for both the gradient and $\ell=0$ cases, illustrating
    the observed convergence behaviour of the consistent
    Newton--Raphson algorithm.
\end{itemize}

% =====================================================================
\subsection{Study~3 -- Internal length-scale parameter study}
\label{ssec:s9_study3}
% =====================================================================

\paragraph{Scientific purpose.}
Characterise the size-dependent response of the model -- the
defining feature of strain-gradient elasticity -- at both small
and finite deformation.  The objective is to quantify how the
apparent stiffness, boundary-layer width, and field distributions
change as the ratio $\ell/L$ is varied between the classical limit
and the strongly gradient-dominated regime, and to determine whether
finite-strain kinematics modify these size effects relative to the
linear prediction.

\paragraph{What will be performed.}
A prototypical size-effect problem (e.g.\ a short beam or plate
strip in bending/shear, with at least one dimension comparable to
$\ell$) will be solved for a range of normalised length scales
$\ell/L\in[0,\ell_{\max}/L]$.  The set of gradient constants
$a_1,\ldots,a_5$ will be held in fixed proportion (so that $\ell$
is varied via a single scalar multiplier on the gradient modulus
tensor $\mathbb{D}$), keeping the isotropic structure unchanged.
Each value of $\ell/L$ will be run at two load levels: a small
load (linear benchmark regime, cross-check against the 2022
linear size-effect trends) and a finite-deformation load.  The
domain dimension $L$ (or equivalently $\ell$) will be varied on
a sufficiently fine mesh (chosen from Study~2 to lie in the
asymptotic convergence range) so that discretisation error does
not contaminate the size-effect trend.

\paragraph{Quantities to be extracted.}
Apparent (effective) stiffness (total load / characteristic
displacement), peak displacement, peak strain, width of the
boundary layer near clamped/free edges, maximum double stress,
and total strain energy, for each $\ell/L$ and each load level.
The $\ell=0$ classical solution will be included as the
reference limit.

\paragraph{Acceptance criteria.}
(i)~Convergence of displacement, strain energy, peak strain,
peak double stress, and apparent stiffness with mesh
refinement will be verified against the Study~2 mesh-convergence
tolerance at each $\ell/L$; (ii)~for each $\ell/L$ the small-load
response will be compared against the linear strain-gradient
solution obtained from Study~1/2 at the same $\ell/L$;
(iii)~differences between the finite-load response and the
small-load response will be reported, and will be quantified
relative to the $\ell=0$ finite-strain control and the
$\ell>0$ small-strain control.

\paragraph{Planned tables.}
\begin{itemize}
  \item Table~\ref{tab:s9_sizeeffect} -- normalised apparent
    stiffness and peak displacement vs.\ $\ell/L$ at small and
    finite load levels.
\end{itemize}

\paragraph{Planned figures.}
\begin{itemize}
  \item Figure~\ref{fig:s9_sizeeffect_curve} -- apparent
    stiffness normalised by the classical ($\ell=0$) stiffness
    vs.\ $\ell/L$, on log--linear axes, at small and finite
    load.
  \item Figure~\ref{fig:s9_bl_profiles} -- displacement and
    cross-sectional strain profiles through the boundary layer
    at several $\ell/L$ values, illustrating the emergence
    and scaling of gradient boundary layers.
  \item Figure~\ref{fig:s9_l_field} -- contour plots of
    $|\bm{\mathcal{G}}|$ at small and finite load for a
    representative $\ell/L$, visualising the spatial
    distribution of gradient activity.
\end{itemize}

% =====================================================================
\subsection{Study~4 -- Finite-strain demonstration: geometric
  nonlinearity and load--deflection response}
\label{ssec:s9_study4}
% =====================================================================

\paragraph{Scientific purpose.}
Demonstrate the primary novelty of the present formulation
relative to the 2022 benchmark: its ability to handle
geometrically nonlinear, finite-rotation, moderate-stretch
response while retaining full strain-gradient character.  This
study produces the signature load--deflection curves and
deformed configurations that anchor the paper's claim of a
finite-strain extension.

\paragraph{What will be performed.}
Two canonical rectangular-domain problems that exhibit strong
geometric nonlinearity on a BFS-friendly mesh will be analysed:
\begin{enumerate}
  \item[(a)] \emph{Cantilever strip under end shear.}  A
    rectangular strip of length $L$, height $h$, fully clamped
    on one edge ($u_i=0$, $\partial u_i/\partial N=0$) and
    subjected to a distributed tangential dead load on the
    opposite (free) edge, producing large end rotations and
    nonlinear load--deflection response.  This is the
    classical geometric-nonlinearity benchmark and allows
    direct comparison with classical finite-strain SVK beam
    solutions when $\ell=0$.
  \item[(b)] \emph{Block under simple shear / extension.}  A
    square block with one edge fully clamped and the opposite
    edge subjected to prescribed displacement (and/or
    distributed traction) producing a finite homogeneous or
    near-homogeneous deformation, used to verify that the code
    reproduces homogeneous finite strain correctly and to
    expose any unphysical locking or boundary artefacts from
    the gradient terms under large stretch.
\end{enumerate}
For each problem, a full incremental load--deflection history
will be computed up to a target maximum displacement (or load)
chosen to produce appreciable but still moderate finite
deformation (consistent with the local-stability discussion of
Section~2 for Saint-Venant--Kirchhoff materials, which are not
globally polyconvex).  Each problem will be run for (i)~the
classical $\ell=0$ case (SVK reference), (ii)~a representative
gradient case ($\ell/L$ chosen from Study~3 to show measurable
size effect), and (iii)~a strongly gradient-dominated case.
Small-load responses will be compared against Study~1/3 for
consistency.

\paragraph{Quantities to be extracted.}
Load--deflection curves (reaction force vs.\ tip/edge
displacement), deformed configurations at several load
levels, Green--Lagrange strain and strain-gradient contours,
Cauchy stress (push-forward) distributions, Newton iteration
statistics, and any path-dependent features such as stiffening
or softening relative to the linear and classical predictions.

\paragraph{Acceptance criteria.}
(i)~Mesh resolution will be confirmed sufficient (via the
Study~2 mesh-convergence check) for each problem and each
$\ell/L$ before reporting load--deflection data; (ii)~the
small-load tangent stiffness extracted from each
load--deflection curve will be compared against the linear
stiffness computed at the same $\ell/L$ from the Study~3
data; (iii)~global equilibrium (reactions vs.\ applied load)
will be satisfied to within the Newton--Raphson tolerance at
every reported increment; (iv)~the $\ell=0$ SVK response
will be reported alongside every $\ell>0$ curve so that
gradient effects can be isolated.

\paragraph{Planned tables.}
\begin{itemize}
  \item Table~\ref{tab:s9_loaddisp} -- load factor vs.\
    characteristic displacement at selected increments, for
    each $\ell/L$ and each benchmark problem.
  \item Table~\ref{tab:s9_newton_stats} -- average and
    worst-case Newton iterations per step, step cutbacks,
    and final load factor achieved, as a function of
    $\ell/L$ and load level, documenting solver robustness.
\end{itemize}

\paragraph{Planned figures.}
\begin{itemize}
  \item Figure~\ref{fig:s9_cant_ld} -- load vs.\ tip
    displacement for the cantilever strip for several $\ell/L$
    values, with the classical ($\ell=0$) reference and the
    linear gradient prediction overlaid.
  \item Figure~\ref{fig:s9_cant_deformed} -- overlaid deformed
    configurations of the cantilever at several load levels
    for a representative $\ell/L$, illustrating finite
    rotations.
  \item Figure~\ref{fig:s9_block_contours} -- contour plots of
    $E_{11}$, $E_{12}$, and $|\bm{\mathcal{G}}|$ for the
    block problem at finite load for $\ell=0$ and $\ell>0$.
\end{itemize}

% =====================================================================
\subsection{Study~5 -- Sensitivity to the five gradient-elastic
  constants}
\label{ssec:s9_study5}
% =====================================================================

\paragraph{Scientific purpose.}
Determine which of the five isotropic gradient-elastic constants
$a_{1},\ldots,a_{5}$ dominate the finite-strain response, and
assess whether sensitivities observed in the linear regime
carry over to finite deformation or are modified by geometric
nonlinearity.  This supports both physical interpretation
(identifying the materially relevant gradient modes) and
future model reduction (e.g.\ reduced one- or two-length-scale
specialisations).

\paragraph{What will be performed.}
A single representative boundary-value problem (e.g.\ the
cantilever strip of Study~4 at a fixed finite load level and a
fixed $\ell/L$ chosen to place the response in the
gradient-sensitive regime) will be chosen as the baseline.
A one-at-a-time (OAT) sensitivity sweep will be performed in
which each constant $a_i$ is individually scaled above and
below its baseline value (e.g.\ multiplicative factors
selected to span the positive-definite admissible range for
$\mathbb{D}$) while the other four are held fixed.  For each
perturbed parameter set the problem will be re-solved on the
Study~2 mesh.  To screen for two-factor interactions, a
small two-level factorial sweep on the most influential
constants (identified from the OAT sweep) will additionally
be performed if computational budget permits; otherwise this
will be deferred to future work (Subsection~\ref{ssec:future}).

\paragraph{Quantities to be extracted.}
Normalised sensitivity of tip displacement, peak strain, peak
double stress, and total strain energy to each $a_i$;
variation in apparent stiffness and boundary-layer width as
each $a_i$ is varied; identification of any parameter whose
perturbation triggers loss of positive definiteness or
Newton convergence failure (to be reported as a finding, not
pre-judged).

\paragraph{Acceptance criteria.}
(i)~Symmetry and baseline reproduction (each sweep returns to
baseline at unit factor); (ii)~ranking of constants by
sensitivity will be reported as observed; (iii)~any
finite-strain-specific sensitivity amplification not present
in the linear regime will be identified by comparing
small-load and finite-load sweeps.

\paragraph{Planned tables.}
\begin{itemize}
  \item Table~\ref{tab:s9_sensitivity} -- normalised change
    in characteristic response quantities per multiplicative
    factor on each $a_i$, at small and finite load.
\end{itemize}

\paragraph{Planned figures.}
\begin{itemize}
  \item Figure~\ref{fig:s9_tornado} -- tornado-style
    horizontal-bar sensitivity chart ranking $a_1$--$a_5$ by
    their influence on apparent stiffness at small and finite
    load.
  \item Figure~\ref{fig:s9_sensitivity_curves} -- response
    curves (apparent stiffness vs.\ multiplicative factor)
    for the two most influential constants, showing linear
    vs.\ finite-strain sensitivity.
\end{itemize}

% =====================================================================
\subsection{Study~6 -- Higher-order boundary-condition comparison}
\label{ssec:s9_study6}
% =====================================================================

\paragraph{Scientific purpose.}
Because the formulation is $C^1$ and supports both essential
prescription of the normal slope $\partial u_i/\partial N$ and
natural prescription of the double traction $\mathcal{D}_i$
(Section~4), the physical interpretation of boundary conditions
near free and supported edges requires explicit documentation.
This study examines how the choice of higher-order BC affects
the finite-strain solution, particularly the boundary layers
that distinguish gradient from classical elasticity.

\paragraph{What will be performed.}
The cantilever strip (or a simplified shorter strip) will be
re-analysed with the same geometry, material, and body/edge
loading but with the supported edge modelled under two
distinct higher-order BC specifications admitted by the
$C^1$ formulation:
\begin{enumerate}
  \item[(i)]  \emph{Fully clamped (essential/essential):}
    $u_i=\bar{u}_i$ \emph{and}
    $\partial u_i/\partial N=\bar{g}_i=0$ on the supported edge
    (zero displacement and zero normal slope);
  \item[(ii)] \emph{Soft simple support (essential/natural):}
    $u_i=\bar{u}_i$ with zero double traction
    $\mathcal{D}_i=0$ on the supported edge (displacement
    prescribed but normal slope free to develop);
  \item[(iii)] \emph{Traction-free edge (natural/natural):}
    $\bar{T}_{0i}=0$ and $\mathcal{D}_i=0$ on a reference
    free edge, included as the control for comparison.
\end{enumerate}
Each BC will be exercised at a fixed $\ell/L$ (from Study~3)
and at both small and finite load, on the Study~2 convergent
mesh.

\paragraph{Quantities to be extracted.}
Displacement and rotation profiles across the boundary layer,
strain and strain-gradient intensity near the supported edge,
reaction forces and double reactions on the supported edge,
and apparent global stiffness for each BC type.

\paragraph{Acceptance criteria.}
(i)~Mesh resolution will be confirmed sufficient (via the
Study~2 mesh-convergence check) for each BC type and load
level; (ii)~at $\ell=0$ the clamped and soft-simple-support
responses will be compared and their difference will be
reported; (iii)~for $\ell>0$, the peak boundary gradient,
boundary-layer width, apparent stiffness, and edge reactions
will be tabulated for each BC type at both small and finite
load, enabling direct comparison with the underlying
essential/natural BC partition of Section~4.

\paragraph{Planned tables.}
\begin{itemize}
  \item Table~\ref{tab:s9_bc} -- peak boundary gradient,
    boundary-layer width, and global stiffness for each BC
    type, at small and finite load.
\end{itemize}

\paragraph{Planned figures.}
\begin{itemize}
  \item Figure~\ref{fig:s9_bc_layer} -- close-up profiles of
    displacement, rotation, and strain through the first
    $O(\ell)$ distance from the supported edge, for each BC
    type, at small and finite load, illustrating the
    boundary-layer structure enforced by the $C^1$
    formulation.
\end{itemize}

% =====================================================================
\subsection{Interpretation and reporting framework}
\label{ssec:s9_interp}
% =====================================================================

When MATLAB results are available, they will be interpreted
through the following structured comparisons, which will be
applied uniformly across Studies~1--6:

\begin{enumerate}
  \item \emph{Benchmark recovery (Study~1).}  Every
    finite-strain result will be traceable, in the limit
    $\|\gradz\mathbf{u}\|\to 0$, to a 2022-benchmark
    counterpart; deviations at finite load will be attributed
    to geometric nonlinearity rather than to coding or
    constitutive error.
  \item \emph{Gradient vs.\ classical.}  For any reported
    finite-strain result, the $\ell=0$ SVK control will be
    plotted/tabled alongside it, so that the reader can
    isolate the contribution of the gradient terms at
    finite deformation.
  \item \emph{Finite vs.\ infinitesimal.}  For any reported
    gradient result, the small-strain gradient prediction
    will be overlaid to isolate the contribution of
    geometric nonlinearity.
  \item \emph{Mesh resolution.}  Where results depend
    quantitatively on mesh density, they will be reported on
    at least two meshes (or on the convergent Study~2 mesh)
    with an indication of discretisation uncertainty.
  \item \emph{Solver robustness.}  Newton--Raphson performance
    (iteration counts, cutbacks, final load factor attained)
    will be reported alongside every parametric result, so
    that any apparent physical trend that is in fact a
    solver artefact can be identified.
  \item \emph{Units and normalisation.}  All graphs and
    tables will use the non-dimensionalisation of
    Subsection~\ref{ssec:s9_setup} so that results can be
    compared across studies and with published
    strain-gradient benchmarks.
\end{enumerate}

No physical conclusion will be drawn from any simulation that
has not first passed (a)~the Study~1 verification gate and
(b)~the Study~2 mesh-convergence check on a problem-specific
basis.

% =====================================================================
\subsection{Consolidated list of planned tables}
\label{ssec:s9_tables}
% =====================================================================

The following tables will be generated from MATLAB results;
numerical entries are placeholders and will be populated after
computations are complete.

% Table: displacement error norms vs benchmark
\begin{table}[h]
\centering
\caption{Study~1 -- Displacement, gradient, strain, and energy
  error norms versus element size $h$ for the 2022 benchmark
  recovery; convergence rates will be computed from successive
  refinements.}
\label{tab:s9_disp_bench}
\begin{tabular}{llllllll}
\toprule
Mesh & $N_e$ & DOF & $h$ & $\|\mathbf{e}\|_{L^2}$ &
Rate & $\|\mathbf{e}\|_{H^1}$ & Rate \\
\midrule
TBD & TBD & TBD & TBD & TBD & -- & TBD & -- \\
TBD & TBD & TBD & TBD & TBD & TBD & TBD & TBD \\
TBD & TBD & TBD & TBD & TBD & TBD & TBD & TBD \\
TBD & TBD & TBD & TBD & TBD & TBD & TBD & TBD \\
\bottomrule
\end{tabular}
\end{table}

% Table: energy comparison
\begin{table}[h]
\centering
\caption{Study~1 -- Total strain energy $\Pi_{\mathrm{int}}^{h}$
  versus the 2022 benchmark value $U^{\mathrm{exact}}$, with
  relative error.}
\label{tab:s9_energy_bench}
\begin{tabular}{llll}
\toprule
Mesh & $\Pi_{\mathrm{int}}^{h}$ & $U^{\mathrm{exact}}$ &
Relative error \\
\midrule
TBD & TBD & TBD & TBD \\
TBD & TBD & TBD & TBD \\
TBD & TBD & TBD & TBD \\
TBD & TBD & TBD & TBD \\
\bottomrule
\end{tabular}
\end{table}

% Table: pointwise field comparison at probes
\begin{table}[h]
\centering
\caption{Study~1 -- Pointwise comparison of displacement, strain,
  strain gradient, stress, and double stress at selected probe
  points between BFS and the 2022 benchmark on the finest mesh.}
\label{tab:s9_field_bench}
\begin{tabular}{lllllll}
\toprule
Probe & Quantity & BFS & Benchmark & Abs.\ error & Rel.\ error \\
\midrule
TBD & $u_1$ & TBD & TBD & TBD & TBD \\
TBD & $u_2$ & TBD & TBD & TBD & TBD \\
TBD & $E_{11}$ & TBD & TBD & TBD & TBD \\
TBD & $\mathcal{G}_{111}$ & TBD & TBD & TBD & TBD \\
TBD & $S_{11}$ & TBD & TBD & TBD & TBD \\
TBD & $\Sigma_{111}$ & TBD & TBD & TBD & TBD \\
\bottomrule
\end{tabular}
\end{table}

% Table: nonlinear convergence
\begin{table}[h]
\centering
\caption{Study~2 -- Mesh convergence of the finite-strain
  formulation at small and finite load levels, including
  Newton--Raphson iteration statistics.}
\label{tab:s9_conv_nl}
\begin{tabular}{llllllllll}
\toprule
Load & Mesh & DOF & $h$ & $L^2$ err & Rate
  & Energy err & Rate & Total iter & Cutbacks \\
\midrule
Small & TBD & TBD & TBD & TBD & -- & TBD & -- & TBD & TBD \\
Small & TBD & TBD & TBD & TBD & TBD & TBD & TBD & TBD & TBD \\
Small & TBD & TBD & TBD & TBD & TBD & TBD & TBD & TBD & TBD \\
Finite & TBD & TBD & TBD & TBD & -- & TBD & -- & TBD & TBD \\
Finite & TBD & TBD & TBD & TBD & TBD & TBD & TBD & TBD & TBD \\
Finite & TBD & TBD & TBD & TBD & TBD & TBD & TBD & TBD & TBD \\
\bottomrule
\end{tabular}
\end{table}

% Table: size effect
\begin{table}[h]
\centering
\caption{Study~3 -- Normalised apparent stiffness and peak
  displacement versus normalised length scale $\ell/L$ at
  small and finite load levels.}
\label{tab:s9_sizeeffect}
\begin{tabular}{llllll}
\toprule
$\ell/L$ &
$K/K_{\ell=0}$ (small) & $u_{\max}/L$ (small) &
$K/K_{\ell=0}$ (finite) & $u_{\max}/L$ (finite) \\
\midrule
0    & 1 (def.) & TBD & 1 (def.) & TBD \\
TBD  & TBD & TBD & TBD & TBD \\
TBD  & TBD & TBD & TBD & TBD \\
TBD  & TBD & TBD & TBD & TBD \\
TBD  & TBD & TBD & TBD & TBD \\
\bottomrule
\end{tabular}
\end{table}

% Table: load-displacement
\begin{table}[h]
\centering
\caption{Study~4 -- Load factor versus characteristic
  displacement for the cantilever strip at selected increments,
  for several $\ell/L$ values.}
\label{tab:s9_loaddisp}
\begin{tabular}{lllll}
\toprule
Load factor & $u_{\mathrm{tip}}/L$ ($\ell=0$)
  & $u_{\mathrm{tip}}/L$ ($\ell/L=*$)
  & $u_{\mathrm{tip}}/L$ ($\ell/L=**$) \\
\midrule
TBD & TBD & TBD & TBD \\
TBD & TBD & TBD & TBD \\
TBD & TBD & TBD & TBD \\
TBD & TBD & TBD & TBD \\
TBD & TBD & TBD & TBD \\
\bottomrule
\end{tabular}
\end{table}

% Table: Newton stats
\begin{table}[h]
\centering
\caption{Study~4 -- Representative Newton--Raphson solver
  statistics (average iterations per step, worst-case
  iterations, cutbacks, final load attained) for the
  cantilever and block problems at several $\ell/L$.}
\label{tab:s9_newton_stats}
\begin{tabular}{lllllll}
\toprule
Problem & $\ell/L$ & Load regime
  & Avg.\ iter/step & Worst iter & Cutbacks & Final $\lambda$ \\
\midrule
Cantilever & 0 & Finite & TBD & TBD & TBD & TBD \\
Cantilever & * & Finite & TBD & TBD & TBD & TBD \\
Cantilever & ** & Finite & TBD & TBD & TBD & TBD \\
Block & 0 & Finite & TBD & TBD & TBD & TBD \\
Block & * & Finite & TBD & TBD & TBD & TBD \\
\bottomrule
\end{tabular}
\end{table}

% Table: sensitivity
\begin{table}[h]
\centering
\caption{Study~5 -- Normalised sensitivity of characteristic
  responses to one-at-a-time multiplicative perturbation of
  each gradient constant $a_i$, at small and finite load.}
\label{tab:s9_sensitivity}
\begin{tabular}{lllll}
\toprule
$a_i$ & Factor
  & $\Delta K/K$ (small)
  & $\Delta K/K$ (finite)
  & $\Delta\Pi_{\mathrm{int}}/\Pi_{\mathrm{int}}$ (finite) \\
\midrule
$a_1$ & TBD & TBD & TBD & TBD \\
$a_2$ & TBD & TBD & TBD & TBD \\
$a_3$ & TBD & TBD & TBD & TBD \\
$a_4$ & TBD & TBD & TBD & TBD \\
$a_5$ & TBD & TBD & TBD & TBD \\
\bottomrule
\end{tabular}
\end{table}

% Table: BC comparison
\begin{table}[h]
\centering
\caption{Study~6 -- Effect of higher-order boundary-condition
  specification on peak boundary gradient, boundary-layer
  width, and apparent stiffness at small and finite load.}
\label{tab:s9_bc}
\begin{tabular}{lllllll}
\toprule
BC type & Load
  & $|\gradz\mathbf{u}|_{\max}$ at edge
  & $\delta_{\mathrm{BL}}/\ell$
  & $K/K_{\mathrm{clamped}}$ \\
\midrule
Clamped  & Small & TBD & TBD & 1 (def.) \\
Soft SS  & Small & TBD & TBD & TBD \\
Clamped  & Finite & TBD & TBD & 1 (def.) \\
Soft SS  & Finite & TBD & TBD & TBD \\
\bottomrule
\end{tabular}
\end{table}

% =====================================================================
\subsection{Consolidated list of planned figures}
\label{ssec:s9_figures}
% =====================================================================

The following figures will be generated from MATLAB results.
Exact axis ranges, contour levels, and annotations will be
set when numerical data are available.

\begin{enumerate}
  \item[\textbf{Fig.\ref{fig:s9_conv_bench}.}]
    Log--log plot of $L^2$, $H^1$, and $L^\infty$ displacement
    error norms versus element size $h$ for the Study~1
    benchmark recovery, with observed convergence rates
    annotated.
  \item[\textbf{Fig.\ref{fig:s9_disp_field_bench}.}]
    Side-by-side contour plots of the in-plane displacement
    components $(u_1,u_2)$ on the finest BFS mesh and from
    the 2022 benchmark solution.
  \item[\textbf{Fig.\ref{fig:s9_err_field_bench}.}]
    Contour plot of the pointwise displacement error
    magnitude $\|\mathbf{u}^{h}-\mathbf{u}^{\mathrm{exact}}\|$
    over the domain on the finest mesh.
  \item[\textbf{Fig.\ref{fig:s9_conv_nl_loglog}.}]
    Log--log convergence of displacement and energy errors
    versus $h$ for the finite-strain problem at small and
    finite load levels, with the $\ell=0$ reference overlaid.
  \item[\textbf{Fig.\ref{fig:s9_newton_hist}.}]
    Representative Newton--Raphson residual history
    ($\|\mathbf{R}^{(k)}\|_2$ versus iteration $k$) at a
    mid-range load step, illustrating the observed
    convergence behaviour of the consistent Newton--Raphson
    algorithm.
  \item[\textbf{Fig.\ref{fig:s9_sizeeffect_curve}.}]
    Normalised apparent stiffness $K/K_{\ell=0}$ versus
    $\ell/L$ on log--linear axes, at small and finite load,
    showing the size-effect curve.
  \item[\textbf{Fig.\ref{fig:s9_bl_profiles}.}]
    Displacement and cross-sectional strain profiles through
    the boundary layer at several $\ell/L$ values,
    illustrating gradient boundary layers.
  \item[\textbf{Fig.\ref{fig:s9_l_field}.}]
    Contours of $|\bm{\mathcal{G}}|$ at small and finite
    load for a representative $\ell/L$, showing spatial
    gradient activity.
  \item[\textbf{Fig.\ref{fig:s9_cant_ld}.}]
    Load versus tip displacement for the cantilever strip
    for several $\ell/L$ values, with the $\ell=0$ classical
    SVK reference and the linear gradient prediction
    overlaid.
  \item[\textbf{Fig.\ref{fig:s9_cant_deformed}.}]
    Overlaid deformed configurations of the cantilever at
    several load levels for a representative $\ell/L$,
    illustrating finite rotations.
  \item[\textbf{Fig.\ref{fig:s9_block_contours}.}]
    Contour plots of $E_{11}$, $E_{12}$, and
    $|\bm{\mathcal{G}}|$ for the block problem at finite
    load for $\ell=0$ and a representative $\ell>0$.
  \item[\textbf{Fig.\ref{fig:s9_tornado}.}]
    Tornado chart ranking $a_1$--$a_5$ by their influence
    on apparent stiffness at small and finite load.
  \item[\textbf{Fig.\ref{fig:s9_sensitivity_curves}.}]
    Response curves of apparent stiffness versus
    multiplicative factor on the most influential
    gradient constants, comparing small and finite load.
  \item[\textbf{Fig.\ref{fig:s9_bc_layer}.}]
    Close-up profiles of displacement, rotation, and strain
    within $O(\ell)$ of the supported edge, for clamped
    versus soft simple-support BCs, at small and finite load.
\end{enumerate}

\noindent (Figures will be renumbered at final submission to fit
the 8--12 figure budget; the current list reserves slots for all
planned outputs and will be trimmed to the most informative subset
once results are available.)

% =====================================================================
\subsection{Future extensions (out of scope for the present paper)}
\label{ssec:future}
% =====================================================================

The following extensions are acknowledged as natural follow-ons
but will not be pursued in the present paper, in order to keep
the numerical programme within the scope of a single Q1 journal
article:
\begin{itemize}
  \item dynamic and eigenvalue (free-vibration / buckling)
    analysis using the mass matrix $\mathbf{M}^{e}$ and
    inertial term $\mathbf{f}_{\mathrm{iner}}^{e}$ of
    Section~7;
  \item adoption of polyconvex or other globally stable
    constitutive models (neo-Hookean, Mooney--Rivlin, Ogden)
    in place of the Saint-Venant--Kirchhoff classical part,
    extending the stable deformation range;
  \item three-dimensional implementation using $C^1$
    solid-shell or trivariate $C^1$ elements;
  \item contact, impact, and large-rotation structural
    applications;
  \item model reduction to one- or two-length-scale
    gradient models informed by the Study~5 sensitivity
    results;
  \item frictional/interface effects and non-smooth
    boundary conditions;
  \item experimental validation against size-effect test
    data (validation in the sense of
    \cite{OberkampfTrucano2002,ASME_Verification_2012});
  \item adaptive mesh refinement driven by a posteriori
    error indicators for the fourth-order problem;
  \item optimisation / uncertainty quantification of the
    gradient constants against experimental or synthetic
    data.
\end{itemize}

% =====================================================================
\subsection{Closing statement}
\label{ssec:s9_closing}
% =====================================================================

The numerical programme specified in this section defines a
disciplined sequence of simulations that (i)~verifies the
code against the mandatory 2022 Mindlin Form-II benchmark,
(ii)~establishes mesh convergence in the nonlinear regime,
(iii)~characterises size effects through an internal-length-scale
sweep, (iv)~demonstrates geometrically nonlinear capability
through canonical load--deflection problems, (v)~quantifies
sensitivity to the five gradient constants, and (vi)~documents
the role of higher-order boundary conditions.  The programme
is sized for one Q1 journal paper (six studies, eight tables
and approximately a dozen figures, with explicit provision
for trimming) and remains fully compatible with the validated
2022 benchmark through the Study~1 gate.  No numerical value,
convergence rate, or observation is pre-judged; all will be
extracted from MATLAB computations and reported in the final
paper.

% =====================================================================
\bibliographystyle{unsrt}
\bibliography{references_calculation}

\end{document}
